% =====================================================================
%  Compression and Verification
%  Layered Emergence, Structural Convergence, and Verifiable Execution
%  in Complex Software Systems
%
%  English build — authoritative text.
%  The Chinese build in ../zh is a controlled rendering of this text;
%  the term mapping is fixed in ../TERMINOLOGY.md.
%
%  Compile: xelatex main.tex (three passes)
% =====================================================================
\documentclass[11pt,a4paper,oneside,openany]{book}

% =====================================================================
%  preamble.tex (English build — authoritative text)
%  Compression and Verification
%  Compile with: xelatex main.tex  (three passes)
% =====================================================================

% ---------- Page ----------
\usepackage{geometry}
\geometry{a4paper,left=2.7cm,right=2.7cm,top=2.9cm,bottom=2.9cm}
\setlength{\emergencystretch}{1.5em}
\tolerance=2000
\hbadness=10000

% ---------- Mathematics ----------
\usepackage{amsmath}
\usepackage{amssymb}
\usepackage{mathtools}
\usepackage{amsthm}

% ---------- Tables and lists ----------
\usepackage{array}
\setlength{\tabcolsep}{4pt}
\usepackage{booktabs}
\usepackage{longtable}
\usepackage{enumitem}
\setlist{nosep,leftmargin=2.2em}


% ---------- Colour, captions ----------
\usepackage{xcolor}
\usepackage{caption}
\captionsetup{font=small,labelfont=bf}

% ---------- Theorem environments ----------
% A run-in style is used for statements that are claims about practice rather
% than pure mathematics, so that they read as labelled paragraphs rather than
% as displayed blocks.
\newtheoremstyle{runin}%
  {0.7em}{0.2em}{\normalfont}{}{\bfseries}{.}{0.5em}{}

\theoremstyle{plain}
\newtheorem{theorem}{Theorem}[chapter]
\newtheorem{proposition}[theorem]{Proposition}
\newtheorem{lemma}[theorem]{Lemma}
\newtheorem{corollary}[theorem]{Corollary}
\newtheorem{conjecture}[theorem]{Conjecture}

\theoremstyle{definition}
\newtheorem{definition}[theorem]{Definition}
\newtheorem{hypothesis}[theorem]{Hypothesis}
\newtheorem{example}[theorem]{Example}
\newtheorem{problem}[theorem]{Problem}

\theoremstyle{runin}
\newtheorem{principle}[theorem]{Principle}

\theoremstyle{remark}
\newtheorem{remark}[theorem]{Remark}
\newtheorem{observation}[theorem]{Observation}

\renewcommand{\proofname}{Proof}
\renewcommand{\qedsymbol}{$\blacksquare$}

% ---------- Macros ----------
\newcommand{\Beh}{\mathcal{B}}
\newcommand{\Obs}{\mathcal{O}}
\newcommand{\Env}{\mathcal{A}}
\newcommand{\Cost}{\mathcal{C}}
\newcommand{\Inv}{\mathcal{I}}
\newcommand{\Mod}{\mathcal{M}}
\newcommand{\Tr}{\operatorname{Tr}}
\newcommand{\Allow}{\operatorname{Allowed}}
\newcommand{\Observe}{\alpha}
\newcommand{\Conc}{\gamma}
\newcommand{\Obl}{\operatorname{Obl}}
\newcommand{\Ker}{\operatorname{ker}}
\newcommand{\TCB}{\ensuremath{\mathrm{TCB}}}
\newcommand{\DL}{\ell}
\newcommand{\Entropy}{H}
\newcommand{\code}[1]{\texttt{#1}}
\newcommand{\term}[1]{\textbf{#1}}
\newcommand{\Lk}[1]{\ensuremath{\mathrm{L}_{#1}}}
\newcommand{\Dom}{\operatorname{Dom}}
\newcommand{\Conv}{\operatorname{Conv}}

\newcolumntype{L}[1]{>{\raggedright\arraybackslash}p{#1}}
\newcolumntype{C}[1]{>{\centering\arraybackslash}p{#1}}

% ---------- Hyperlinks (loaded last) ----------
\usepackage[colorlinks=true,
            linkcolor=blue!55!black,
            citecolor=teal!70!black,
            urlcolor=blue!65!black,
            bookmarksnumbered=true,
            unicode=true]{hyperref}


\title{Compression and Verification}
\author{The Compression--Assurance Working Notes}
\date{2026}
\hypersetup{pdfauthor={The Compression--Assurance Working Notes},
            pdftitle={Compression and Verification}}

\begin{document}

\frontmatter

\begin{titlepage}
\centering
\vspace*{2.6cm}
{\Huge\bfseries Compression and Verification\par}
\vspace{1.0cm}
{\LARGE Layered Emergence, Structural Convergence,\\[0.35em]
 and Verifiable Execution in Complex Software Systems\par}
\vspace{1.8cm}
{\large The Compression--Assurance Working Notes\par}
\vspace{0.45cm}
{\normalsize\itshape An independent technical report\par}
\vspace{0.7cm}
{\small\url{https://github.com/ghwinter/assay}\par}
\vspace{0.25cm}
{\footnotesize Correspondence via the repository issue tracker.\\
 Source of this report: \url{https://github.com/ghwinter/assay/tree/master/docs/paper}\par}
\vfill
{\normalsize 2026\par}
\vspace{1.0cm}
\end{titlepage}

\chapter*{Abstract}
\addcontentsline{toc}{chapter}{Abstract}

We formalise one question. \emph{Do large software systems converge in their high-level
abstractions, structures and architectures, in the way that low-level ABIs, calling
conventions and stack frames converge?} The intuition is decomposed into three
propositions of differing strength, each given a decidable or falsifiable form.

Four theses organise the work. The \term{compression thesis}: a complex system is
comprehensible to a bounded intelligence exactly to the extent that it is compressible;
comprehension is task-constrained lossy compression whose optimal descent depth is the
stationarity condition of a rate--distortion Lagrangian. The \term{necessity thesis}: any
system with concurrently live entities and contended finite resources, whose operational
semantics is well defined, has already defined six functions---identity and lifecycle,
allocation and reclamation, scheduling, naming and addressing, communication and
synchronisation, and fault isolation and recovery; these may be built or adopted, but not
omitted. The \term{separation thesis}: generation may be free, probabilistic and untrusted,
while admission must be deterministic, mechanical and exhaustively inspectable; formal-layer
guarantees come from a minimal trusted kernel, intent-layer guarantees from adversarial
consensus, and the two are stitched by artefact binding. The \term{design thesis}:
architecture is manufactured separation of scales---nature obtains levels free through
renormalisation, software must pay for them---and design is forward compression while
comprehension is inverse compression, two reading directions of one structure.

The formal results are as follows. A level's \emph{oblivious set} as the kernel of its
abstraction map, making ``what this level may ignore'' computable and reducing the reality
of a level to a compression ratio; Galois connections for inter-level semantic
preservation; a \emph{contract operad} whose well-definedness reduces the composability
of local guarantees to well-typed colouring and acyclicity of the interface graph,
thereby explaining why feedback and retry invalidate existing composition arguments; a
\emph{level family} distinguishing cross-domain isomorphism of the six functions from
transport along problem shape; near-decomposability as interface sufficiency, yielding a
diagnostic for \emph{leak pairs}; a \emph{coarse-graining criterion} making the
convergence thesis falsifiable within a single project; an amortisation inequality
$\Theta(nm)\to\mathcal{O}(n+m)$ giving the phase-transition threshold of internal
infrastructuralisation; expected read-set as a measurable architectural objective, with
invariant centralisation proved to minimise it; an \emph{observability rank} criterion
for whether comprehension can converge at all, separating insufficient observation from
insufficient sampling; \emph{governable complexity}
$C=\mathbb{E}[\rho]\cdot\lambda$ against comprehension capacity $K=hb$, yielding a
stability condition and a proof that architecture fixes the scaling exponent of required
headcount; technical debt as a description-length gap with a rate--distortion reading of
refactoring; a Hoeffding bound $\exp(-2n(1/2-p)^2)$ for consensus reliability, dominated
by a correlated error mode $q$ that is in turn reduced to an incentive-design problem;
and a consolidated assurance condition combining two-layer trust with artefact binding.

Twelve falsifiable hypotheses with operationalisations, controls, falsifiers and risks are
given, together with a six-phase roadmap, a paper-centred interdisciplinary literature map,
and an explicit list of non-claims rejecting the use of G\"odel incompleteness, Rice's
theorem, the open status of $\mathrm{P}$ versus $\mathrm{NP}$, near-decomposability, and
analogies to the Langlands programme or Landauer's principle, as load-bearing universal
laws.

\vspace{1em}
\noindent\textbf{Keywords:} software architecture; structural convergence; abstraction
levels; information compression; minimum description length; formal verification;
proof-carrying code; semantic preservation; trusted computing base; consensus;
intent gap

\tableofcontents
\listoftables

\mainmatter

% =====================================================================
\chapter{Introduction}
% =====================================================================

\section{The questions}

The following three questions share one structure. They are stated here in the form in
which they can be discussed, proved and falsified.

\begin{problem}[Generality of low-level determinism]\label{prob:low}
Software at the execution level is strikingly uniform: ABIs and calling conventions, the
call stack, heap and object lifetime, control and data flow, synchronisation primitives
and thread scheduling. That uniformity is imposed by a shared computational model,
hardware architecture, language runtime and operating-system interface. Does the same
determinism extend upwards? Do most large software systems, normally biased and normally
engineered, converge on the same or similar abstractions, structures and architectures?
\end{problem}

\begin{problem}[Coercion by scale]\label{prob:scale}
When a system crosses some order of magnitude, does it necessarily begin to design and
optimise for recurrent problems of one kind, and thereby grow a unified internal
infrastructure? An operating-system kernel defines how memory is managed and allocated,
how threads and processes are defined and scheduled, and what runtime interface is
offered upwards. A browser, a game engine, an asynchronous runtime, a distributed store
and a relational engine address the same class of questions at different levels. Once the
code base is large enough, do all of them exhibit the same kind of reasoning---unified
function, unified interfaces, managed planning, runtime structure?
\end{problem}

\begin{problem}[Comprehension and design]\label{prob:mind}
Software is written by humans and machines, but it is not easy to understand. A large
system can hardly be read sequentially from the micro level. At what level of abstraction
should it be approached, and how does comprehension lead into development without
dissolving into a sea of code? The question bears equally on machines: an intelligence
cannot descend into unbounded low-level information, and must retain an understanding of
first principles---call this \Lk{0}---together with the levels above it, \Lk{1}, \Lk{2},
\Lk{3}, and so on. What an intelligence needs is compressed information. Conversely, how
should a large complex system be designed so that its purpose is achieved?
\end{problem}

The third question has a contemporary variant that moves it from comprehension to
warranty:

\begin{problem}[Correctness without review]\label{prob:trust}
If machine generation drives the cost of code towards zero and no human reads the
implementation line by line, how is correctness guaranteed \emph{mechanically} and
\emph{non-self-referentially}? How is it guaranteed that the stated purpose is achieved,
and that no unauthorised side effect accompanies it? A concrete instance: a hand-written
call takes 1\,ns; a generated implementation takes 1\,$\mu$s. Functionally identical,
purpose achieved, and a performance side effect introduced.
\end{problem}

\section{Three gaps in existing treatments}

Each question has local counterparts in the literature: architectural convergence in
pattern languages and complex-systems theory, coercion by scale in information hiding and
near-decomposability, comprehension and design in cognitive models and multi-view
modelling, correctness in program logic, refinement calculi and verified compilation.
Treating them as a whole exposes three structural gaps.

\paragraph{The strength of the convergence claim is not graded.}
``Large software architectures converge'' contains at least three propositions of very
different strength: recurrent local architectural patterns (well supported by engineering
experience); measurable structural similarity above a random baseline after abstraction
(plausible and worth systematic testing); and convergence of nearly all large systems on
one high-level architecture (unsupported, and probably too strong). Conflating them makes
the claim untestable.

\paragraph{Necessity claims lack a load-bearing structure.}
``Scale forces a unified infrastructure'' is a necessity claim, but it is usually stated
as an empirical observation, and therefore cannot answer where the necessity comes from,
what fixes the threshold, or why some equally large systems grow no kernel. The answers
given below are that the necessity comes from the completeness of the operational
semantics, the threshold from an amortisation inequality, and the apparent exceptions
from systems that \emph{adopt} a kernel rather than build one.

\paragraph{Comprehension, design and assurance are three readings of one structure.}
Comprehension is treated as a cognitive problem, design as an engineering problem and
assurance as a logical problem, in three separate literatures. They are in fact different
reading directions of one object: a level is the unit into which assurance obligations
decompose; comprehension recovers that decomposition and design prescribes it.

\section{Core theses}

Four theses organise the paper. Each is annotated with its formal location and its
strength: \textbf{[P]} proved, \textbf{[A]} argued under stated assumptions,
\textbf{[H]} advanced as a falsifiable hypothesis.

\begin{principle}[Compression]\label{th:compression}\textbf{[A]}
A complex system is comprehensible to a bounded intelligence if and only if it is
compressible. Comprehension is task-constrained lossy compression: description length is
minimised subject to preserved task-relevant predictive power. The optimal descent depth
is therefore the stationarity condition of a rate--distortion Lagrangian, that is, the
point at which marginal predictive gain vanishes.
\end{principle}

\begin{principle}[Necessity]\label{th:necessity}\textbf{[A]}
Any system with concurrently live entities and contended finite resources, whose
operational semantics is well defined, has already defined six functions: identity and
lifecycle, allocation and reclamation, scheduling, naming and addressing, communication
and synchronisation, and fault isolation and recovery. These may be built or adopted from
a host---an operating system, a language runtime, a framework---but they cannot be omitted
from the semantics.
\end{principle}

\begin{principle}[Separation]\label{th:separation}\textbf{[A]}
Generation may be free, probabilistic, open and untrusted; admission must be
deterministic, mechanical and exhaustively inspectable. Guarantees on the formal layer
(specification $\to$ implementation $\to$ binary $\to$ execution) come from mechanical
decision by a minimal trusted kernel. Guarantees on the intent layer (intent $\to$
specification) come from adversarial consensus, and that layer admits no mechanical
solution. The two are stitched by artefact binding: only the verified binary is admitted
to run, and the environmental assumptions are checked or monitored on it.
\end{principle}

\begin{principle}[Design]\label{th:design}\textbf{[A]}
Architecture is manufactured separation of scales. Nature obtains levels free from scale
separation and renormalisation; software levels do not emerge spontaneously and must be
paid for in interfaces, encapsulation and invariants, and they leak. Design is therefore
forward compression and comprehension is inverse compression, two reading directions of
one structure. The measurable objective of architecture is the expected read-set
$\mathbb{E}[\rho]$ of a typical change.
\end{principle}

The four theses and their formal results are referred to collectively as the
\term{compression--assurance framework}. The name is taken from the two ends of the claim:
\emph{compression} is the mechanism of comprehension and design, \emph{assurance} is the
mode in which trust is anchored. The central assertion is that a complex system is
comprehensible to a bounded agent if and only if it is compressible, and that it is
compressible if and only if it admits a separation of scales in which every level has an
explicitly declared oblivious set.

\section{Epistemic discipline}

The material crosses engineering intuition, complex-systems theory, information theory,
mathematical logic and social epistemology, and the risk of overreach is high. Five rules
are observed throughout.

\begin{enumerate}[label=(\roman*)]
\item \term{Strength labelling.} Every claim is labelled as a theorem, proposition,
      conjecture, hypothesis or heuristic. Heuristics are not written as theorems.
\item \term{Non-claims list.} Chapter \ref{ch:objections} states explicitly what the paper
      does \emph{not} claim, including several ``authoritative theorems'' that are declined.
\item \term{Explicit assumptions.} Any conclusion that depends on an environmental
      assumption states it, and any guarantee that has a boundary marks it.
\item \term{Falsifiability.} Chapter \ref{ch:program} reduces the core claims to fourteen
      falsifiable hypotheses with operationalisations, controls, falsifiers and risks.
\item \term{Verification not isomorphic to the conclusion.} Where the paper argues that
      generation must be untrusted and verification independent, that rule applies to the
      paper itself: its load-bearing propositions must be independently checkable and must
      not depend on any unproved conjecture.
\end{enumerate}

Rule (v) constrains the paper's own rhetoric. Several claims have been removed from the
load-bearing structure: G\"odel incompleteness and Rice's theorem are not used to prove
that a gap between intent and specification is principled
(Chapter \ref{ch:intent}, \S\ref{sec:intent-gap}); the open status of $\mathrm{P}$ versus
$\mathrm{NP}$ is not used to support ``verification is cheaper than generation''
(\S\ref{sec:asymmetry}); and near-decomposability is not treated as a universal law
(Chapters \ref{ch:convergence} and \ref{ch:objections}). Each of these gains strength by
being unloaded rather than reinforced.

\section{Contributions}

Formal contributions:

\begin{enumerate}[label=(C\arabic*)]
\item The \emph{oblivious set} of a level as the kernel of its abstraction map, which makes
      ``what this level need not care about'' a computable object
      (Definition \ref{def:oblivious}).
\item Galois connections as the condition for inter-level semantic preservation
      (Definitions \ref{def:galois}, \ref{def:soundness}, Proposition \ref{prop:soundness}).
\item A \emph{contract operad} whose well-definedness reduces the composability of local
      guarantees to well-typed colouring and acyclicity of the interface graph
      (Definition \ref{def:operad}, Proposition \ref{prop:operad-compose}).
\item A \emph{level family} $\pi:\mathcal{L}\to\mathcal{P}$ that separates cross-domain
      isomorphism of the six functions from transport along problem shape
      (Definition \ref{def:fibration}, Proposition \ref{prop:transport}).
\item A \emph{coarse-graining criterion} under which the convergence thesis can be
      falsified inside a single project (Definition \ref{def:coarsegrain},
      Conjecture \ref{conj:relevant}).
\item An amortisation inequality comparing $\Theta(nm)$ direct interactions with
      $\mathcal{O}(nk+m)$ substrate interactions, giving the break-even threshold of internal
      infrastructuralisation (Theorem \ref{th:integration}).
\item Conditional necessity of the six functions from the well-definedness of the
      operational semantics (Theorem \ref{th:six}), and the resulting build-or-adopt
      criterion (Corollary \ref{cor:buildadopt}).
\item A criterion for over- and under-constraint that renders ``unique below, manifold
      above'' as solution-set geometry (Proposition \ref{prop:overconstrain}).
\item Expected read-set as the measurable objective of architecture, with a bound showing
      that invariant centralisation lowers it, stated relative to the enforcement sites that
      no independent evidence settles (Definition \ref{def:readset}, Definition
      \ref{def:enforce}, Theorem \ref{th:centralise}).
\item Technical debt as a description-length gap (Definition \ref{def:debt}), with a
      rate--distortion reading of refactoring.
\item An identifiability criterion for whether comprehension can converge, separating
      insufficient observation from insufficient sampling, of which the linear observability
      rank is the finite-dimensional specialisation (Definition \ref{def:observability},
      Proposition \ref{prop:observability}).
\item A governable-complexity stability condition and a scaling law for required
      headcount (Definition \ref{def:governable}, Proposition \ref{prop:exponent}).
\item A Hoeffding bound for consensus reliability and the proof that a correlated error
      mode dominates it, reducing the governing parameter to an incentive-design problem
      (Theorem \ref{th:condorcet}, Proposition \ref{th:correlated}, \S\ref{sec:mechanism}).
\item The observation that the mechanical and consensus layers are two enrichments of one
      structure, and that no information-preserving change of base runs from the second to the
      first; consequently reliability must be \emph{budgeted} across the stages that must all hold
      rather than gated stage by stage, a per-stage threshold giving a composite bound that grows
      linearly and is vacuous once the number of stages reaches the reciprocal of the
      threshold (\S\ref{sec:enrichment}, Principle \ref{pr:budget}). The mathematical content here
      is Lawvere's, not ours; the contribution is its application to the two-layer structure and
      the budget rule that follows.
\item A consolidated assurance condition combining two-layer trust with artefact binding
      (Definition \ref{def:assurance}, Proposition \ref{prop:assurance}).
\end{enumerate}

Methodological contributions:

\begin{enumerate}[label=(M\arabic*)]
\item A descent protocol that converts comprehension of a large system from sequential
      reading into layered modelling and active identification.
\item An executable list of seam detectors for locating real module boundaries in an
      unfamiliar code base.
\item A design procedure that replaces adjectival goals with an optimisable objective and
      a stability condition.
\item Fourteen falsifiable hypotheses with experimental designs.
\item A cross-disciplinary research position that couples social epistemology to
      proof-carrying code, which the literature map identifies as unoccupied.
\end{enumerate}

\section{Structure}

Chapter \ref{ch:framework} establishes the formal language: levels, abstraction maps,
oblivious sets, contracts, traces, assurance and trust. Chapter \ref{ch:convergence}
addresses Problem \ref{prob:low}, grading the convergence claim and stating its strength
law. Chapter \ref{ch:infrastructure} addresses Problem \ref{prob:scale}, proving the
conditional necessity of the six functions and deriving the break-even threshold at which a substrate pays for itself.
Chapter \ref{ch:comprehension} addresses the first half of Problem \ref{prob:mind}:
comprehension as inverse compression. Chapter \ref{ch:design} addresses the second half:
design as forward compression. Chapter \ref{ch:intent} addresses Problem \ref{prob:trust}
with a verifiable refinement chain and a two-layer trust structure. Chapter
\ref{ch:levels} promotes levels from a comprehension aid to the unit into which assurance
obligations decompose. Chapter \ref{ch:program} states the falsifiable programme. Chapter
\ref{ch:literature} positions the work. Chapter \ref{ch:objections} lists objections and
limits. Chapter \ref{ch:conclusion} concludes.

\section{Terminology}

Several terms drift between communities. They are fixed once here; the full glossary is
Appendix \ref{app:glossary}, and the Chinese rendering is fixed term by term in the
companion terminology file.

\begin{longtable}{L{3.1cm}L{11.4cm}}
\toprule
\textbf{Term} & \textbf{Meaning in this paper} \\
\midrule
\endhead
level & A tuple (representation space, behaviour set, observation map, environment
  assumptions, cost model, invariants), defined relative to a study object and a task;
  a cognitive scale, not a physical directory. \\
oblivious set & The differences a level can never distinguish; the kernel of its
  abstraction map. \\
separation of scales & The condition under which adjacent levels each support a valid
  description and the lower level's implementation details do not appear in the upper
  level's semantics. \\
internal infrastructuralisation & The process by which a system promotes recurrent
  management problems to unified mechanisms on which the upper layers then depend. \\
the six functions & Identity and lifecycle; allocation and reclamation; scheduling;
  naming and addressing; communication and synchronisation; fault isolation and recovery. \\
read-set & The set of modules that must be read to complete a typical change. \\
assurance & A traceable evidential relation holding a property within stated premises and
  a stated scope. \\
mechanical assurance & Assurance by a single deterministic, exhaustively inspectable
  decision procedure. \\
consensus assurance & Assurance by multiple independent, adversarial participants. \\
artefact binding & The identity relation between the verified artefact and the executed
  artefact. \\
intent gap & The distortion between an intent and its formal projection. \\
\bottomrule
\end{longtable}

One terminological caution is load-bearing. A \emph{level} is a cognitive scale, not a
physical tier. The same component occupies different levels in different tasks: an
asynchronous runtime is infrastructure to an application developer and a complex system in
its own right to its maintainer. Levels are defined relative to the study object and the
task (Chapter \ref{ch:levels}).

% =====================================================================
\chapter{Formal framework}\label{ch:framework}
% =====================================================================

This chapter establishes the shared formal language. It does not invent new mathematics;
it names objects that engineering practice already uses but does not define: what a level
may ignore, under what conditions semantic preservation holds across levels, how far
comprehension should descend, and within what scope a guarantee holds.

\section{Levels and abstraction maps}\label{sec:level}

\begin{definition}[Level]\label{def:level}
A \term{level} is a six-tuple
\[
\Lambda \;=\; \bigl(X,\; \Beh,\; \Obs,\; \Env,\; \Cost,\; \Inv\bigr),
\]
where $X$ is the \term{representation space} of states admissible at that level;
$\Beh \subseteq X^{\mathbb{N}}$ is the \term{behaviour set} of admissible execution
traces; $\Obs$ is the \term{observation map} projecting the level's behaviour onto the
observable algebra of the level above; $\Env$ is a set of \term{environmental assumptions}
on which the level's conclusions depend; $\Cost : \Beh \to \mathbb{R}_{\ge 0}^{k}$ is a
\term{cost model} with $k$ components (time, space, energy, bandwidth); and
$\Inv \subseteq \Beh$ is the \term{invariant set} that must hold under every admissible
transition.
\end{definition}

The six components are not of equal weight. $X$ and $\Beh$ describe what exists, $\Obs$
describes what cannot be seen, $\Env$ describes what is depended upon, $\Cost$ describes
what it costs, and $\Inv$ describes what must not break. Chapter \ref{ch:levels} argues
that $\Obs$ and $\Inv$ decide whether a level is real: a level that cannot state its
observation map and its invariants is not a level but a directory.

\begin{definition}[Abstraction and concretisation]\label{def:galois}
Let $\Lambda_i$ and $\Lambda_{i+1}$ be adjacent levels, with the smaller index the upper
level. The \term{abstraction map} $\Obs_i : \Beh_{i+1} \to \Beh_i$ projects lower-level
behaviour onto upper-level behaviour; the \term{concretisation map}
$\Conc_i : \Beh_i \to 2^{\Beh_{i+1}}$ expands an upper-level behaviour into the set of its
lower-level implementations. If there are partial orders $\sqsubseteq_i$ and
$\sqsubseteq_{i+1}$ such that
\begin{equation}\label{eq:galois}
b \sqsubseteq_{i+1} \Conc_i(b') \quad\Longleftrightarrow\quad \Obs_i(b) \sqsubseteq_i b',
\qquad \forall b \in \Beh_{i+1},\; b' \in \Beh_i,
\end{equation}
then $(\Obs_i, \Conc_i)$ is a \term{Galois connection}, written $\Obs_i \dashv \Conc_i$.
\end{definition}

A Galois connection states precisely what abstraction discards and what it retains, and is
the central tool of the abstract-interpretation framework \cite{cousot1977}. Where an
inter-level relation can be expressed as a Galois connection, the semantic-preservation
condition becomes mechanically decidable; where it cannot, one must fall back on
simulation relations or on an explicit human argument.

\begin{definition}[Inter-level soundness]\label{def:soundness}
Let $A \in \Env_i$ be an environmental assumption declared at the upper level, and let
$\Beh_{i+1}\!\upharpoonright\! A$ denote the lower-level behaviours admissible under $A$.
Then $\Lambda_i$ is \term{sound} with respect to $\Lambda_{i+1}$ under $A$ if
\begin{equation}\label{eq:soundness}
\Obs_i\bigl(\Beh_{i+1}\!\upharpoonright\! A\bigr) \;\subseteq\; \Beh_i .
\end{equation}
\end{definition}

\begin{proposition}[Composition of soundness]\label{prop:soundness}
Let there be a chain $\Lambda_0, \Lambda_1, \dots, \Lambda_n$, with increasing index
denoting increasing concreteness, and suppose that for each $i$ there exists
$A_i \in \Env_i$ satisfying \eqref{eq:soundness}. If
\begin{enumerate}[label=(\alph*)]
\item \emph{observation compatibility}: $\Obs_i \circ \Obs_{i+1} = \Obs_i^{(i+2)}$, where
      $\Obs_i^{(i+2)}$ is the direct observation from $\Beh_{i+2}$ to $\Beh_i$; and
\item \emph{assumption closure}:
      $\Obs_i(\Beh_{i+1}\!\upharpoonright\! A_i) \subseteq \Beh_i\!\upharpoonright\! A_{i-1}$,
      that is, the assumptions satisfied below suffice to discharge those above,
\end{enumerate}
then $\Obs_0^{(n)}\bigl(\Beh_n \upharpoonright \bigcap_i A_i\bigr) \subseteq \Beh_0$.
\end{proposition}

\begin{proof}
Induction on $n$. The case $n = 1$ is \eqref{eq:soundness}. Condition (b) ensures that at
each inductive step the upper level's environmental assumption holds of the upper-level
behaviour set, so that \eqref{eq:soundness} can be applied level by level. Condition (a)
ensures that the composite observation agrees with the level-by-level observations, so
that the conclusion holds at the top level.
\end{proof}

\begin{remark}[Why composition fails]\label{rmk:why-compose-fails}
Proposition \ref{prop:soundness} is the precise form of ``how do local guarantees compose
into an end-to-end guarantee'', and equally the precise form of why they often do not:
composition requires both conditions simultaneously. In practice (a) fails when observation
maps disagree---for instance when ``the transaction committed'' denotes different
observations at the application layer and the storage layer---and (b) fails when
environmental assumptions are not closed, for instance when the upper layer assumes that
``a write is durable'' while the lower layer's durability depends on an undeclared flush
policy. Chapter \ref{ch:intent} uses these two failure modes as the skeleton of its
evidence-boundary table.
\end{remark}

\begin{remark}[Relation to the source model of the companion report]\label{rmk:source}
Definition \ref{def:level} extends the four-tuple
$M_i = (X_i, \mathrm{Beh}_i, \mathrm{Obs}_i, A_i)$ used in the companion research report
by adding $\Cost$ and $\Inv$, which Chapters \ref{ch:design} and \ref{ch:intent} require in
order to place resource guarantees and invariant enforcement inside the level object. The
companion report describes inter-level semantics by an \emph{interpretation relation}
$R_i \subseteq \mathrm{Beh}_{i+1} \times \mathrm{Beh}_i$, requiring that for every
admissible target behaviour $b_{i+1}$ of a successful transformation
$x_{i+1} = T_i(x_i)$ there exists a source behaviour $b_i$ with $R_i(b_{i+1}, b_i)$. That
is a \emph{simulation condition} and is more general than the set inclusion of
\eqref{eq:soundness}; the two coincide when $R_i$ is induced by a function, that is when
$R_i = \{(b, \Obs_i(b))\}$. This paper uses $R_i$ where a general relation is needed and
$\Obs_i$ where the abstraction is functional. The companion report stresses that, across
different languages and observation domains, $R_i$ may be a forward simulation, a backward
simulation, a bisimulation, a contextual equivalence or a refinement relation, and that
one must not assume a single relation throughout. That qualification is accepted here.
\end{remark}

\section{Oblivious sets and compression}\label{sec:oblivious}

The next two definitions are the formal centre of the paper. They turn ``what this level
need not care about'' from an engineering phrase into a computable object.

\begin{definition}[Oblivious set]\label{def:oblivious}
The \term{oblivious set} of level $\Lambda_i$ with respect to the level below is the kernel
of the abstraction map:
\begin{equation}\label{eq:oblivious}
\Obl(\Lambda_i) \;:=\; \Ker \Obs_i \;=\;
\bigl\{(b,b') \in \Beh_{i+1}^{2} \;:\; \Obs_i(b) = \Obs_i(b')\bigr\}.
\end{equation}
If $(b,b') \in \Obl(\Lambda_i)$, the level is \term{oblivious} to the difference between
$b$ and $b'$.
\end{definition}

\begin{proposition}[Basic properties]\label{prop:oblivious}
$\Obl(\Lambda_i)$ is an equivalence relation on $\Beh_{i+1}$, and the quotient satisfies
\[
\Beh_{i+1} \big/ \Obl(\Lambda_i) \;\cong\; \Obs_i(\Beh_{i+1}) .
\]
\end{proposition}

\begin{proof}
The three properties of an equivalence relation follow from those of $=$. The surjection
$\Obs_i: \Beh_{i+1} \to \Obs_i(\Beh_{i+1})$ induces a bijection
$\Beh_{i+1}/\Ker\Obs_i \to \Obs_i(\Beh_{i+1})$, by the universal property of the quotient.
\end{proof}

Proposition \ref{prop:oblivious} yields a direct consequence: \emph{abstraction is
quotienting by the kernel}. Everything a level exposes to the level above is the quotient
of lower-level behaviour by that kernel. Hence:

\begin{itemize}
\item ``what this level guarantees'' is an element of the quotient
      $\Obs_i(\Beh_{i+1})$;
\item ``what this level does not guarantee'' is what has been quotiented away, namely
      $\Obl(\Lambda_i)$;
\item ``what this level may ignore'' is every lower-level detail compatible with
      $\Obl(\Lambda_i)$.
\end{itemize}

The third item has the deepest consequence: it makes ignorance a property that can be
declared, checked and held to account. A module that has never declared its oblivious set
has an incomplete contract.

\begin{remark}[A criterion for pattern realism]\label{rmk:dennett}
Proposition \ref{prop:oblivious} is isomorphic to an independent philosophical criterion.
Dennett's criterion for a \emph{real pattern} states that a pattern is real exactly when
describing the system in terms of it is more economical than describing the raw data
\cite{dennett1991}. Under that criterion, whether a level exists is not an ontological
question but a question of compression ratio---precisely the quantity measured by
Definition \ref{def:compression}.

Three consequences follow, and are used in later chapters.
\begin{enumerate}[label=(\alph*)]
\item Levels are neither discovered natural kinds nor arbitrary conventions; they are fixed
      points of a compression operator.
\item The convergence thesis of Chapter \ref{ch:convergence} needs no separate
      philosophical defence: if two systems compress independently under similar
      constraints, that their fixed points are similar follows from properties of the
      compression operator rather than from an additional assumption about the world.
\item The criterion is quantitative, so whether a given level is real is measurable rather
      than arguable.
\end{enumerate}
\end{remark}

\begin{definition}[Compression ratio]\label{def:compression}
Let $\mu$ be a reference measure on $\Beh_{i+1}$ and let $\Entropy$ denote the
corresponding entropy functional. The \term{compression ratio} of level $\Lambda_i$ is the
information annihilated:
\begin{equation}\label{eq:compression}
\kappa_i \;=\; \Entropy_{\mu}\bigl(\Beh_{i+1}\bigr) \;-\;
\Entropy_{\Obs_i \# \mu}\bigl(\Obs_i(\Beh_{i+1})\bigr),
\end{equation}
that is, the information that cannot be recovered under $(\Obs_i, \Conc_i)$. The wording
is borrowed from information theory \cite{shannon1948} and algorithmic information theory
\cite{kolmogorov1965}.
\end{definition}

Compression is not itself the objective: a system that guarantees nothing attains
arbitrarily large $\kappa$. The objective is to maximise $\kappa$ \emph{subject to
preserving task-relevant predictive power}, which is the subject of
\S\ref{sec:comprehension-formal}.

\section{Architecture as a typed relational graph}\label{sec:relations}

Describing a system by levels alone is insufficient. Chapter \ref{ch:convergence} argues
that cross-system similarity resides mainly in \emph{relations} rather than in components,
so the relations themselves must be typed.

\begin{definition}[Architecture relational graph]\label{def:relgraph}
The \term{architecture relational graph} of a system is a typed directed multigraph
\begin{equation}\label{eq:relgraph}
\mathcal{G} \;=\; \bigl(V,\; E_{\mathrm{abs}} \sqcup E_{\mathrm{dep}} \sqcup
E_{\mathrm{cmp}} \sqcup E_{\mathrm{ctl}} \sqcup E_{\mathrm{ctr}}\bigr),
\end{equation}
where $V$ is a set of components and the five edge classes are:
\begin{enumerate}[label=(\alph*)]
\item $E_{\mathrm{abs}}$, \term{abstraction}: $u \to v$ when several details of $v$ are
      expressed uniformly by $u$; annotated with which semantics of $v$ are retained and
      which are hidden.
\item $E_{\mathrm{dep}}$, \term{dependency}: $u \to v$ when $u$ requires $v$ to discharge
      its own duty; annotated with whether the dependency is on an abstract interface or on
      a concrete implementation.
\item $E_{\mathrm{cmp}}$, \term{composition}: $u \to v$ when $u$ is constituted by $v$;
      annotated with ownership versus mere use.
\item $E_{\mathrm{ctl}}$, \term{control}: $u \to v$ when execution passes from $u$ to $v$;
      annotated with synchronous call, asynchronous message, event callback or data-driven
      dispatch, and with who schedules.
\item $E_{\mathrm{ctr}}$, \term{contract}: $u \to v$ when $u$ promises something to $v$ and
      requires something of $v$; annotated with preconditions, postconditions, invariants
      and failure semantics.
\end{enumerate}
\end{definition}

Keeping the five classes separate is not formalism for its own sake. Consider a single
call graph: it would compress four distinct facts into one edge class---that a thread pool
composes threads and queues, that it depends on synchronisation primitives, that it is
invoked through a submission interface, and that it promises to accept no new tasks after
shutdown---and would lose all structural information in doing so. The seam detectors of
Chapter \ref{ch:comprehension} are, in effect, procedures for recovering this five-way
distinction.

\begin{definition}[Abstract skeleton]\label{def:skeleton}
A \term{skeleton map} $\phi$ replaces each component by its role label in some duty
classification and retains the type of each relation while discarding its instance. The
image $\phi(\mathcal{G})$ is the \term{abstract skeleton} of $\mathcal{G}$.
\end{definition}

The central proposition of Chapter \ref{ch:convergence} is stated with $\phi$:
$\phi(\mathcal{G}_1) \approx \phi(\mathcal{G}_2)$ can hold even when
$\mathcal{G}_1$ and $\mathcal{G}_2$ have nothing in common at the source level.

\section{Comprehension as task-constrained lossy compression}\label{sec:comprehension-formal}

\begin{definition}[Task]\label{def:task}
A \term{task} is a pair $T = (Q, d)$, where $Q$ is a family of \term{queries}
$q : \Beh \to \mathcal{A}_q$ mapping system behaviour to observable answers, and $d$ is a
loss functional on answers.
\end{definition}

Typical queries are: if I change this, which behaviours change? under what conditions does
this path fail? who enforces this invariant? which modules must change to add a new entity
type?

\begin{definition}[$T$-sufficiency]\label{def:sufficient}
Let $M$ be a model of system $S$ at some level of abstraction. $M$ is
\term{$T$-sufficient} for $S$ if
\begin{equation}\label{eq:sufficient}
\forall q \in Q:\quad q\bigl(\Beh(M)\bigr) = q\bigl(\Beh(S)\bigr),
\end{equation}
or, in the approximate form,
$\mathbb{E}_{q \sim Q}\bigl[d\bigl(q(\Beh(M)), q(\Beh(S))\bigr)\bigr] \le D_T$.
\end{definition}

\begin{definition}[Description length and the rate--distortion objective]\label{def:mdl}
The definition rests on the minimum description length principle
\cite{rissanen1978,grunwald2006,grunwald2019}. Let $\DL(M)$ be the description length of
$M$ under a fixed encoding, and let
$L_T(M) = \mathbb{E}[d(q(\Beh(M)), q(\Beh(S)))]$ be the task loss. Comprehending $S$ means
solving
\begin{equation}\label{eq:ratedistortion}
M^{\star} \;=\; \arg\min_{M} \; \DL(M)
\qquad \text{s.t.} \qquad L_T(M) \le D_T ,
\end{equation}
or equivalently the Lagrangian form $\min_M \; \DL(M) + \lambda\, L_T(M)$.
\end{definition}

\begin{proposition}[Stopping rule for descent]\label{prop:stopping}
Parameterise the model family by descent depth, $M_0 \succ M_1 \succ \cdots$, with
increasing index meaning increasing retention of lower-level detail, and suppose
$\DL(M_{k+1}) \ge \DL(M_k)$ and $L_T(M_{k+1}) \le L_T(M_k)$. Then the optimal depth
$k^{\star}$ satisfies the stationarity condition
\begin{equation}\label{eq:stopping}
\left.\frac{\partial L_T}{\partial \DL}\right|_{k^{\star}} \;=\; -\frac{1}{\lambda} ,
\end{equation}
and once $\partial L_T / \partial \DL = 0$ beyond some depth, further descent is strictly
dominated.
\end{proposition}

\begin{proof}
The Lagrangian is $\mathcal{L}(k) = \DL(M_k) + \lambda L_T(M_k)$. First-order optimality
gives $\mathrm{d}\DL/\mathrm{d}k + \lambda\, \mathrm{d}L_T/\mathrm{d}k = 0$, that is
$\mathrm{d}L_T/\mathrm{d}\DL = -1/\lambda$. If
$\mathrm{d}L_T/\mathrm{d}\DL = 0$ on an interval while $\mathrm{d}\DL/\mathrm{d}k > 0$,
then $\mathrm{d}\mathcal{L}/\mathrm{d}k = \mathrm{d}\DL/\mathrm{d}k > 0$ and the objective
increases monotonically, so further descent is strictly dominated.
\end{proof}

Proposition \ref{prop:stopping} yields a rule usable in practice.

\begin{principle}[Descent depth]\label{pr:descent}
To make a change correctly it suffices to descend until the invariant governing that change
becomes visible, and then to stop. Further descent increases description length without
reducing task loss.
\end{principle}

The rule explains both why reading all the source is not comprehension and why reading only
an architecture diagram is not comprehension either: the former passes the stationary point
into the dominated region, the latter has not yet reached $T$-sufficiency.

\section{Contracts, traces and assurance}\label{sec:contract}

\begin{definition}[Behaviour contract]\label{def:contract}
A \term{behaviour contract} is an eight-tuple
\begin{equation}\label{eq:contract}
S \;=\; \bigl(P,\; Q,\; E^{+},\; E^{-},\; R,\; I,\; F,\; A\bigr),
\end{equation}
where $P$ is a precondition; $Q$ a postcondition; $E^{+}$ the set of permitted side
effects; $E^{-}$ the set of forbidden side effects; $R$ a resource budget (time, space,
bandwidth, energy); $I$ the invariants that must hold throughout execution; $F$ the
failure, cancellation, retry and recovery semantics; and $A$ the environmental assumptions
(hardware, operating system, load, external services).
\end{definition}

Equation \eqref{eq:contract} is not claimed to be a recognised specification format. Its
purpose is to state that \emph{correctness is not compressed into a return value but
defined as a set of constraints on the whole execution}. That is the starting point of
Chapter \ref{ch:intent}.

\begin{definition}[Traces and the admissible set]\label{def:traces}
Let $B$ be an executable artefact and $E$ a runtime environment (processor, instruction
set, operating system, loader, configuration, external services and devices). Let
$\Tr(B,E)$ be the set of \emph{observable execution traces} $B$ may produce in $E$,
including input and output, system calls, state mutations, I/O, concurrency events and
failure behaviour, and let $\Allow(S)$ be the traces permitted by contract $S$. Then $B$
\term{satisfies} $S$ in $E$ if
\begin{equation}\label{eq:inclusion}
\Tr(B,E) \;\subseteq\; \Allow(S).
\end{equation}
\end{definition}

Inclusion \eqref{eq:inclusion} is far stronger than output correctness: it constrains
\emph{all possible behaviours}, including side effects that were never intended.

\begin{definition}[Assurance and its condition]\label{def:assurance}
An \term{assurance} is a four-tuple (property assured, object, premises, evidence). A
property $\varphi$ holds within scope $\Sigma$ exactly when
\begin{equation}\label{eq:assurance}
\begin{aligned}
&\text{environmental assumptions hold} \\
&\land\; \text{the contract is explicit and authorised} \\
&\land\; \text{the verifier accepts the implementation and its proof} \\
&\land\; \text{the transformations preserve the required semantics} \\
&\land\; \text{the execution boundary is effective} \\
&\land\; \text{the object executed is the object verified} \\
&\qquad\Longrightarrow\; \varphi \text{ holds within the declared scope } \Sigma .
\end{aligned}
\end{equation}
\end{definition}

The last conjunct---called here the \term{fourth pillar}---is frequently omitted, but it
is not a detail: even if a verifier correctly accepts a binary, the preceding conclusion
does not transfer if the deployment system actually runs a different file, or if loading
changes a critical premise. Its formal expression is an identity assertion:
\begin{equation}\label{eq:binding}
\mathrm{digest}\bigl(B_{\text{deployed}}\bigr) \;=\;
\mathrm{digest}\bigl(B_{\text{verified}}\bigr)
\quad \land \quad
E_{\text{deployed}} \models A_{\text{claimed}} .
\end{equation}
Equation \eqref{eq:binding} is why content addressing, reproducible builds and load-chain
integrity belong to the theory rather than to deployment housekeeping.

\section{Two modes of trust generation}\label{sec:trust}

\begin{definition}[Trusted computing base]\label{def:tcb}
For an assurance regime, the \term{trusted computing base} $\TCB$ is the set of components
that must be trusted unconditionally for the assurance to hold: the proof checker and its
execution, the logic and semantic definitions in use, the identity between the program a
proof is about and the program that executes, the compilation, linking, loading and
runtime mechanisms, and the assumed hardware and environmental conditions.
\end{definition}

\term{TCB minimisation} is the precise restatement of the intuition that a system must not
verify itself. Its exact content is this: a mechanical system does not \emph{betray}, but
it can be defective, can fail, and can embody an incomplete model. The basis of trust is
therefore not that the mechanism is error-free, but that its decision does not depend on
the generator's will. The two are not the same.

\begin{definition}[Mechanical and consensus assurance]\label{def:twotrust}
\term{Mechanical assurance} is provided by a single deterministic, exhaustively inspectable
decision procedure whose objects are finite, static and unambiguous symbolic objects, such
as proof terms, typing derivations or constraint systems. \term{Consensus assurance} is
provided by the adversarial scrutiny of multiple independent participants, and its objects
are high-dimensional, dynamic and partly self-referential judgements, such as whether a
specification faithfully expresses an intent.
\end{definition}

The two modes are complementary rather than competing. Table \ref{tab:twotrust} states the
difference.

\begin{table}[htbp]
\centering
\caption{The two modes of trust generation}\label{tab:twotrust}
\begin{tabular}{L{2.7cm}L{5.3cm}L{5.3cm}}
\toprule
\textbf{Dimension} & \textbf{Mechanical assurance} & \textbf{Consensus assurance} \\
\midrule
Generation & A minimal decision procedure examines symbolic objects &
  Independent parties scrutinise adversarially and sign \\
Object & Formalised properties & Judgements that cannot be formalised \\
Strength & Absolute under the declared premises & Statistically reliable, not absolute \\
Failure mode & Defective decision procedure; lapsed premise &
  Collusion, capture, \emph{cognitive isomorphism} \\
Failure shape & Discrete (defective or not) & Continuous (higher independence, lower failure) \\
Cost & Proportional to the complexity of the property & Proportional to participant count and independence \\
Disciplinary analogue & Checking a mathematical proof & Peer review, a jury, the judiciary \\
\bottomrule
\end{tabular}
\end{table}

The load-bearing row is failure shape. Mechanical assurance fails discretely---not because
it is more reliable, but because its object is a finite static object and the check can be
exhausted. Consensus assurance fails continuously, because its object is intentional and
can only be approached by a network. \S\ref{sec:consensus} gives the precise form of that
continuity and proves a counter-intuitive consequence: the governing variable of consensus
reliability is not the number of participants but the independence of their errors.

\section{Contract operads and composition}\label{sec:operad}

Proposition \ref{prop:soundness} gives sufficient conditions for composition, but it lists
them one by one. This section supplies the algebraic structure of contracts, so that the
composition conditions become well-definedness conditions on that structure rather than a
checklist.

\begin{definition}[Contract operad]\label{def:operad}
A \term{contract operad} $\mathcal{O}$ consists of:
\begin{enumerate}[label=(\roman*)]
\item a set of \term{colours} $\mathrm{Col}$, whose elements are behaviour contracts
      (Definition \ref{def:contract});
\item for each $n \ge 0$ and colours $S_0, S_1, \dots, S_n$, a set of \term{operations}
      $\mathcal{O}(S_1, \dots, S_n;\, S_0)$ whose elements are \emph{composites taking
      $S_1,\dots,S_n$ as subcomponent contracts and satisfying $S_0$};
\item \term{substitution}: from $f \in \mathcal{O}(S_1,\dots,S_n;\,S_0)$ and
      $g_i \in \mathcal{O}(T_{i,1},\dots,T_{i,m_i};\,S_i)$ for $i = 1..n$,
      \[
      f \circ (g_1, \dots, g_n) \;\in\;
      \mathcal{O}\bigl(T_{1,1},\dots,T_{1,m_1},\dots,T_{n,1},\dots,T_{n,m_n};\,S_0\bigr);
      \]
\item \term{units} $\mathrm{id}_S \in \mathcal{O}(S;\,S)$.
\end{enumerate}
Substitution is required to be associative and unital. Informally, a colour is what a part
promises, an operation is how parts are assembled, and substitution refines a part without
altering the promise of the whole. The definition is the standard one for an operad,
coloured by contracts \cite{leinster2004}.
\end{definition}

\begin{definition}[Interface graph and freeness]\label{def:interface-graph}
Let $f \in \mathcal{O}(S_1,\dots,S_n;\,S_0)$. Take the subcomponents as vertices and draw a
directed edge $u \to v$ when the guarantee of $u$ is required by an assumption of $v$. The
resulting digraph is the \term{interface graph} of $f$. If it is acyclic, $f$ is said to be
\term{freely generated} by its interface graph.
\end{definition}

\begin{proposition}[Structural form of the composition conditions]\label{prop:operad-compose}
Conditions (a) and (b) of Proposition \ref{prop:soundness} are respectively equivalent to
two structural conditions on the contract operad:
\begin{enumerate}[label=(\alph*)]
\item \term{Well-typed colouring}: observation compatibility (condition (a) of
      Proposition \ref{prop:soundness}) holds if and only if substitution is well defined
      on colours, that is, the composite never demands of a subcomponent an observation its
      colour does not declare;
\item \term{Environmental closure}: assumption closure (condition (b) of
      Proposition \ref{prop:soundness}) holds if and only if the interface graph is
      acyclic, that is, the composite can be built uniquely bottom-up from its
      subcomponents.
\end{enumerate}
\end{proposition}

\begin{proof}
For (a): if the observation maps are compatible, the projection of each subcomponent onto
the interface agrees with the composite's observation, so substitution introduces no new
colour requirement; conversely, if substitution is well typed on colours, every observable
of the composite is traceable to an observation declared by some subcomponent colour, which
is observation compatibility.

For (b): if the interface graph is acyclic, a topological order exists and substitution
along that order yields the unique construction, with assumption closure verified level by
level. Conversely, if there is a cycle $u_1 \to u_2 \to \cdots \to u_k \to u_1$, then the
assumption of $u_1$ depends on the guarantee of $u_k$, whose guarantee depends (around the
cycle) on the assumption of $u_1$; the two sides of the substitution define each other and
the composite has no well-founded construction. To give it meaning one must supply a
fixed-point semantics, that is, a separate existence-and-uniqueness condition for that
cycle.
\end{proof}

Part (b) accounts for two of the three composition failures in \S\ref{sec:composition}, and
explains why two established practices must introduce an additional hypothesis:

\begin{itemize}
\item contract composition for dynamical systems must assume that the interconnection is
      \emph{well posed}, which is exactly the fixed-point semantics supplied for a feedback
      cycle;
\item software contracts fail on network-retry loops because retry turns an acyclic call
      graph into a cyclic interface graph, and the colouring is not re-verified once the
      cycle appears.
\end{itemize}

\begin{remark}[Feedback is a structural change]\label{rmk:feedback}
Proposition \ref{prop:operad-compose} reduces composability to acyclicity of the interface
graph. This has a non-trivial consequence: \emph{introducing feedback is not a local
change}. Any edit that turns an acyclic interface graph into a cyclic one---retry, cache
write-back, bidirectional synchronisation, an event loop---invalidates the existing
composition argument, regardless of whether the components remain individually correct.
This agrees with engineering experience: such edits are a high-incidence region for
regressions.
\end{remark}

\section{Notation}\label{sec:notation}

\begin{longtable}{L{3.2cm}L{11.3cm}}
\caption{Notation used throughout.}\label{tab:notation}\\
\toprule
\textbf{Symbol} & \textbf{Meaning} \\
\midrule
\endfirsthead
\toprule
\textbf{Symbol} & \textbf{Meaning} \\
\midrule
\endhead
\bottomrule
\endlastfoot
$\Lambda_i$ & level: the six-tuple (Definition \ref{def:level}) \\
$\Obs_i \dashv \Conc_i$ & Galois connection between levels (Definition \ref{def:galois}) \\
$\Obl(\Lambda_i) = \Ker \Obs_i$ & oblivious set (Definition \ref{def:oblivious}) \\
$\kappa_i$ & compression ratio (Definition \ref{def:compression}) \\
$\mathcal{G}$ & architecture relational graph (Definition \ref{def:relgraph}) \\
$\phi(\mathcal{G})$ & abstract skeleton (Definition \ref{def:skeleton}) \\
$T = (Q,d)$ & task (Definition \ref{def:task}) \\
$\DL(M)$ & description length (Definition \ref{def:mdl}) \\
$L_T(M)$ & task loss (Definition \ref{def:mdl}) \\
$S$ & behaviour contract: the eight-tuple (Definition \ref{def:contract}) \\
$\Tr(B,E)$, $\Allow(S)$ & observable traces, permitted traces
  (Definition \ref{def:traces}) \\
$\mathfrak{A}$ & assurance (Definition \ref{def:assurance}) \\
$\TCB$ & trusted computing base (Definition \ref{def:tcb}) \\
$\mathcal{O}$, $\mathrm{Col}$ & contract operad and its colours (Definition \ref{def:operad}) \\
$R_i$ & inter-level interpretation relation (Remark \ref{rmk:source}) \\
$\rho(c)$ & read-set of change $c$ (Definition \ref{def:readset}) \\
$\mathcal{K}_6$ & the six functions (Definition \ref{def:six}) \\
$D$ & technical debt (Definition \ref{def:debt}) \\
$C(t)$, $K$ & governable complexity, comprehension capacity
  (Definition \ref{def:governable}) \\
$\mathcal{J}$ & observation operator (Definition \ref{def:observability}) \\
$\operatorname{rank}_{\mathcal{T}}\mathcal{J}$ & observability rank on the task subspace
  (Definition \ref{def:observability}) \\
$q$ & correlated error mode (\S\ref{sec:consensus}) \\
\end{longtable}

% =====================================================================
\chapter[Structural Convergence]{Structural Convergence: Over-Constrained Solution Sets and Engineering Attractors}\label{ch:convergence}
% =====================================================================

This chapter addresses Problem \ref{prob:low}. High-level structure does converge, but by a
different mechanism from the low level, and claims of different strength must be stated
separately; the strength of convergence depends on domain binding, and its carrier is
relations rather than components. The chapter closes with a falsifiable verification design.

\section{Three Convergence Claims of Different Strength}\label{sec:three-claims}

\begin{hypothesis}[The weak claim: recurrence of local patterns]\label{hyp:weak}
Software in different domains repeatedly adopts the same set of local structural solutions:
layering and staging, separation of mechanism from policy, separation of control from data,
events and state machines, separation of declarative description from execution engine,
separation of front end from back end, and resource ownership and lifecycle management.
What matters about these structures is not that they share a name, but that they share a
\emph{structural motivation}.
\end{hypothesis}

The weak claim has substantial engineering experience behind it, and it is also the weakest:
it is almost impossible to falsify, and therefore has limited explanatory power.

\begin{hypothesis}[The medium claim: structural similarity after abstraction]\label{hyp:medium}
For large software systems that come from different domains, are developed by different
teams and carry different implementation details, once differences of implementation
language, component naming and domain semantics have been eliminated, their high-level
architectures exhibit a structural similarity \emph{significantly above a random baseline}
in some definable abstract space.
\end{hypothesis}

The medium claim is the core claim, and the operationalisable one. It can be operationalised
(\S\ref{sec:baseline}), and whether it holds depends on no metaphysical assumption.

\begin{hypothesis}[The strong claim: single-point convergence]\label{hyp:strong}
Almost all large software that is normally biased eventually converges on one and the same
high-level architecture.
\end{hypothesis}

The strong claim is unsupported by the evidence and is probably too strong. Its principal
defect is that it conflates ``there exists an attractor'' with ``there is only one
attractor''.

This paper takes Hypothesis \ref{hyp:medium} (the medium claim) as its core claim, treats
Hypothesis \ref{hyp:weak} (the weak claim) as its degenerate form, and rejects
Hypothesis \ref{hyp:strong} (the strong claim) as an overstatement.

\section{Constraint Systems and Solution-Set Geometry}\label{sec:overconstrain}

\begin{definition}[Constraint system and solution set]\label{def:constraint}
Let the design space be a $d$-dimensional manifold $\mathcal{D}$. A design problem is
defined by $c$ independent constraints
$g_1, \dots, g_c : \mathcal{D} \to \mathbb{R}$, and its \term{solution set} is
$\mathcal{S} = \{x \in \mathcal{D} : g_j(x) = 0,\; j = 1..c\}$.
\end{definition}

\begin{proposition}[Over- and under-constraint]\label{prop:overconstrain}
Under the notation of Definition \ref{def:constraint}, if the constraints are mutually
independent and transversal, then
\begin{enumerate}[label=(\alph*)]
\item when $c \ge d$, $\mathcal{S}$ is generically a finite set (or empty);
\item when $c < d$, $\mathcal{S}$ is generically a submanifold of dimension
      $\ge d - c$;
\item when $c \ll d$, the dimension of $\mathcal{S}$ tends to $d$, that is, the solution
      set is almost unconstrained.
\end{enumerate}
\end{proposition}

\begin{proof}
Under transversality, each independent constraint lowers the codimension by $1$. By the
transversality theorem (Thom), codimension is additive:
$\mathrm{codim}\,\mathcal{S} = c$, hence $\dim \mathcal{S} = d - c$ (when $d-c \ge 0$);
when $c \ge d$ the dimension is $0$, that is, a discrete point set; when $c < d$ the
dimension is positive, that is, a submanifold.
\end{proof}

Proposition \ref{prop:overconstrain} gives a clean explanation of ``unique below, manifold
above'', but the wording must be handled with care. It does \emph{not} say that the low
level really is formalised as a constraint system; it says that the effective dimension of
the low-level design space is small and its effective constraints numerous, so that the
solution set is degenerate in the geometric sense.

\begin{example}[The low level]\label{ex:abi}
Consider defining a function calling convention for a fixed instruction-set architecture.
The design freedoms include where arguments are passed (registers or stack), who is
responsible for saving registers, stack frame layout, where the return address is stored,
how variadic arguments are passed, alignment requirements, and compatibility with existing
binaries, debuggers, exception unwind tables, the system-call interface and the dynamic
linker. The great majority of these constraints are not negotiable preferences but are
fixed by \emph{objects that already exist}: existing binaries must keep working, the
debugger must be able to walk the stack, exceptions must unwind, and the kernel must be
able to receive system calls. Hence $c \ge d$, and the solution set degenerates to very few
feasible points, in practice often to a single point.
\end{example}

\begin{example}[The high level]\label{ex:arch}
Consider choosing an architecture for a task-processing system. The design freedoms
include: function calls within a monolith, multi-process message passing, microservice RPC,
event-driven asynchronous processing, the actor model, distributed workflow engines, and
their continuous parameters along the following dimensions: deployment granularity, state
ownership, consistency model, synchronous or asynchronous, failure boundaries, and data
partitioning. Here the problem-level constraints are markedly fewer than the degrees of
freedom, and $\mathcal{S}$ is a high-dimensional manifold.
\end{example}

Why, then, does the high level still converge? Because there is a \emph{second set of
constraints}, which come not from the problem but from the human organisations that build
and maintain the system:

\begin{principle}[Compression of the solution set by organisational constraints]\label{pr:org}
The under-constrained high-dimensional solution manifold at the problem level is compressed
by organisational constraints into a few sharp peaks. These constraints include: the cost of
communication between teams, the ceiling on cognitive load, the limits of deployment and
operations capacity, the market for technical talent, the sunk cost of an existing
technology stack, and regulatory and compliance requirements.
\end{principle}

One corollary follows: \emph{high-level convergence is not towards engineering optimality
but towards structures navigable by a bounded organisation}. The strength law below makes
this precise.

\section{Mechanisms of Convergence}\label{sec:mechanisms}

\begin{enumerate}[label=(M\arabic*)]
\item \term{Complexity constraint}. Let a system have $N$ components. If dependencies are
      left unorganised, the number of potential directed dependencies reaches $N(N-1)$, and
      the undirected count $N(N-1)/2$. Maintainability therefore requires structure to be
      imposed on dependencies: modularisation, layering, explicit interfaces, constraints on
      dependency direction, subsystem boundaries. These techniques are not identical, but
      all of them solve one and the same problem: \emph{limiting the propagation of
      complexity}. The classical argument for modularisation and information hiding is
      \cite{parnas1972}.
\item \term{Information-flow constraint}. A large class of software can be abstracted as
      ``input $\to$ processing $\to$ state change $\to$ output'' (with feedback, concurrent
      events, external side effects and continuously running state machines superimposed).
      Where different duties have different rates of change, error modes and resource
      requirements, separating them yields a concrete benefit: parsing from execution,
      business rules from data access, pure computation from side effects, and the control
      plane from the data plane.
\item \term{State-management constraint}. A system must answer: who owns a given state? who
      may modify it? when is it created and destroyed? how do several components observe its
      changes? how is concurrent access coordinated? how is consistency restored on failure?
      These questions are highly general, and therefore repeatedly produce similar
      structures: transaction and lock management, process and handle management, the state
      owner and event loop of a GUI, replicated state machines, entity lifecycle and system
      update.
\item \term{Cognitive constraint}. Large software is constrained not only by machines but by
      the limits of developer comprehension. If implicit relations hold between all
      components, then even a system that runs in the computational sense leaves its
      developers unable to predict the effect of a change. Humans therefore tend to build
      hierarchies, namespaces, modules, explicit interfaces, composable abstractions and
      locally comprehensible state. \emph{Software architecture organises not only
      computation but comprehensibility under bounded cognition.}
\item \term{Evolutionary constraint}. Large software is not designed in one pass; it passes
      through initial implementation, feature growth, performance optimisation, fault repair,
      team handover, module refactoring and compatibility maintenance. Some decisions reveal
      their cost only as evolution proceeds, which produces a selection effect. Evolution in
      practice does not guarantee optimality: legacy baggage, compatibility requirements,
      organisational structure and sunk cost can all keep an undesirable architecture alive
      for a long time.
\end{enumerate}

Mechanisms (M1)--(M3) come from the objective structure of the problem; mechanisms (M4)--(M5)
come from humans and history. This distinction corresponds to the strength law of
\S\ref{sec:strength}: the first three mechanisms produce cross-domain similarity, the last
two produce within-domain similarity and path dependence.

The classical statements of mechanisms (M4) and (M5) are decomposition into levels under
cognitive load \cite{simon1962} and the determination of system structure by communication
structure \cite{conway1968}. The oscillation between a ``no final state'' view of
architecture and multiple stable forms is treated in \cite{nygard2018}; the practice by
which organisational shape drives architecture back through team topologies appears in
bounded contexts and domain-driven design \cite{evans2003}.

\section{Convergence Strength and Domain Binding}\label{sec:strength}

\begin{definition}[Domain binding]\label{def:dom}
Let the semantics of level $\Lambda$ depend on a set of domain axioms
$\mathrm{Ax}(\Lambda)$ (for example, ``an account balance may be negative'' or ``a frame
must be presented within 16.6\,ms''). Define the \term{domain binding} of $\Lambda$ as
\begin{equation}\label{eq:dom}
\mathrm{Dom}(\Lambda) \;=\; \frac{\bigl|\mathrm{Ax}(\Lambda)\bigr|}
{\bigl|\mathrm{Ax}(\Lambda)\bigr| + \bigl|\mathrm{Ax}_{\text{general}}(\Lambda)\bigr|} ,
\end{equation}
that is, the fraction of the level's axioms that are domain-specific. $\mathrm{Dom} = 0$
means that the level is determined entirely by general axioms (such as memory locality or
the cost of concurrency coordination); $\mathrm{Dom} = 1$ means that the level is entirely
domain convention.
\end{definition}

\begin{conjecture}[The strength law of convergence]\label{conj:strength}
Let $\mathrm{Conv}(\Lambda)$ be the structural similarity of level $\Lambda$ under
cross-system comparison (\S\ref{sec:baseline} gives the operationalisation). Then
$\mathrm{Conv}$ is monotonically decreasing in $\mathrm{Dom}$:
\begin{equation}\label{eq:convlaw}
\Lambda_1 \preceq \Lambda_2 \;\Longrightarrow\;
\mathrm{Conv}(\Lambda_1) \;\ge\; \mathrm{Conv}(\Lambda_2)
\qquad \text{whenever} \qquad
\mathrm{Dom}(\Lambda_1) \le \mathrm{Dom}(\Lambda_2).
\end{equation}
\end{conjecture}

Conjecture \ref{conj:strength} has observable qualitative support: the closer to the
physical and network layers, the stronger the convergence (the eight fallacies of
distributed computing are hard constraints with no room for negotiation); the closer to the
business domain, the weaker the convergence (an insurance core system and a recommender
system are almost never isomorphic in their domain models); the closer to organisational
structure, the more it is governed by communication structure; the closer to the maturity
phase of a technology, the more its form converges; and the closer to core and high-risk
modules, the more conservatively it converges.

This paper does not call Equation \eqref{eq:convlaw} a law, because it depends on the
definability of $\mathrm{Dom}$ and the measurability of $\mathrm{Conv}$, both of which must
be established experimentally. Its value is to convert a vague directional judgement into a
falsifiable monotonicity hypothesis.

\section{Relations Rather Than Components}\label{sec:relations-converge}

\begin{principle}[Relations first]\label{pr:relations}
Cross-system structural commonality resides mainly in \emph{relations} rather than in
\emph{components}. Different systems need not have the same modules, but they may have
similar relational graphs:
\begin{itemize}
\item which parts decide and which parts execute;
\item which parts own state and which parts may only read it;
\item which parts define stable contracts and which parts are allowed to change;
\item which parts coordinate globally and which parts work independently;
\item which dependencies may cross a boundary and which must be forbidden.
\end{itemize}
\end{principle}

In the language of Definitions \ref{def:relgraph} and \ref{def:skeleton}: what must be
sought is not a universal architectural template but a set of structural relations that
recur across different systems, together with the \emph{conditions under which they hold}.
This is the same line of thought as the six functions of Chapter \ref{ch:infrastructure},
viewed from two sides: the six functions supply coordinates for \emph{management duties},
the relational graph supplies coordinates \emph{between duties}.

\begin{example}[Abstract skeletons of a database and a browser]\label{ex:db-browser}
A database engine comprises an SQL parser, a query optimiser, an execution engine, a buffer
pool, a transaction manager and a storage engine; a browser comprises HTML/CSS/JavaScript
parsers, a style and layout engine, a script execution engine, a rendering engine, a network
subsystem and a cache persistence mechanism. Their module names, internal algorithms and
concrete duties all differ. Yet under a suitable skeleton map $\phi$, both present as: an
\emph{input and interface layer} (receiving requests, parsing input, establishing an
execution context); an \emph{interpretation and planning layer} (converting external intent
into an internally executable representation); an \emph{execution and coordination layer}
(scheduling work, coordinating subsystems, managing the execution process); and a
\emph{state and resource layer} (managing data, caches, lifecycle and persistence).
\end{example}

The value of Example \ref{ex:db-browser} is limited: this skeleton is already broad enough
to describe almost any system. An additional requirement on the abstraction level is
therefore needed:

\begin{principle}[Two-way constraint on the abstraction level]\label{pr:abstraction-level}
An abstraction level worth studying must satisfy two directions at once: \emph{upwards}, it
must be able to identify structure common across systems; \emph{downwards}, it must retain
enough discriminating power that different systems remain distinguishable at that level.
The more abstract the level, the more implementation differences are eliminated, but the
more easily the description loses discriminating power; there is a proper middle level.
\end{principle}

The coordinates of the six functions in Chapter \ref{ch:infrastructure} are designed
precisely to satisfy this two-way requirement: they are narrower than
``input--processing--output'' (and therefore discriminating) and broader than any concrete
module list (and therefore comparable across domains).

\section{Verification Design}\label{sec:baseline}

For the medium claim (Hypothesis \ref{hyp:medium}) to be a scientific claim, it must be able
to answer ``compared with what''. A complete test design has five steps.

\paragraph{Step 1: Build the architecture graph.}
For each system, build a graph $\mathcal{G} = (V, E)$ whose nodes are components or duties
and whose edges are calls, dependencies, control flow or data flow. By
Definition \ref{def:relgraph}, \emph{relations of different types must be recorded
separately}, so that entirely different relations are not conflated.

\paragraph{Step 2: Define the abstraction map.}
Map concrete components to more general categories, for example: input and parsing, planning
and decision, execution and coordination, state and persistence, communication and
boundary, resource and lifecycle management. These categories must not be defined so broadly
in advance that any system will appear similar.

\paragraph{Step 3: Choose discriminating structural metrics.}
Candidate metrics include: the degree of modularity of the dependency graph; the depth of
the hierarchy; the proportion of cross-module dependencies; the concentration of state
ownership; the degree of coupling between control flow and data flow; the network centrality
of core components; and the direction of dependencies between duties.

\paragraph{Step 4: Construct a sound control group.}
It is not enough to compare a few well-known systems and find that all of them are modular.
The observed similarity must be compared against an appropriate random baseline, controlling
for system size, domain, platform, architectural style and historical technology stack.

\paragraph{Step 5: Distinguish genuine convergence from common origin.}
Two systems may adopt the same architectural pattern because they independently solved
similar problems; because they used the same framework; because they were shaped by the same
standard; because they drew on the same design literature; or because they inherited the
same historical architecture. \emph{Only after these factors are controlled is there better
reason to speak of independent structural convergence.}

Chapter \ref{ch:program} writes this design out as Hypotheses $H_1$ and $H_2$.

\section{Counterexamples and Limits}\label{sec:limits}

The cases that invalidate or degrade the convergence claim must be listed explicitly:

\begin{enumerate}[label=(\alph*)]
\item \term{Path dependence}. Once a structure is locked in by an ecosystem, a toolchain or
      a standard, it may persist for a long time even if suboptimal. The strength of the
      lock-in is proportional to the cost of migration and independent of the quality of the
      structure locked in.
\item \term{Phase transition in size}. At the two sizes of $10^{2}$ and $10^{9}$ entities,
      the optimal structure is qualitatively different: the consistency model, the sharding
      strategy and the failure model all change. There is therefore no such thing as ``the
      architecture of a large system'', only ``an architecture within some range of sizes''.
      This directly limits any statement that ``all large systems converge''.
\item \term{Organisational politics}. The determination of system structure by communication
      structure (Conway's law) means that even when the technical problem is identical,
      different organisations will produce different system boundaries. This is both part of
      the reason for convergence (when organisational shapes converge) and part of the reason
      for divergence (when organisational shapes differ).
\item \term{Strongly coupled systems}. When the internal coupling of a system cannot be
      localised (for example, a tightly coupled real-time control system or a monolithic
      kernel), module boundaries may be forced rather than chosen, and comparing
      ``architecture'' itself becomes less meaningful.
\item \term{Arbitrariness in the choice of abstraction level}. The choice of metrics in
      Steps 3 and 4 itself contains researcher degrees of freedom. Different duty
      classification systems yield different similarity conclusions. This is the most fragile
      methodological link in this line of work and must be mitigated by pre-registering the
      metrics and by robustness tests across several classification systems.
\end{enumerate}

The falsifiable part of this chapter is Hypothesis \ref{hyp:medium} and
Conjecture \ref{conj:strength}: if, after controlling for size, domain, platform and
framework, cross-system structural similarity is not significantly above a random baseline,
or if similarity is not monotone in $\mathrm{Dom}$, then the central claim of this chapter is
weakened (Hypotheses $H_1$ and $H_2$, Chapter \ref{ch:program}). The weak claim
\ref{hyp:weak} and Principle \ref{pr:relations} are unaffected by this, because they are
interpretive frameworks rather than empirical assertions.

% =====================================================================
\chapter[Internal Infrastructuralisation]{Internal Infrastructuralisation: The Conditional Necessity of the Six Functions}\label{ch:infrastructure}
% =====================================================================

This chapter addresses Problem \ref{prob:scale}. The central claim is that any sufficiently
complex software system reinvents an operating system inside its own problem domain. The
cause is not that engineers imitate a kernel; it is that a set of semantic obligations
cannot be evaded at scale. The chapter gives the premises, the proof and the criteria for
that claim, together with a falsifiable test that can be carried out inside a single
project.

\section{The Phenomenon and the Terms}\label{sec:phenomenon}

The observed phenomenon is this: once a system grows past a certain size, it begins to
promote mechanisms that were previously scattered across it into unified system-level
capabilities. Specifically:

\begin{itemize}
\item several facilities need to allocate and release resources;
\item several subsystems need to manage object lifetime;
\item several components need to execute asynchronous tasks;
\item several modules need to share state and coordinate modifications;
\item several facilities need unified error handling, logging and diagnostics;
\item several execution paths need to obey the same thread-safety and concurrency rules.
\end{itemize}

If each module solved these problems on its own, the system would accumulate a mass of
duplicated implementations, inconsistent behaviour and coupling that is hard to control. The
system therefore begins to handle them centrally, producing common mechanisms, interfaces,
managers, runtimes and execution models.

\begin{definition}[Internal infrastructuralisation]\label{def:infra}
\term{Internal infrastructuralisation} is a system's passage through three successive
stages:
\begin{enumerate}[label=(\roman*)]
\item \term{commonality identification}: recurrent requirements are identified as one class
      of problem (a merger in problem space);
\item \term{mechanism unification}: the commonality is promoted to a common capability, so
      that different facilities access shared resources through common interfaces and obey
      common lifecycle, error-handling and concurrency rules (centralisation and
      standardisation of mechanism);
\item \term{structural unification}: upper-layer facilities no longer solve all lower-level
      problems directly, but build on unified rules (unification of structure).
\end{enumerate}
In shorthand,
{\small$\boxed{\text{commonality identification} \to \text{mechanism unification} \to
\text{upper-layer composition}}$}.
\end{definition}

Mechanism unification does \emph{not} mean that all logic is gathered into one central
module. A mature architecture can achieve unification \emph{through protocols, interfaces
and distributed ownership}. What is unified is the \emph{rules and contracts}, not
necessarily the \emph{implementations and ownership}. That distinction is the key to
avoiding a god object, and Chapter \ref{ch:design} states it as a design principle.

\section{Participant--Resource Structure}\label{sec:participant}

\begin{definition}[Participant--resource structure]\label{def:pr-structure}
A \term{participant--resource structure} is a five-tuple
\begin{equation}\label{eq:pr}
\Pi = \bigl(E,\; \mathcal{R},\; \to,\; \mathrm{Ready},\; \mathrm{Fault}\bigr),
\end{equation}
where $E$ is a set of concurrently live entities,
$\mathcal{R} = \{r_1,\dots,r_m\}$ is a finite set of shared resources, $\to$ is the
transition relation of the operational semantics,
$\mathrm{Ready} \subseteq E \times \mathrm{Conf}$ indicates which entities may execute in a
given configuration, and $\mathrm{Fault}$ is a set of fault events.
\end{definition}

The two quantities that matter are $n = |E|$ and $m = |\mathcal{R}|$. Section
\ref{sec:integration} proves that the independent variable deciding whether a unified
infrastructure is needed is the structural complexity $n \times m$, not lines of code.

\begin{definition}[Unified substrate]\label{def:substrate}
A \term{unified substrate} for $\Pi$ is a set of mechanisms $\mathcal{U}$ such that the
interaction of every entity $e \in E$ with every resource $r \in \mathcal{R}$ proceeds
through an interface declared by $\mathcal{U}$ rather than being implemented directly by
$e$.
\end{definition}

\section{The Integration Cost Theorem and the Phase-Transition Threshold}\label{sec:integration}

\begin{theorem}[Integration cost]\label{th:integration}
Let the participant--resource structure $\Pi$ have $n$ entities and $m$ shared resources.
Assume:
\begin{enumerate}[label=(\alph*)]
\item without a substrate, every entity that participates in a resource $r$ must implement
      correct interaction with $r$ independently, including independent handling of the
      invariants of $r$;
\item under maximal coupling, every entity participates in all $m$ resources;
\item with a substrate, every entity uses at most $k$ substrate interfaces, and the
      substrate itself defines the semantics of each of the $m$ resources once; the fixed
      cost of the substrate is $F$ (design, implementation and long-term maintenance).
\end{enumerate}
Then the number of interacting pairs is $\Theta(nm)$ without a substrate and
$\mathcal{O}(nk + m)$ with one. The substrate is better than no substrate in expectation if
and only if
\begin{equation}\label{eq:threshold}
n \cdot m \;>\; c\,(n k + m) \;+\; F
\end{equation}
where $c$ is the unit cost coefficient of one interaction.
\end{theorem}

\begin{proof}
Without a substrate: correct interaction of entity $e$ with resource $r$ requires an
independent implementation, so the number of interacting pairs is
$\sum_{e} |\mathcal{R}(e)| \le nm$, with equality under assumption (b), that is
$\Theta(nm)$. In addition, for each resource $r$ the $n$ independent implementations must
agree on the invariants of $r$; this is a consistency obligation whose number of failure
modes is $\Omega(n)$ (any single implementation that deviates breaks the global invariant).

With a substrate: the semantics of the $m$ resources are defined once each, at cost
$\Theta(m)$; every entity interacts with the substrate through at most $k$ interfaces, at
cost $\mathcal{O}(nk)$; add the fixed cost $F$. The total is therefore
$\mathcal{O}(nk + m) + F$.

Comparing the two expressions and normalising by the unit cost coefficient $c$ gives
\eqref{eq:threshold}.
\end{proof}

\begin{corollary}[Phase-transition threshold]\label{cor:threshold}
Because $k$ is the \emph{bounded} interface usage of each entity, $\mathcal{O}(nk + m)$ is
linear in $n$, whereas $\Theta(nm)$ is a product linear in both $n$ and $m$. Hence when
either $n$ or $m$ grows substantially, inequality \eqref{eq:threshold} must hold. The phase
transition occurs at $n m \sim F/c$, that is, when the \emph{product of the number of
entities and the number of resources} exceeds the order of the fixed cost of the substrate.
\end{corollary}

Corollary \ref{cor:threshold} makes precise the intuition that ``the code volume reaches
some order of magnitude'', and supplies its correct form:

\begin{principle}[The threshold is not lines of code]\label{pr:notloc}
The independent variable deciding whether a unified substrate must be built is the product
of the \emph{number of live entities} and the \emph{number of shared resources}, together
with the \emph{frequency of recurrent problems}, the \emph{number of interactions between
state and resources}, the \emph{complexity of concurrent execution}, the \emph{rate of
module evolution}, and the \emph{demands on consistency, performance and reliability}.
Lines of code are only a proxy for these quantities.
\end{principle}

This explains a fact that ought to be puzzling: two programs of a million lines each may
have entirely different internal architectures. One may be a large body of mutually
independent generated code, data files or linear business logic ($m$ small, $n$ small); the
other may have complex shared state, asynchronous tasks and many interdependent subsystems
($n$ and $m$ both large). Only the second develops a need for a unified runtime, resource
management and coordination mechanisms.

The central reduction (Section \ref{sec:participant}) supplies a second explanation here:
when the number of live entities in a system jumps from $10^{2}$ to $10^{6}$, naive
management undergoes a \emph{phase transition}. For 100 objects, conditional branching and
raw allocation suffice; for a million concurrent tasks, cache entries or connections, the
potential interactions between entities are $O(N^{2})$ and the complexity is infeasible, so
pooling, indexing, layering, asynchrony, batching and cache locality are forced in. That set
of techniques is exactly the full inventory of ``kernel technology''.

\begin{example}\label{ex:pg-memctx}
A relational database engine builds its own memory-context system because a
general-purpose allocator fragments and cannot track leaks under a load that ``creates and
destroys thousands of tuples per query''. It therefore constructs allocation and bulk
reclamation over a tree of contexts, which is an object pool inside the database domain. An
asynchronous runtime builds its own work-stealing scheduler because operating-system threads
are too heavy in the face of a million concurrent operations, and therefore constructs
$M{:}N$ scheduling, which is user-space time-sharing. A jump in order of magnitude forces
kernel technology in: this is the first necessity.
\end{example}

\section{The Six Functions}\label{sec:six}

\begin{definition}[The six functions]\label{def:six}
For a participant--resource structure $\Pi$, the \term{six functions} $\mathcal{K}_6$ are
defined as
\begin{enumerate}[label=(F\arabic*)]
\item \term{Identity and lifecycle}: the identity, creation, liveness and termination
      relations of entities;
\item \term{Allocation and reclamation}: the allocation and reclamation policies for finite
      resources;
\item \term{Scheduling}: the selection function that chooses the next executor among ready
      entities;
\item \term{Naming and addressing}: the mechanisms that make entities discoverable,
      referenceable and traversable;
\item \term{Communication and synchronisation}: the interaction primitives between entities
      and the composition rules for concurrent transitions;
\item \term{Fault isolation and recovery}: the containment relations of faults and the
      replayability of state.
\end{enumerate}
\end{definition}

\begin{theorem}[Conditional necessity of the six functions]\label{th:six}
Let $\Pi = (E, \mathcal{R}, \to, \mathrm{Ready}, \mathrm{Fault})$ be a
participant--resource structure whose operational semantics is well defined (that is, in
every reachable configuration the next transition is determined, or the non-determinism it
permits is explicitly specified). If $|E| > 1$ and $\mathcal{R}$ is finite and contended,
then the six functions of $\mathcal{K}_6$ have necessarily already been defined, either in
the system's own code or in its host (operating system, language runtime, framework).
\end{theorem}

\begin{proof}[Structural argument]
It is shown item by item that the absence of any single function renders the semantics
incomplete (an undefined behaviour exists).

\emph{(F1) Identity and lifecycle}: if entities were indistinguishable, $\mathrm{Ready}$
could not be defined as a function on the set of entities; if entities had no creation and
termination relations, the variation of $E$ over time would be unspecified and $\to$ would
be undefined at steps where ``an entity appears or disappears''.

\emph{(F2) Allocation and reclamation}: $\mathcal{R}$ is finite and contended, so ``more
resources are requested at once than can be supplied'' is a reachable state. Without an
allocation policy, $\to$ is undefined in that state. If the resources of terminated entities
are not reclaimed, available resources decrease monotonically and the system reaches
resource exhaustion in finitely many steps; $\to$ is then undefined in the forward-looking
sense.

\emph{(F3) Scheduling}: if $|E| > 1$ and several entities are ready at once, $\to$ is no
longer a function but a relation. Without a specified choice among ready entities, the
``permitted non-determinism'' is not delimited and the semantics is incomplete.

\emph{(F4) Naming and addressing}: an interaction primitive (see F5) must designate the
object of the interaction. Without a referencing mechanism, the object of an interaction
cannot be expressed.

\emph{(F5) Communication and synchronisation}: if entities interact, the global transition
is not the simple product of local transitions; without composition rules, the semantics of
concurrent transitions is undefined.

\emph{(F6) Fault isolation and recovery}: $\mathrm{Fault} \ne \varnothing$ (faults are a
physical and engineering given). Without containment relations, a local deviation has global
semantics, which is equivalent to the system having no well-defined behaviour under fault;
without recovery rules, the system never returns to a specified state after the first fault.

In sum, the absence of any one function destroys the premise of well-definedness. \qed
\end{proof}

\begin{remark}[The nature of the theorem]\label{rmk:six-nature}
Theorem \ref{th:six} is a \emph{conditional necessity} result, not an unconditional theorem.
Its conditions are $|E| > 1$, $\mathcal{R}$ finite and contended, and faults reachable. Its
proof is a structural enumeration argument, not a derivation in some formal system. It can
be falsified by exhibiting a well-defined system that satisfies the premises and lacks one
of the six. That judgement is itself open to challenge.
\end{remark}

The wording of Theorem \ref{th:six}, ``either in the system's own code or in its host'', is
essential. It separates two questions: whether the functions necessarily appear, and whether
they must be built in-house.

\section{Build or Adopt}\label{sec:buildadopt}

\begin{corollary}[Build-or-adopt dichotomy]\label{cor:buildadopt}
By Theorem \ref{th:six} the six functions must be defined, but where they are defined is a
design choice. That choice is decided by the \emph{domain criticality} of the invariants:
\begin{enumerate}[label=(\alph*)]
\item if the invariant carried by a function is \emph{domain-critical}, that is, the
      invariant can hold only under a guarantee owned by this system, the function must be
      built in-house;
\item if the invariant carried by a function is \emph{general}, that is, a general-purpose
      host can provide a guarantee of sufficient strength, the function should be adopted;
\item implementing a general function in-house produces the failure mode this paper calls
      \emph{reinvention of commodity parts}: a worse general-purpose mechanism
      re-implemented at a higher long-term cost of comprehension and maintenance.
\end{enumerate}
\end{corollary}

This criterion explains why the boundary between ``application'' and ``platform'' falls
exactly where it does: \emph{the set of invariants you are forced to own is the boundary of
your kernel}.

\begin{example}\label{ex:buildadopt}
A game engine must build its own entity--component storage and frame scheduling, because a
general-purpose host cannot provide the guarantee of a ``sub-millisecond deterministic frame
budget''. A database must build its own storage engine and write-ahead log, because a
general-purpose file system cannot provide the guarantee of ``recovery to a deterministic
state after a crash''. Conversely, a web application should adopt garbage collection, thread
scheduling and memory allocation, because the invariants of these functions are general and
building them in-house yields only a worse version.
\end{example}

\begin{remark}[A sharper criterion: monotonicity]\label{rmk:calm}
Corollary \ref{cor:buildadopt} decides build versus adopt by whether the invariant is
domain-critical. For functions of the ``coordination'' kind there is a sharper criterion that
can be checked statically: \emph{whether the computation is monotone}. Where monotonicity
holds \cite{calm2011} (outputs grow with inputs and do not depend on negation or on
aggregate global knowledge), coordination is not required and can be pushed down to the
participants; non-monotonicity (uniqueness constraints, negation, threshold aggregation,
non-commutative updates) forces some site to carry coordination.

This sharpens the claim for function (F5) (communication and synchronisation) and part of
(F3) (scheduling) of Section \ref{sec:six} from ``centralise once the scale is reached'' to
``centralise once non-monotonicity appears''. For monotone subsystems the $k$ of the
integration cost theorem (Theorem \ref{th:integration}) can be reduced to nearly a constant,
because participants need not agree on shared state.
\end{remark}

\section{Cross-Domain Comparison}\label{sec:instances}

Theorem \ref{th:six} predicts the pattern of Table \ref{tab:organs}: each column is a system
from a different company, a different language, a different era and a different domain, while
each row is the same organ. The table does not claim that these mechanisms can be identified
one with another; it compares them along the same dimension of \emph{management duty}.

\begin{table}[htbp]
\centering
\scriptsize
\caption{Independent implementations of the same set of management duties in different domains}\label{tab:organs}
\begin{tabular}{L{1.5cm}L{2.4cm}L{2.4cm}L{2.2cm}L{2.1cm}L{2.3cm}L{2.4cm}}
\toprule
\textbf{Function} & \textbf{Linux} & \textbf{Chromium} & \textbf{Tokio} & \textbf{Bevy}
& \textbf{Unreal} & \textbf{PostgreSQL} \\
\midrule
Identity/\newline lifecycle & task\_struct, inode, file & Frame, Document, Node &
Task, JoinHandle & Entity, Component & UObject, Actor & backend, xid, tuple \\
Allocation/\newline reclamation & buddy, slab, page cache & PartitionAlloc, Oilpan &
slab-style task allocation & archetype column store & FMemory, GC &
MemoryContext, shared\_buffers \\
Scheduling & EEVDF, workqueue & task queue, RendererScheduler &
work-stealing, waker & Schedule/System parallel executor & TaskGraph, GameThread &
postmaster, autovacuum \\
Naming/\newline addressing & kobject, sysfs, fd & DOM, Web IDL & Handle, waker &
Entity ID & FName, reflection metadata & OID, system catalog \\
Communi\-cation/\newline synchroni\-sation & syscall, netlink, IPC & Mojo IPC &
mpsc, oneshot &
Event, Resource & delegates, replication & libpq, LWLock, extension hooks \\
Isolation/\newline recovery & namespace, cgroup, OOM & site isolation, renderer process sandbox &
panic propagation, cancellation & state machine isolation & crash reporter &
MVCC snapshots, WAL replay \\
\bottomrule
\end{tabular}
\end{table}

Conversely, the \emph{form} these six organs take in different hosts varies widely, and that
variation is the concrete content of ``the trunk converges while the branches diverge'' in
Chapter \ref{ch:convergence}. The sources of that variation in form can be listed
systematically:

\begin{enumerate}[label=(\alph*)]
\item \term{Host memory model}. Systems with manual memory management (C/C++/Rust) must
      build their own allocator and even their own reclaimer (engine object reclamation,
      Blink's garbage collection in the browser); systems with managed memory spend the same
      effort on pause behaviour, because allocation is given and the pain point becomes the
      tail of the latency distribution.
\item \term{Host concurrency model}. The runtime model determines the shape of the
      scheduling organ: a game engine uses frame-driven cooperative scheduling (predictable
      latency over throughput), a server event-driven preemptive scheduling (throughput
      first). Frame schedulers and work-stealing schedulers are therefore two subspecies of
      one organ.
\item \term{Granularity of domain primitives}. A general-purpose kernel offers pages,
      threads, file descriptors and event polling; but a game engine needs frames,
      components and assets, a database needs buffer pools, transactions and version chains,
      and a browser needs document nodes, security contexts and cross-process objects. A
      mismatch in primitive granularity forces each domain to build its own infrastructure.
\item \term{Evolutionary stage}. These internal operating systems are themselves
      continuously refactored, and each refactoring converges towards ``a more complete set
      of the six functions'' rather than diverging from it.
\end{enumerate}

Even the operating system itself is circumvented by its users: high-performance network and
storage stacks redo I/O by polling in user space, asynchronous I/O interfaces are redesigned,
and the operating system cannot escape being ``reinvented'' by its own user space. This is
not anomalous; it is a direct consequence of Corollary \ref{cor:threshold}: when the
$n \times m$ of user space crosses the threshold, user space grows a new kernel.

\section{Transport and Re-Derivation}\label{sec:transport}

Table \ref{tab:organs} shows the same set of functions recurring across domains. But
``recurrence'' admits two entirely different explanations, with entirely different
predictions:

\begin{enumerate}[label=(\alph*)]
\item \term{Isomorphism}: the six functions are six implementations of one structure and can
      be translated into one another;
\item \term{Transport along the base}: the six functions are the fibres of some more general
      structure over different \emph{problem shapes}, and when the shape changes the
      structure must be re-derived, not carried over.
\end{enumerate}

The original treatment gives only (a). Distinguishing the two requires making the base
explicit, namely ``problem shape''.

\begin{definition}[Problem shape and level family]\label{def:fibration}
A pair $\pi : \mathcal{L} \to \mathcal{P}$ is a \term{level family}, where
\begin{itemize}
\item $\mathcal{P}$ is a set of \term{problem shapes}, determined by the resource form,
      concurrency structure, failure modes and time scales of the domain;
\item $\mathcal{L}$ is a set of layer implementations;
\item the fibre $\pi^{-1}(p)$ is the family of layer implementations feasible at problem
      shape $p$.
\end{itemize}
$\pi$ maps ``an implementation'' to ``what shape of problem it solves''. Reindexing $\pi$
gives \term{transport along the base}: given $p \to p'$, an implementation at $p$ is moved
to $p'$.
\end{definition}

\begin{proposition}[What transport preserves]\label{prop:transport}
Let $p \to p'$ be a map between problem shapes. If the map is an \emph{equivalence} (the same
resource form, concurrency structure, failure modes and time scales), then the \emph{slots}
of $\mathcal{K}_6$ (the question each function must answer) are preserved under transport;
the form of the implementation may change, but the slots are neither added nor removed. If
the map is not an equivalence, the semantic requirement of at least one slot changes, and
the implementation must be re-derived.
\end{proposition}

\begin{proof}[Argument]
Each item of $\mathcal{K}_6$ is uniquely determined by the semantic obligation of
Theorem \ref{th:six}: identity by the distinguishability of entities, allocation by the
finiteness and contention of resources, scheduling by the cardinality of the ready set,
naming by the object of interaction, communication by the mode of composition of concurrent
transitions, and isolation by the propagation range of faults. These six obligations depend
only on the four components of problem shape. When the map preserves the four components,
the obligation set is unchanged; when any component changes, the semantic requirement of the
corresponding obligation changes with it.
\end{proof}

\begin{corollary}[An observable discriminator]\label{cor:transport-test}
By Proposition \ref{prop:transport}, differences in form fall into two classes:
\emph{implementation differences of the same slot over different bases} (for instance
frame-driven and event-driven schedulers, whose bases have different time scales), and
\emph{differences of slot}. The former is not a counterexample; the latter is.
\end{corollary}

Corollary \ref{cor:transport-test} makes the middle proposition of Chapter
\ref{ch:convergence} (Hypothesis \ref{hyp:medium}) testable more rigorously: two systems
must first be aligned on the base, and only then compared slot by slot. Comparing modules
directly without aligning the base is the most common error in existing comparative studies
of architecture.

\begin{remark}[The status of this section]\label{rmk:fibration-status}
Definition \ref{def:fibration} and Proposition \ref{prop:transport} are a \emph{structural
proposal}, not an application of an established theorem of category theory. Their role is to
make explicit a question that has not previously been posed, whether the recurrence is
isomorphism or transport, and to give a discriminating corollary
(Corollary \ref{cor:transport-test}). Promoting $\pi$ to a strict fibration and proving the
properties of the reindexing functor is unfinished work (item O in
Chapter \ref{ch:objections}).
\end{remark}

\section{The Coarse-Graining Criterion}\label{sec:coarsegrain}

Theorem \ref{th:six} establishes that the six functions necessarily \emph{exist}, but not
that they are the most \emph{salient} structure. A stronger proposition is that the six
functions are the \emph{surviving set} of system structure under coarse-graining.

\begin{definition}[Coarse-graining sequence]\label{def:coarsegrain}
Let the architecture relational graph of a system be $\mathcal{G}$
(Definition \ref{def:relgraph}). A \term{coarse-graining} is a quotient map
$q : \mathcal{G} \to q(\mathcal{G})$ that merges components while preserving relation types.
A sequence of quotient maps $q_1, q_2, \dots$ (each step strictly coarser than the last) is
a \term{coarse-graining sequence}. A structure $\mathcal{S}$ \term{survives} under that
sequence if there exists $N$ such that for all $n \ge N$, $\mathcal{S}$ remains recognisable
in $q_n(\mathcal{G})$.
\end{definition}

\begin{conjecture}[The six functions are the relevant structure]\label{conj:relevant}
Under the coarse-graining sequence of a normally biased large system, the structures that
survive are exactly the slots of $\mathcal{K}_6$ and their relational graph; the remaining
structures (concrete modules, algorithm choices, data-layout details) are washed out at a
sufficiently coarse granularity.
\end{conjecture}

The value of Conjecture \ref{conj:relevant} is that it yields a falsification procedure
\emph{inside a single project}: no cross-project statistical comparison is needed, only
repeated coarse-graining of one code base, observing the surviving set. If the surviving set
is not the six functions, the conjecture is refuted and the convergence thesis of
Chapter \ref{ch:convergence} requires a different explanation.

\begin{principle}[Form of implementation]\label{pr:coarsegrain-impl}
Coarse-graining can be implemented entirely mechanically: starting from a module dependency
graph, three operations are applied repeatedly: (i) merge nodes whose fan-in and fan-out
fall below a threshold; (ii) collapse transitive dependency edges; (iii) erase nodes that
carry no invariant ($\mathrm{S}(\iota) = \varnothing$ in Definition \ref{def:enforce}).
Each step records the change in the surviving structure. The whole experiment needs one
repository and one dependency graph, with no labels and no manual classification, and is
therefore unaffected by the ``freedom in choosing a duty-classification system'' noted in
Section \ref{sec:limits}.
\end{principle}

Chapter \ref{ch:program} states Conjecture \ref{conj:relevant} as Hypothesis $H_{10}$.

\begin{remark}[Relation to renormalisation]\label{rmk:rg}
The operations of Principle \ref{pr:coarsegrain-impl} are formally the same as
coarse-graining in the renormalisation group: short-range degrees of freedom are eliminated
and one observes whether long-range structure converges to a fixed point. The fixed point
here is the six-function slots. It must be stated plainly that this does \emph{not} mean
that software possesses a renormalisation-group flow: software has no separation of scales
(the discussion preceding Principle \ref{pr:vm}), hence no exponential decay of genuinely
irrelevant operators. What is the same is the \emph{operation}, not the \emph{mechanism}.
\end{remark}

\section{Virtualisation and Deep Modules}\label{sec:virtualisation}

The phrase ``compress complexity and provide the upper layers with a high-level language and
symbols'' can be made precise. The core work of every infrastructure layer is one and the
same action:

\begin{principle}[Virtualisation]\label{pr:vm}
An infrastructure layer translates what is \emph{physical, continuous and volatile} into
\emph{discrete, stable, composable} domain-level symbols.
\end{principle}

\begin{example}\label{ex:vm}
Operating system: physical memory $\to$ virtual address space; block devices $\to$ files;
hardware interrupts $\to$ signals.
Database: physical pages $\to$ multi-version tuples $\to$ ``rows'' in SQL.
Browser: byte streams $\to$ document trees; socket events $\to$ event-loop callbacks.
Game engine: raw bytes in memory $\to$ components and assets; the time of one frame $\to$ a
schedule table.
Asynchronous runtime: event-poll bitmap $\to$ the Future of \code{async}/\code{await}.
\end{example}

In the vocabulary of module design \cite{parnas1972,ousterhout2018}, these are all \emph{deep
modules}: minimal interface, very deep implementation. Behind one allocation call stands an
entire allocator system; behind one task creation stands an entire scheduler. The duty of
an infrastructure layer is to keep producing deep modules, sealing complexity below and
exposing clean symbols above.

This connects directly to the theory of comprehension in Chapter \ref{ch:comprehension}: a
deep module is precisely a level of very high compression ratio
(Definition \ref{def:compression}) that quotients away a large amount of lower-level
difference while retaining the predictive power the upper layer needs. It also explains why
good infrastructure is at once ``high-performance'' and ``easy to understand'': both are
consequences of the same quotient map.

\section{Functional Differentiation in Agent Systems}\label{sec:agent}

If Theorem \ref{th:six} holds, the most recently emerged complex computational domain should
be re-enacting this evolution. Table \ref{tab:agent} gives the observation.

\begin{table}[htbp]
\centering
\small
\caption{Agent systems rediscovering the six functions}\label{tab:agent}
\begin{tabular}{L{3.4cm}L{11.0cm}}
\toprule
\textbf{Function} & \textbf{Counterpart in the agent domain} \\
\midrule
Allocation and reclamation & the context window as fast storage; retrieval augmentation as
page-fault handling; external memory as swap space; context compression as memory
compaction \\
Communication and synchronisation & tool calls or function calls as system-call
primitives \\
Scheduling & the agent main loop, multi-agent orchestration \\
Identity and lifecycle & lifecycle management of sessions, tasks and subgoals \\
Naming and addressing & tool registration and discovery protocols \\
Isolation and recovery & sandboxing, retries, checkpoints and rollback \\
\bottomrule
\end{tabular}
\end{table}

The value of this observation is that it is a \emph{natural experiment}: no one is
deliberately imitating a kernel; complexity itself has pressed out the same anatomy. It also
forms the empirical basis of Hypothesis $H_3$ (Chapter \ref{ch:program}).

Theorems \ref{th:integration} and \ref{th:six} hold under their stated premises; what can be
challenged is the premises and the rigour of the structural argument. What is empirically
falsifiable is whether the appearance of the six functions is predicted by $n \times m$
rather than by lines of code (Hypotheses $H_3$ and $H_{10}$), and whether there exists a
well-defined system that satisfies the premises yet lacks one of the functions over the long
run.

% =====================================================================
\chapter{Comprehension: Inverse Compression}\label{ch:comprehension}
% =====================================================================

This chapter addresses the first half of Problem \ref{prob:mind}: for a large system that
cannot be read through from the micro level, at what level should it be understood, and how
does comprehension lead into development. The central claim is that comprehension is inverse
compression, whose optimal depth is given by the stationarity condition of a rate--distortion
objective and whose feasibility is decided by observability rank.

\section{Failure of Sequential Reading}\label{sec:why-sequential}

A natural way to read is to follow the physical order and the control flow of the source code:

\begin{quote}
\code{main()} $\to$ \code{module\_a} $\to$ \code{module\_b} $\to$
\code{function\_c} $\to$ \code{helper\_d} $\to$ \code{allocator} $\to$
\code{operating\_system} $\to \cdots$
\end{quote}

One then keeps following downwards and ends in infinite recursion. The failure is not a lack
of patience but structural:

\begin{principle}[Calling relations are not cognitive dependencies]\label{pr:call-vs-cog}
The function-call relation is the topology of \emph{implementation}; the cognitive-dependency
relation is the topology of \emph{comprehension}. The two are not isomorphic. In studying
transaction isolation one may need to know how the transaction manager uses locks and
snapshots, but not necessarily understand the entire virtual-memory subsystem of the operating
system; conversely, in investigating anomalous memory usage the high-level business semantics
of the transaction manager may be temporarily irrelevant.
\end{principle}

The more fundamental reason concerns the relation between the direction of generation and the
direction of comprehension. A system is designed \emph{top-down} (intent first, then
architecture, and only then the unfolding of code); the inverse process should therefore take
the same direction. Reading code bottom-up is like understanding a cat by counting all its
molecules: feasible in principle, absurd in practice.

The concession ``feasible in principle'' has a practical consequence: sequential reading is
not \emph{impossible} but \emph{dominated}. Proposition \ref{prop:stopping} gives its precise
form: beyond the stationary point, description length increases while task loss no longer
falls.

\section{Cognitive Maps}\label{sec:maps}

Confronted with a large monolithic repository, one should build a conceptual map before a file
map. The division of labour is:

\begin{itemize}
\item a file map answers ``where the code is'';
\item a conceptual map answers ``what the system is made of, why these parts exist, who
      depends on whom, and what rules allow them to collaborate''.
\end{itemize}

\begin{definition}[The four cognitive maps]\label{def:fourmaps}
For a system $S$ and a task $T$, the four cognitive maps are:
\begin{enumerate}[label=(\roman*)]
\item \term{system panorama}: compresses $S$ into $5$--$15$ major responsibility domains and
      answers: which responsibility domain is the most important? Why does each
      responsibility domain exist? What are the most important dependency relations among
      them? Which mechanisms are shared by several responsibility domains? Which states and
      behaviours decide the correctness of the whole system?
\item \term{key subsystem map}: for each key subsystem, answers eight questions:
      responsibility; boundary (what it is responsible for and what it is not); state (what
      state it owns and who may modify it); mechanism; interface; constraints; failure
      behaviour; and implementation entry points (which source files and tests can verify the
      above conclusions).
\item \term{key execution path map}: converts the static architecture into a model of dynamic
      behaviour. For example, the real questions behind a query path ``input $\to$ parsing
      $\to$ planning $\to$ execution $\to$ storage access $\to$ return'': when is the query
      plan created? Where is state kept during execution? Who manages the transaction
      snapshot? How is an error propagated upwards when a storage access fails? Which steps
      hold locks? Which operations may trigger disk I/O? Which resources must be released
      when a query is cancelled?
\item \term{contract map}: records the relations among components by each of the five edge
      classes of Definition \ref{def:relgraph}, in particular the \emph{oblivious set} of each
      layer, that is, what that layer is allowed never to care about.
\end{enumerate}
\end{definition}

Map (iv) is this paper's addition to conventional architecture documentation. Its necessity
follows directly from Proposition \ref{prop:oblivious}: a contract that has not declared its
oblivious set has not declared its compression boundary, and hence one cannot judge how far
the level above may rely on it.

\section{Operational Criteria for Comprehension}\label{sec:five-capabilities}

``Comprehension'' must be operationalised, or it cannot be decided whether comprehension has
been achieved. This paper adopts a five-capability criterion.

\begin{definition}[Criteria for comprehension]\label{def:understand}
An agent (human or machine) is said to \term{understand} a subsystem $M$ to the degree
required by a task $T$ if and only if it can do the following:
\begin{enumerate}[label=(U\arabic*)]
\item \term{Explain}: state why this subsystem exists;
\item \term{Predict}: predict what behaviour a given input, state change or error will cause;
\item \term{Locate}: decide which components to inspect when a fault occurs;
\item \term{Modify}: find the correct extension points and the constraints that must be
      respected when a feature is to be added;
\item \term{Verify}: know which tests or implementation details can substantiate its own
      judgement.
\end{enumerate}
This paper calls a state that attains (U1)--(U5) \term{operational abstract comprehension}.
\end{definition}

The relation between Definition \ref{def:understand} and Definition \ref{def:sufficient} is
direct: (U2) is the behavioural statement of $T$-sufficiency; (U1), (U3) and (U4) are concrete
query families in $T$; (U5) is the requirement that one can trace back to lower-level evidence,
namely the two clauses ``sufficiently accurate'' and ``verifiable'' of Definition
\ref{def:mdl}.

One diagnostic corollary matters here: if one can only recite the architecture diagram but
cannot predict what a change will affect, the current abstraction is not yet sufficient. This
and the seam detectors of \S\ref{sec:seams} are one thing seen from two directions.

\section{Reading Order}\label{sec:protocol}

\begin{principle}[Descent protocol]\label{pr:protocol}
Form the model top-down, descend to verify on demand, and never read sequentially.
\end{principle}

The steps are as follows.

\paragraph{The Discipline Layer.}
Read the README, architecture documents, architecture decision records, design papers and the
kernel documentation tree. About $1\%$ of the time establishes about $80\%$ of the skeleton
model. The output is a hand-drawn organ diagram: which large blocks the system has, what each
is responsible for, and how data flows.

\paragraph{Runtime Observation.}
Run the system first, read it afterwards. Runtime behaviour exposes structure and is far more
honest than static code: observe the process and thread model; set breakpoints or trace the
whole journey of one typical request; read logs and metrics, \emph{the vocabulary of the logs
is the system's self-description}. One end-to-end trace is worth ten thousand lines of static
reading.

\paragraph{Interface Boundaries.}
Read the header files, public interfaces and exported symbols of modules in dependency order,
not their implementations. An interface is a compression boundary (Proposition
\ref{prop:oblivious}) and states what this layer promises; how it discharges that promise is
not yet needed. What is wanted at this point is an organ atlas, not a tissue section.

\paragraph{Data Shape.}
The data model is the most stable abstraction and control flow the most volatile. Establish
first: what the core entities are, who owns their lifecycle, which state machines exist, and
what is writable and what is read-only. Once the shape of the data is in hand, control flow is
only data rearranged over time.

\paragraph{Descent on Demand.}
Descend into code only with a concrete task in hand: fix this defect, add this feature.
Reading is then directed retrieval rather than wandering. Descent follows the V-shape ``call
site $\to$ implementation $\to$ underlying mechanism'': stop as soon as the current problem
can be explained, then return to the upper layer and reassess the model.

\paragraph{Prediction and Verification.}
At each layer, write the mental model down first (``I guess the executor fetches the next
tuple like this''), then verify it. \emph{Being able to predict behaviour is comprehension}.
If the prediction is right, the compression model is effective; if it is wrong, that spot is
exactly where the system is most counter-intuitive.

\paragraph{Write-Back to the Contract Map.}
Write the newly acquired knowledge back into the contract map of Definition \ref{def:fourmaps},
filling in that layer's new guarantees, new assumptions and new oblivious set. This step is
the only means of preventing cognitive drift: without write-back, the layered model gradually
drifts away from the real system.

\begin{remark}[An operational ratio]
Understand about $10\%$ of the code in depth while knowing where the other $90\%$ is and who
owns it. The ratio is not a precise target but a quantified reminder of the principle that
comprehension is not exhaustion.
\end{remark}

\section{Detection Criteria for Module Boundaries}\label{sec:seams}

Architecture is ``the places where a change is localised''. The key operation of comprehension
is therefore \emph{finding seams}. The criteria below are ordered by reliability, and all can
be executed mechanically or semi-mechanically.

\begin{enumerate}[label=(S\arabic*)]
\item \term{Where types stop crossing}. Where is data re-encoded, wrapped or degraded to a
      generic type (byte arrays, strings, maps)? This is a strong signal of a module boundary:
      where the type system stops providing guarantees is where a human must supply the
      contract.
\item \term{Units that can be tested in isolation}. Where does a test suite exist that can be
      run without starting the whole system? This is the sharpest criterion: \emph{only a unit
      that can be verified in isolation is a real module}. Its reliability comes from
      Definition \ref{def:sufficient}: a unit test is essentially the assertion that ``the
      interface observation is a sufficient statistic of behaviour''.
\item \term{Serialisation, version fields and copy boundaries}. Where are there version
      fields, serialisation, deep copies, cross-process or cross-network calls? Crossing such
      a boundary makes a contract visible and falsifiable.
\item \term{Centralised enforcement sites of invariants}. Where are there assertions, database
      constraints, type-level constraints or contract checks? Where invariants are enforced
      centrally is usually where a level boundary lies.
\item \term{Independent targets of the build system}. Where does the build configuration
      define an independent compilation target or package? This is a \emph{necessary} but not a
      \emph{sufficient} condition for a boundary: an independent target is not necessarily a
      real module, but a real module nearly always has one.
\item \term{Co-change graphs and incident history}. This is the most underrated technique.
      Examine the commit history: which files tend to be modified together in the same change?
      \emph{A co-change graph gives semantic coupling, whereas an import graph gives syntactic
      coupling}. Where the two disagree is exactly where the architecture documentation
      disagrees with reality. Incident and defect histories then give the real invariant gaps.
\end{enumerate}

\begin{principle}[Summary of the seam criteria]\label{pr:seam-summary}
A module is real if and only if its interface observation is a sufficient statistic of its
external behaviour; a boundary is real if and only if, once it has been crossed, changes on
the other side no longer need to be understood.
\end{principle}

Chapter \ref{ch:design} inverts Principle \ref{pr:seam-summary} and uses it as a design
objective.

\section{Active Identification and Passive Reading}\label{sec:active}

\begin{principle}[Comprehension as system identification]\label{pr:sysid}
Comprehending a complex system is essentially \emph{system identification}: constructing a
model that predicts its behaviour and quantifying the uncertainty of that model. The most
efficient form of system identification is the ``perturb--observe'' loop.
\end{principle}

Passive reading is an extremely poor way to sample, because it \emph{has no error signal}.
Every judgement it produces is unfalsifiable, and therefore cannot be updated. Active probing
is the opposite: form a falsifiable prediction first (``if I change this, only these two
modules should be affected''), make a minimal change or run a test, then observe the
prediction error. \emph{The error is the part of the structure you cannot see.}

\begin{proposition}[The advantage of active probing]\label{prop:active}
Let comprehension be an iterative update of a model $M$, in which each observation $o$
supplies information $I(\text{structure}; o \mid M)$ about the true structure. The observation
distribution of passive reading does not depend on $M$, so its expected information gain is
constant in $M$; the observation distribution of active probing depends on the current
prediction $\hat{y}(M)$, and its expected information gain is proportional to the prediction
error $\mathbb{E}\,|y - \hat{y}(M)|$. Hence, wherever the prediction error is non-zero, active
probing has strictly higher marginal information gain.
\end{proposition}

\begin{proof}[Argument]
The samples drawn by passive reading are independent of $M$, hence $p(o \mid M) = p(o)$ and
the information gain is constant. Active probing selects the observations that maximise
predictive disagreement (D-optimal or uncertainty sampling in experimental design), so
$p(o \mid M)$ concentrates on the region where the current predictions of $M$ are uncertain;
by the definition of information gain $I = \mathbb{E}[\log p(o\mid M) - \log p(o)]$, that
quantity is non-negative and vanishes when $p(o\mid M) = p(o)$, that is, when prediction and
truth agree. Hence active probing is strictly better wherever the prediction error is non-zero.
\end{proof}

For AI the conclusion carries more weight: the advantage of AI is not that it can ``read'' more
code but that it can \emph{probe} at high throughput, generating hypotheses, running tests,
making minimal changes, inspecting co-changes, reading incident records, and then updating the
model. That is a frequency humans cannot match.

\section{The Observability Criterion}\label{sec:observability}

\S\ref{sec:comprehension-formal} defines comprehension as task-constrained compression, and
Proposition \ref{prop:stopping} gives the stopping point of descent. Both presuppose that
continued observation can improve the model. This section identifies a stronger limitation:
\emph{when observation is insufficient, increasing the number of observations changes
nothing}.

\begin{definition}[Observation operator and task subspace]\label{def:observability}
Let $\Beh$ be the behaviour space of the system and $\mathcal{T} \subseteq \Beh$ the
\emph{task-relevant difference subspace}, that is, the subspace spanned by the differences
that the query family $Q$ of Definition \ref{def:task} can distinguish. Let
$\mathcal{J} : \Beh \to \mathcal{Y}$ be the observation operator \emph{available} to the
agent: it maps behaviour to what the agent can actually read (source code, logs, traces, test
results, runtime metrics). Define the \term{observability rank} on $\mathcal{T}$ as
\begin{equation}\label{eq:obsrank}
\mathrm{rank}_{\mathcal{T}}\mathcal{J} \;=\;
\dim \mathcal{J}\bigl(\mathcal{T}\bigr),
\end{equation}
the dimension of the task-relevant differences that remain distinguishable after
$\mathcal{J}$. The construction is isomorphic to the observability criterion of linear system
theory \cite{kalman1960}.
\end{definition}

\begin{proposition}[The observability criterion]\label{prop:observability}
An agent can identify the system to $T$-sufficiency (Definition \ref{def:sufficient}) if and
only if
\begin{equation}\label{eq:obscond}
\mathrm{rank}_{\mathcal{T}}\mathcal{J} \;=\; \dim \mathcal{T}.
\end{equation}
When \eqref{eq:obscond} fails, the kernel of $\mathcal{J}$ on $\mathcal{T}$ is non-trivial:
there exist two executions $e_1, e_2$ whose available observations are all identical while
their task-relevant behaviour differs. In that case \emph{no number of additional observations
can reduce task loss}, because new observations still lie in the range of the same
$\mathcal{J}$.
\end{proposition}

\begin{proof}
If \eqref{eq:obscond} holds, then $\mathcal{J}$ is injective on $\mathcal{T}$, task-relevant
differences are uniquely recoverable from observations, and $T$-sufficiency follows from
Definition \ref{def:sufficient}. If it fails, then
$\Ker \mathcal{J} \cap \mathcal{T} \ne \{0\}$; take a pair of executions $e_1 \ne e_2$
corresponding to a non-zero element of that kernel (in the task-relevant sense). The two are
indistinguishable under $\mathcal{J}$, so every $\mathcal{J}$-measurable decision rule gives
the same answer and is wrong on at least one of them. Adding observations does not change
$\mathcal{J}$, and therefore does not change that kernel.
\end{proof}

Proposition \ref{prop:observability} separates ``insufficient comprehension'' into two cases
that are entirely different in engineering terms:

\begin{principle}[Distinguishing two kinds of comprehension failure]\label{pr:two-failures}
\begin{enumerate}[label=(\alph*)]
\item \term{rank deficiency}: the observation operator itself cannot distinguish task-relevant
      differences. The only remedy is to \emph{change the observation}: add invariant checks,
      add trace points, add assertions, expose internal state. This is a design action, not a
      reading action.
\item \term{full rank but insufficient samples}: observation suffices to distinguish, but has
      not yet been collected. The remedy is to \emph{probe more}, with its efficiency governed
      by Proposition \ref{prop:active}.
\end{enumerate}
Confusing the two is the main reason advice such as ``read more code'' fails: when the problem
belongs to (a), additional reading is information-theoretically ineffective.
\end{principle}

\begin{remark}[Measuring the cost of comprehension correctly]\label{rmk:sample-complexity}
The cost of case (b) is a \emph{sample complexity}, not an amount of reading. If the contract
hypotheses are confined to a hypothesis class $\mathcal{H}$ (the set of operations in
Definition \ref{def:operad} provides a natural candidate), then the number of probes required
to identify the system to $T$-sufficiency at error $\varepsilon$ has a lower bound decided by
the distinguishing power of $\mathcal{H}$ on $T$, typically of the form
$\Omega(d/\varepsilon)$, where $d$ is the capacity dimension of the hypothesis class.

This accounting yields a testable prediction: \emph{the factor by which designed probing beats
exhaustive reading is quantifiable}, and that factor grows with the capacity of the hypothesis
class. Hypothesis $H_5$ of Chapter \ref{ch:program} thereby acquires a quantitative form.
\end{remark}

\begin{example}[A typical shape of rank deficiency]\label{ex:rank-deficient}
A service keeps all its state in process memory and exposes only the return values of requests.
If the task is to decide ``whether two concurrent requests can read each other's partial
updates'', then no number of return-value observations suffices, because the difference lies
entirely outside the return values. The correct action is not to read more code but to
introduce an observable: a runtime assertion of a concurrency invariant, or a test with
interleaving injection. The example shows that the distinction of Principle
\ref{pr:two-failures} is practical rather than merely taxonomic.
\end{example}

\section{Budget and Stopping Conditions}\label{sec:budget}

\S\ref{sec:comprehension-formal} has given the formal stopping criterion (Proposition
\ref{prop:stopping}). This section gives its engineering version and the allocation of budget.

\begin{principle}[Allocation of the descent budget]\label{pr:budget}
Allocate the comprehension budget in order of information gain per unit cost, not by code
volume or directory order. The typical high-yield entry points, in order, are: external
interface definitions and the data model $>$ one trace of an end-to-end execution path $>$
invariant enforcement sites $>$ co-change hotspots $>$ specific implementation details.
Low-yield entry points are: sequential reading of implementation bodies, reading through a
utility function library, and reading through generated code.
\end{principle}

\begin{principle}[Stopping conditions]\label{pr:stop}
Stop descending as soon as any of the following holds:
\begin{enumerate}[label=(\alph*)]
\item you can name every invariant governing the current change;
\item further descent no longer changes your prediction of ``what this change will affect'';
\item the cost of obtaining the next piece of evidence exceeds the expected rework cost of the
      current change.
\end{enumerate}
\end{principle}

Condition (a) corresponds to Principle \ref{pr:descent}; condition (b) to Proposition
\ref{prop:stopping}; condition (c) is an economic judgement: comprehension is not a matter of
more being better, but of \emph{enough}.

\section{Implications for Machine Comprehension}\label{sec:ai-comprehension}

The difficulty facing AI is not only a limited context length. Even when a model can read a
large amount of code, three distinct problems remain: too much information drowns the
important information; local pieces of information depend on one another in complex ways; and
comprehension of details does not automatically form a correct system-level model. Enlarging
the context window alone therefore does not solve every problem.

\begin{principle}[A software cognitive architecture for AI]\label{pr:ai-cog}
An effective machine cognitive architecture should have four layers:
\begin{enumerate}[label=(\roman*)]
\item \term{global cognitive model}: system goals, architecture diagram, key invariants, main
      data flows, core design decisions;
\item \term{subsystem cognitive model}: component responsibilities, interfaces, state
      ownership, dependency relations, lifecycle and error semantics;
\item \term{implementation-level index}: key functions, type definitions, call chains, source
      locations, tests and runtime evidence;
\item \term{verification and feedback}: compilation, tests, static analysis, execution traces,
      profiling, used to correct the models above.
\end{enumerate}
Its key property is: \emph{an upper layer model need not replicate all lower-level information,
but need only retain enough information for the agent to find the correct lower-level evidence
when it is required.}
\end{principle}

\begin{definition}[Traceable subsystem cognitive record]\label{def:record}
The cognitive record of each subsystem should contain at least nine fields, listed in Table
\ref{tab:record}.
\end{definition}

\begin{table}[htbp]
\centering
\small
\caption{The minimum field set of a subsystem cognitive record}\label{tab:record}
\begin{tabular}{L{3.0cm}L{11.4cm}}
\toprule
\textbf{Field} & \textbf{Role} \\
\midrule
Purpose & Why this subsystem exists \\
Contract & What guarantees it provides upwards (that is, $\Obs_i(\Beh_{i+1})$) \\
Oblivious & What this layer is allowed never to care about (that is, $\Obl(\Lambda_i)$) \\
State & What state it owns and who may modify it \\
Dependencies & Which other components it depends on \\
Invariants & Which conditions must never be violated, and where they are enforced \\
Lifecycle & How key objects are created, run and destroyed \\
Execution paths & How important requests or tasks pass through the system \\
Evidence & Which source files, tests or experiments support the above conclusions \\
Uncertainty & What has not been verified, and which assumptions may be wrong \\
\bottomrule
\end{tabular}
\end{table}

The last two entries of Table \ref{tab:record}, together with the Oblivious entry, are
additions of this paper, for the following reasons.

\emph{On Uncertainty}: AI cannot treat an architecture summary it has generated itself as
fact. It needs to know which conclusions come from actual source code, which come from tests,
and which are only provisional inferences. This is the same requirement as the epistemic
discipline of Chapter \ref{ch:objections}.

\emph{On Oblivious}: as stated above, it is the compression boundary itself. Without this
entry, a cognitive record describes only ``what there is'' and not ``what may be
disregarded'', and so cannot support the modification capability (U4).

\emph{On dynamic update}: when the code changes, the cognitive record must be updated in step.
A usable machine comprehension system should therefore not be a static knowledge base but a
\emph{bidirectionally traceable} relation between the source code and the abstract model.

\begin{principle}[Readability turns from a virtue into infrastructure]\label{pr:readability}
The ``comprehension interface'' of a system is becoming a first-class engineering deliverable.
Documentation used to be an appendix written for human readers; now machine-readable
architecture descriptions, the practice of interface-as-documentation, and agent-facing
working conventions are becoming components that a system must design, because the readership
of code now includes a non-human intelligence that relies only on the compression layer.
\end{principle}

The empirical content of this chapter is Hypothesis $H_5$ (Chapter \ref{ch:program}): a
comprehension process that uses a layered model together with active probing achieves a higher
task success rate per unit cost than exhaustive reading. Propositions \ref{prop:stopping},
\ref{prop:active} and \ref{prop:observability} hold under the stated model; Definition
\ref{def:understand} is a working definition whose value depends on whether it can be used to
distinguish ``comprehension'' from ``recitation'', which requires experimental test.

% =====================================================================
\chapter{Design: forward compression}\label{ch:design}
% =====================================================================

If comprehension is inverse compression, design is forward compression: compression
boundaries are built in advance for future readers. This chapter turns that symmetry into
an optimisable objective function, states its dynamical constraint, and shows where
layering fails.

\section{Read-sets and the objective function}\label{sec:readset}

The measurable objective of architecture is not an adjective such as ``decoupling'' or
``high cohesion'' but a number that can be counted.

\begin{definition}[Change read-set]\label{def:readset}
Let the code base consist of a set of modules $\Mod$, namely
$\bigcup_{c \in \mathcal{C}} \rho(c)$, where $\mathcal{C}$ is the set of changes. For a
change $c$, its \term{read-set} $\rho(c) \subseteq \Mod$ is the set of modules that must be
read in order to complete $c$ correctly. The \term{architectural quality metric} of the
system is
\begin{equation}\label{eq:readset}
Q(\Mod) \;=\; \mathbb{E}_{c \sim \mathcal{C}}\bigl[\,|\rho(c)|\,\bigr],
\end{equation}
that is, the mean read-set size weighted by the historical change distribution.
\end{definition}

Definition \ref{def:readset} is operational: it can be estimated directly from version
history, by counting for each merged change the number of top-level modules, files and
directories it touches and taking the resulting distribution. This is not a perfect proxy
(touching is not having to read), but it is \emph{measurable}, and ``degree of
modularity'' is not.

\begin{definition}[Enforcement sites of an invariant]\label{def:enforce}
Let $\Inv$ be the set of invariants of the system. For $\iota \in \Inv$, its
\term{enforcement-site set} $\mathrm{S}(\iota) \subseteq \Mod$ is the set of all modules
that \emph{must be implemented correctly for $\iota$ to be guaranteed}. Enforcement is
\term{effective} if $\iota$ is not violated in any execution. Its \term{unattested part}
$\mathrm{S}_{\mathrm{un}}(\iota) \subseteq \mathrm{S}(\iota)$ is the subset of those sites whose
correctness cannot be established without reading them: sites with no machine-checked proof, no
enforced type-level constraint, no separately checkable contract and no accepted certificate. Sites
outside $\mathrm{S}_{\mathrm{un}}(\iota)$ are still \emph{necessary}---if they were wrong the
invariant would fail---but they are \emph{substitutable by their evidence}, in the sense that an
agent may rely on the guarantee without inspecting the implementation.
\end{definition}

\begin{theorem}[Centralisation of invariants]\label{th:centralise}
Suppose that enforcement of invariant $\iota$ is effective and that a change $c$ may break
$\iota$, that is, there exists a legal sequence of changes that breaks $\iota$ in the
absence of an understanding of $\mathrm{S}(\iota)$. Then
\begin{equation}\label{eq:central}
|\rho(c)| \;\ge\; \bigl|\mathrm{S}_{\mathrm{un}}(\iota)\bigr| .
\end{equation}
Hence, other things equal, minimising $\mathbb{E}[|\rho|]$ requires minimising
$\sum_{\iota \in \Inv} |\mathrm{S}_{\mathrm{un}}(\iota)|$, whose lower bound is attained when every
enforcement site is either the only one or independently checkable.
\end{theorem}

\begin{proof}
If $c$ may break $\iota$, an agent executing $c$ must be able to judge whether the change is safe.
For each site in $\mathrm{S}_{\mathrm{un}}(\iota)$ that judgement requires reading that site, since
by definition no other evidence settles it; hence
$\rho(c) \supseteq \mathrm{S}_{\mathrm{un}}(\iota)$, which gives \eqref{eq:central}. Sites in
$\mathrm{S}(\iota) \setminus \mathrm{S}_{\mathrm{un}}(\iota)$ may be relied upon through their
evidence and need not be read. Summing over $\iota$:
$\mathbb{E}[|\rho|] \ge \mathbb{E}[|\bigcup_\iota \mathrm{S}_{\mathrm{un}}(\iota)|]$, with equality
in the optimistic case in which those sets are disjoint. The bound is attained at
$|\mathrm{S}_{\mathrm{un}}(\iota)| = 0$, that is, when every enforcement site carries evidence that
can be checked without being read---a type system, a private constructor, a capability token, a
database constraint or a proof.
\end{proof}

The change from $|\mathrm{S}(\iota)|$ to $|\mathrm{S}_{\mathrm{un}}(\iota)|$ is not a weakening
introduced for safety. It is the correction of a genuine conflation, and the distinction it draws
was missing from the argument. Three different sets were being treated as one: the modules whose
correct behaviour is necessary for $\iota$ (that is $\mathrm{S}(\iota)$); the evidence needed to
show a given change safe; and the source files an agent must actually read. The first does not bound
the third except through the second, and the second is exactly what $\mathrm{S}_{\mathrm{un}}$
isolates. A module that has a machine-checked proof or a stable, independently checkable interface
does not enter the read-set, no matter how necessary its correctness is.

Two consequences follow, and they are the practically useful ones. First, the bound explains why
encoding invariants in types has engineering value beyond expressiveness, and does so quantitatively:
it lowers $|\mathrm{S}_{\mathrm{un}}(\iota)|$ rather than merely relocating the invariant. Second, it
tells a maintainer where the read-set actually comes from: not from the number of places an invariant
matters, but from the number of those places whose correctness rests on nothing checkable. That is a
smaller and more actionable number, and it is the one a refactoring should target.

\begin{corollary}[Optimal location of enforcement sites]\label{cor:site}
Among all candidates satisfying the effectiveness constraint, the enforcement sites that minimise
the expected read-set of \eqref{eq:central} are the modules of \emph{least fan-in}, that is, those on
which fewest other modules depend and which are therefore least likely to be circumvented by
modification. The scope of this optimality is the read-set objective alone; see the remark below.
\end{corollary}

\begin{remark}[Limits of Theorem \ref{th:centralise}]\label{rmk:central-limit}
The theorem has one premise: enforcement must be \emph{effective}. In practice,
concentrating an invariant in a single place often fails, because callers can circumvent
that place. This is the real reason why ``scattered invariants'' recur in engineering. It
is not ignorance but a defence against the risk of circumvention. The correct reading of
Theorem \ref{th:centralise} is therefore: \emph{scattering invariants is a defence with a
cost; where invariants can genuinely be centralised (for instance by a type system, private
constructors, capability tokens or database constraints that make circumvention impossible),
centralisation is strictly better than scattering.}
This explains why ``encoding invariants in types'' has engineering value beyond
expressiveness: it lowers $|\mathrm{S}_{\mathrm{un}}(\iota)|$ towards $0$ without increasing the risk
of circumvention. The two mechanisms are independent and both matter---types remove the need to read,
effective enforcement removes the need to trust.
\end{remark}

\begin{remark}[Centralisation optimises one objective among several]\label{rmk:central-multi}
Corollary \ref{cor:site} identifies the read-set-optimal enforcement site. It should not be read as
identifying the best site, because the read-set is one objective among several, and the others are
not represented in the model at all. Moving an invariant into a single module also moves \emph{load}
onto that module, adds a hop to every operation that would otherwise be local, concentrates
\emph{failure impact}---if the enforcement site fails, every dependent fails together---and creates a
\emph{change bottleneck}, since every change touching the invariant now serialises through one
owner.

These are the same dimensions the paper insists must not be collapsed elsewhere
(Principle \ref{pr:crosscutting}), and the honest statement of the corollary is therefore
conditional: \emph{given} that effectiveness is achievable at the site in question and \emph{given}
that the other four dimensions are acceptable there, least fan-in is optimal for the read-set. When
they are not acceptable, the correct problem is a multi-objective one---correctness, change impact,
latency, and failure domain---and the read-set bound is a constraint within it rather than the
objective. A design review that treats ``centralise the invariant'' as an unconditional
recommendation is applying the corollary outside its scope.
\end{remark}

\section{Near-decomposability}\label{sec:neardecomp}

Simon's near-decomposability \cite{simon1962} is the classical explanation of why a level
structure is stable: interaction within a layer is dense, interaction between layers is
sparse. But ``sparse interaction'' is not itself testable. This section gives a testable
equivalent characterisation.

\begin{definition}[Near-decomposable]\label{def:neardecomp}
Let module $M$ have an interface $I$. For an execution $e$ of $M$, let
$\mathrm{trace}_I(e)$ be the sequence of interactions of $e$ on $I$ and
$\mathrm{behav}_{\mathrm{out}}(e)$ the behaviour of $e$ observable outside $M$. Then $M$ is
\term{near-decomposable} with respect to $I$ if there exists a function $f$ such that
\begin{equation}\label{eq:neardecomp}
\forall e:\quad \mathrm{behav}_{\mathrm{out}}(e) \;=\;
f\bigl(\mathrm{trace}_I(e)\bigr),
\end{equation}
that is, the observation on $I$ is a \emph{sufficient statistic} for the external behaviour
of $M$.
\end{definition}

The relation between Definition \ref{def:neardecomp} and Definition \ref{def:sufficient}
($T$-sufficiency) is immediate: take the query family $Q$ of a task $T$ to be
``observations of external behaviour''; near-decomposability is then exactly the
$T$-sufficiency of $M$ with respect to its interface. Simon's qualitative principle is
thereby translated into a decidable condition.

\begin{proposition}[Leak pairs]\label{prop:leakpair}
$M$ is \emph{not} near-decomposable with respect to $I$ if and only if there exist two
executions $e_1, e_2$ such that
\begin{equation}\label{eq:leakpair}
\mathrm{trace}_I(e_1) = \mathrm{trace}_I(e_2)
\qquad\text{and}\qquad
\mathrm{behav}_{\mathrm{out}}(e_1) \ne \mathrm{behav}_{\mathrm{out}}(e_2).
\end{equation}
Such a pair is called a \term{leak pair}.
\end{proposition}

\begin{proof}
$(\Leftarrow)$ If such $e_1, e_2$ exist, then no function $f$ can satisfy both
$f(\mathrm{trace}_I(e_1)) = \mathrm{behav}_{\mathrm{out}}(e_1)$ and
$f(\mathrm{trace}_I(e_2)) = \mathrm{behav}_{\mathrm{out}}(e_2)$, since the arguments on the
two sides are equal while the assigned values are not; hence \eqref{eq:neardecomp} fails.
$(\Rightarrow)$ If \eqref{eq:neardecomp} fails, no $f$ satisfies the condition, and by the
definition of a function there must exist two points with the same argument and different
assigned values, which is \eqref{eq:leakpair}.
\end{proof}

\begin{corollary}[A diagnostic procedure]\label{cor:diagnose}
Proposition \ref{prop:leakpair} yields an \emph{executable} diagnostic: construct two
scenarios identical at the interface but differing in internal state or external
environment, and observe whether the external behaviour diverges. If it diverges, the
interface is not a sufficient statistic, and the module boundary leaks there.
\end{corollary}

Corollary \ref{cor:diagnose} turns ``does this module's abstraction leak'' from an aesthetic
judgement into a differential test. Its practical value is that \emph{whether an abstraction
leaks is measurable, not merely felt}. The engineering commonplaces of ``shadow state'',
``implicit globals'' and ``cross-request cache pollution'' are all concrete instances of
leak pairs.

\begin{principle}[An equivalent design objective]\label{pr:design-eq}
By Corollary \ref{cor:diagnose}, the testable objective of design can be written as:
\emph{for every claimed module boundary, the interface observation should be a sufficient
statistic for external behaviour; the existence of a leak pair is the failure of that
boundary.}
\end{principle}

\begin{principle}[Two criteria for a claimed boundary, not one]\label{pr:two-criteria}
A claimed boundary between a module and its environment makes two separate claims, and they fail
separately.
\begin{enumerate}[label=(\alph*)]
\item \term{Sufficiency}: the interface observation determines the external behaviour. This is
      Definition \ref{def:neardecomp}; its failure is a leak pair (Proposition \ref{prop:leakpair}),
      found by differential testing (Corollary \ref{cor:diagnose}).
\item \term{Completeness}: the observation map the boundary declares is the whole of what is
      observable across it. Its failure is a difference the boundary declared irrelevant becoming
      visible anyway---through timing, resource occupancy or failure behaviour---and it is found by
      the techniques of information-flow analysis rather than by differential testing
      (Remark \ref{rmk:noninterference}).
\end{enumerate}
The two are independent, in both directions. A boundary can be sufficient and incomplete: every pair
of interface-identical executions has the same external behaviour, and yet a declared-irrelevant
difference is readable from above through a timing channel. It can equally be complete and
insufficient: nothing leaks through undeclared channels, while the declared interface is too coarse
to determine the behaviour.

Corollary \ref{cor:diagnose} and Principle \ref{pr:design-eq} address the first criterion only.
The second is the one that fails silently, because a boundary that leaks through an undeclared
channel passes every differential test written against the declared interface---the test never
constructs the offending pair, since the pair differs in something the boundary said did not matter.
\end{principle}

\section{Technical debt}\label{sec:debt}

``Technical debt'' is usually a metaphor in engineering usage. This section gives a
definition that is precise and still useful.

\begin{definition}[Technical debt]\label{def:debt}
Given an encoding, let the description length of the system code be $\DL(\text{code})$ and
let its \term{shortest generating description} be
$\DL^{\star} = \min\{\,|d| : d \text{ generates code}\,\}$, that is, the shortest program
that generates that code in the Kolmogorov sense. Define the \term{technical debt}
\begin{equation}\label{eq:debt}
D \;=\; \DL(\text{code}) \;-\; \DL^{\star}.
\end{equation}
\end{definition}

The definition is defensible but non-computable ($\DL^{\star}$ is not computable), so its
value lies in the \emph{direction} and the \emph{criterion} it supplies, not in a numerical
value.

\begin{principle}[How technical debt accrues and is repaid]\label{pr:debt}
Technical debt grows when \emph{a change is not re-abstracted}: if a modification is not
re-abstracted, the gap between the code and its ``shortest generating description'' widens.
Therefore:
\begin{enumerate}[label=(\alph*)]
\item refactoring $=$ recompression, a survival necessity for holding the compression ratio
      above the comprehensibility threshold;
\item a system that is not refactored does not stop growing; it accumulates
      incomprehensibility;
\item deciding whether a refactoring is worthwhile is equivalent to deciding whether it
      significantly lowers $D$, that is, whether it makes the system easier to generate
      from a shorter description.
\end{enumerate}
\end{principle}

Principle \ref{pr:debt}(c) gives a practical criterion: if a refactoring does not let you
say in fewer words what the module does, it is probably only reformatting.

\begin{remark}[A rate--distortion reading of refactoring]\label{rmk:rd-refactor}
Regard $(\DL(M), L_T(M))$ as a point in the rate--distortion plane \cite{shannon1959}
(Definition \ref{def:mdl}). The current state of a system is a point on the curve or outside
it: \emph{outside the curve} means that a model of the same description length attains lower
loss, that is, that there is pure room for improvement; refactoring is then moving along the
curve or moving the point back onto it.

This reading yields a testable prediction: \emph{the slope of the rate--distortion curve,
$\partial L_T / \partial \DL$, should predict which refactorings pay off.} A refactoring at
a point where the slope is zero does not lower task loss and is therefore not worth doing;
where the slope is large, it is worth doing even at high cost. If that predictive power does
not hold in real projects, the choice of objective function in this chapter requires
revision (Chapter \ref{ch:program}, hypothesis $H_{12}$).
\end{remark}

\section{Governable complexity and comprehension capacity}\label{sec:governable}

Section \ref{sec:readset} defines architectural quality as a static quantity
$\mathbb{E}[\rho]$. Real systems, however, have a \emph{rate}: changes arrive at some
frequency, and an organisation's reading capacity is finite. This section supplies that
dynamics; it is the only place in the paper where time is a variable.

\begin{definition}[Governable complexity and comprehension capacity]\label{def:governable}
Let $\lambda(t)$ be the change arrival rate (changes per unit time) and let
$\mathbb{E}[\rho]$ be the expected read-set (Definition \ref{def:readset}). Define the
\term{governable complexity}
\begin{equation}\label{eq:governable}
C(t) \;=\; \mathbb{E}\bigl[\rho\bigr] \cdot \lambda(t),
\end{equation}
that is, the information throughput required to maintain the system, in units of
module-reads per unit time. Let the \term{comprehension capacity} of the organisation be
\begin{equation}\label{eq:capacity}
K \;=\; h \cdot b,
\end{equation}
where $h$ is the number of people who can be applied to comprehension and $b$ is the
comprehension bandwidth of one person (module-reads per unit time). The system is in a
\term{governable state} if $C(t) \le K$ for all relevant $t$.
\end{definition}

\begin{proposition}[The scaling exponent of required headcount, under the capacity model]\label{prop:exponent}
Let a system have $n$ features, each changing independently at rate $\lambda_0$, so that
$\lambda \propto n$. Let the read-set $\mathbb{E}[\rho]$ scale as $n^{\alpha}$ in $n$, with
$\alpha \ge 0$ a parameter of the model, $\alpha = 0$ corresponding to complete
localisation and $\alpha = 1$ to unconstrained coupling. Then the governable complexity
scales as
\begin{equation}\label{eq:exponent}
C \;\propto\; n^{1+\alpha}.
\end{equation}
From $C \le K$ and $K = hb$ with $b$ constant, the headcount needed to keep the system governable
scales as $h \propto n^{1+\alpha}$. Therefore, \emph{within this model}:
\begin{enumerate}[label=(\alph*)]
\item if $\alpha$ is pressed to $0$, headcount grows linearly with system size and the system can
      persist indefinitely;
\item if $\alpha = 1$, headcount must grow as the \emph{square} of size, and the system must
      become unstable at some size.
\end{enumerate}
\end{proposition}

\begin{proof}
By Definition \ref{def:governable}, $C = \mathbb{E}[\rho]\cdot\lambda$. Substituting
$\lambda = n\lambda_0$ and $\mathbb{E}[\rho] = \Theta(n^{\alpha})$ gives
$C = \Theta(n^{1+\alpha})$. From $C \le K = hb$ it follows that
$h \ge C/b = \Theta(n^{1+\alpha})$.
\end{proof}

Two things the proposition does not establish, and that the earlier title---\emph{architecture fixes
the scaling exponent}---asserted by implication. First, it does not show that the exponent is fixed
by the architecture; it \emph{assumes} that $\mathbb{E}[\rho]$ scales as a power of $n$ and treats
$\alpha$ as the parameter of that scaling. Whether $\alpha$ is architecture-determined, and whether
the read-set proxy used to estimate it is valid, are both empirical questions---the second is
Hypothesis $H_4$, and the first is what Hypothesis $H_{14}$ asks directly. Second, it does not show
that real organisations' headcount follows the exponent. It shows that headcount must follow it
\emph{if} the organisation is to remain governable under the stated capacity model, which is a
conditional statement about a model of comprehension capacity, not a forecast about hiring. The
empirical content is that systems with higher estimated $\alpha$ should require disproportionately
more coordination effort per unit of growth; that is testable, and Chapter \ref{ch:program} proposes
how.

\begin{corollary}[The actual role of architecture]\label{cor:arch-role}
Proposition \ref{prop:exponent} assigns architecture a precise duty: it does not lower the
cost of an individual change but \emph{lowers the exponent with which cost grows with size}.
This explains why architectural benefit is invisible in small systems and decisive for
survival in large ones, and why the charge of ``architectural over-engineering'' always
appears to hold at small scale.
\end{corollary}

\begin{corollary}[Derivation of two classical conclusions]\label{cor:brooks}
Equation \eqref{eq:exponent} yields two classical conclusions:
\begin{enumerate}[label=(\alph*)]
\item \term{Relation between organisation size and system size}: from
      $h \propto n^{1+\alpha}$, organisation size must grow with system size, and the
      growth exponent is fixed by the architecture. This is the \emph{dynamical} version of
      the correspondence between organisational structure and system structure
      \cite{conway1968}; that correspondence is a static statement (structure mirrors
      organisation), whereas the dynamical version gives the governability bound.
\item \term{Adding people makes a project slower}: increasing $h$ raises $K$ linearly, but
      at the same time raises the number of changes in flight and hence $\lambda$, and
      raises $\mathbb{E}[\rho]$ through the interfaces it adds. When the growth of the
      latter two exceeds that of the former, $C/K$ worsens rather than improves. This is an
      in-model explanation of the empirical rule that adding people to a late project makes
      it later.
\end{enumerate}
\end{corollary}

\begin{principle}[Design implication]\label{pr:governable-design}
By Proposition \ref{prop:exponent}, the criterion of an architecture review should not be
``is the read-set small'' but \emph{``what is the scaling exponent of the read-set in
size''}. A system whose read-set is currently small but grows rapidly with size is more
dangerous than one whose read-set is currently larger but scales flatly. This requires
hypothesis $H_4$ to be strengthened: the value of a read-set lies mainly in its
\emph{scaling}, not in its \emph{magnitude}.
\end{principle}

\begin{remark}[The only dynamics]\label{rmk:only-dynamics}
The paper has so far been static: structure, invariants, read-sets and contracts are all
time-invariant quantities. Definition \ref{def:governable} is the only dynamical variable,
because it is the only place where ``what an organisation can process'' and ``what the
system demands'' must match in \emph{rate}. Elsewhere time can be eliminated.
\end{remark}

\section{The design procedure}\label{sec:seven-steps}

\begin{enumerate}[label=(D\arabic*)]
\item \term{Fix the invariants first}. Before writing any code, write down the propositions
      that must hold forever of this system (for example ``the total balance is
      conserved'', ``an I/O thread never blocks'', ``the version number increases
      monotonically''). Invariants are the gravitation of a level; all structure
      crystallises around them. By Theorem \ref{th:centralise}, each invariant should
      designate a single enforcement site and state why circumvention is impossible.
\item \term{Choose the primitives}. Atoms, components, messages, transactions, entities,
      and so on. Once the primitives are chosen they become the shared language of all
      subsequent discussion of the system. \emph{Choosing the primitives is the most
      important single decision at this level}, because it fixes the shape of every
      expression thereafter.
\item \term{Fix the data model}. Entities, relations, lifecycles, ownership. A data model
      is an order of magnitude more stable than logic, so it should be frozen first.
\item \term{Cut the organs}. The basis of modularisation is not ``similar
      functionality'' but \emph{information that changes together being hidden together}
      \cite{parnas1972}. The test for each module is that its interface is small enough to
      hold in mind at once and its implementation deep enough to be replaced wholesale
      \cite{ousterhout2018}. The interface is the compression summary you deliberately
      maintain; the implementation is the compressed content that may be rewritten.
\item \term{Make the physical laws executable}. An invariant written in a document rots;
      only encoded in a type system, assertions, contract tests and property tests does it
      endure. \emph{Types are architectural documentation enforced by the compiler}. This
      is the confluence of ``designing for comprehensibility'' and ``designing for
      correctness'': an invariant that a machine can check need not be maintained from
      memory by a human or an AI.
\item \term{Record decisions}. Write down why each choice was made and what was given up.
      Code is the result; a decision record is part of the generating rules, and it spares
      the reader performing inverse compression the most expensive guessing.
\item \term{Refactoring $=$ recompression}. See Principle \ref{pr:debt}.
\end{enumerate}

One meta-rule runs through all seven:

\begin{principle}[Design for the reader's mental bandwidth]\label{pr:reader}
Design for the mental bandwidth of your reader, not for your typing speed. Code is read far
more often than it is written, and the users of an architecture are future comprehenders.
\end{principle}

\section{Cross-cutting concerns and the limits of layering}\label{sec:crosscutting}

The failure modes of layering must be acknowledged explicitly. One class of concerns cannot
be localised to any layer:

\begin{itemize}
\item security and authorisation (every module may introduce a vulnerability);
\item multi-tenant isolation (spanning storage, computation and network);
\item distributed tracing (which must run through every call path);
\item transactions and consistency (across the boundaries of state owners);
\item observability and quotas (required at every point of resource consumption).
\end{itemize}

\begin{principle}[Cross-cutting concerns are independent dimensions]\label{pr:crosscutting}
For cross-cutting concerns, the mature treatment is to acknowledge them as
\emph{independent dimensions} (policy plane, capability model, taint propagation, service
mesh) rather than forcing them into the hierarchy and pretending that they are a layer.
\end{principle}

This principle follows directly from the discussion in Chapter \ref{ch:convergence}: a
level structure is supported by near-decomposability, and near-decomposability presupposes
sparse interaction between layers. When a concern has a non-trivial interaction with every
module, it breaks near-decomposability by definition, and therefore cannot be represented as
``a layer''. The result of forcing the representation is that the level structure fails, and
the failure is usually misdiagnosed as ``insufficiently thorough layering''.

\section{Design checklist}\label{sec:checklist}

The following checklist turns this chapter into executable items. Each item should be
answered explicitly at a design review.

{\small\begin{longtable}{L{1.0cm}L{14.0cm}}
\caption{Design checklist for large systems}\label{tab:checklist}\\
\toprule
\textbf{No.} & \textbf{Check} \\
\midrule
\endfirsthead
\toprule
\textbf{No.} & \textbf{Check} \\
\midrule
\endhead
\bottomrule
\endlastfoot
A1 & Which propositions must hold forever of the system? Where is each enforced? Why is circumvention impossible?\\
A2 & Is the enforcement site of each invariant unique? If not, has the cost (the read-set) been explicitly accepted?\\
A3 & What are the primitives of this level? Are they the shared language of all subsequent discussion?\\
A4 & Has the data model been frozen before the logic? Are entities, relations, lifecycles and ownership explicit?\\
\midrule
B1 & Is each module's interface small enough to hold in mind at once?\\
B2 & Is each module deep enough to be replaced wholesale?\\
B3 & For each claimed boundary, does a leak pair exist? (That is: two executions with the same interface trace and different external behaviour.)\\
B4 & Has each module's \emph{oblivious set} been written down explicitly?\\
B5 & Which modules change together? Has this been verified against co-change history rather than intuition alone?\\
\midrule
C1 & How many unrelated modules must be modified to add a new feature? (The measured read-set.)\\
C2 & Does replacing the underlying implementation force large-scale modification of the business logic?\\
C3 & When a subsystem fails, where does the failure propagate?\\
C4 & Does adding a task type require rewriting the shared scheduling logic?\\
C5 & Does concurrent access break state invariants? Where is this detected?\\
C6 & Are performance bottlenecks concentrated in a few critical paths?\\
\midrule
D1 & Have the invariants been encoded in executable form (types, assertions, contract tests, property tests)?\\
D2 & Can each module's contract be verified in isolation? If not, is it a module?\\
D3 & Does each critical subsystem have a traceable record of what is known about it (Table \ref{tab:record})?\\
D4 & Does the decision record cover why each choice was made and what was given up?\\
D5 & Are cross-cutting concerns treated as independent dimensions rather than disguised as a layer?\\
\midrule
E1 & Does the system need a unified substrate? Judge by $n \times m$, not by lines of code.\\
E2 & Of the six functions, which are built and which are adopted? The criterion is the domain-criticality of the invariants.\\
E3 & For the parts that are built, are their invariants ones that a general-purpose host cannot provide?\\
E4 & For the parts that are adopted, is the strength of the guarantee offered by the host sufficient to support this layer's contract?\\
\bottomrule
\end{longtable}}

Theorem \ref{th:centralise} and Propositions \ref{prop:leakpair} and \ref{prop:exponent}
hold under the stated models. The empirical content is hypotheses $H_4$ and $H_{12}$: that
the scaling exponent of the read-set outperforms read-set magnitude and conventional size
metrics; that centralising invariants (where circumvention is impossible) lowers both the
read-set and the defect rate; and that the rate--distortion slope predicts the benefit of a
refactoring. If all three are refuted, the choice of objective function in this chapter
requires revision.

% =====================================================================
\chapter[From Intent to Verifiable Execution]{From Intent to Verifiable Execution}\label{ch:intent}
% =====================================================================

This chapter addresses Problem \ref{prob:trust}: if the cost of coding tends to zero and no
human reviews the implementation line by line, how is correctness guaranteed
mechanically. The answer has three parts: decompose correctness into dimensions that can be
treated separately, define the intermediate process as a verifiable refinement chain, and
organise trust as a composite of two modes of generation. The chapter also identifies an
entry point on that chain which cannot be mechanically closed, and the sense in which
consensus can approximate it.

\section{Dimensions of Correctness}\label{sec:five-correctness}

An AI rewrites a 1\,ns call into a 1\,$\mu$s call: functionally identical, but a thousand
times slower. The example shows that ``correct'' is not a single dimension.

\begin{definition}[The five dimensions of correctness]\label{def:five-correctness}
See Table \ref{tab:five-correctness}.
\end{definition}

\begin{table}[htbp]
\centering
\small
\caption{The five dimensions of correctness and their mechanical checkability}\label{tab:five-correctness}
\begin{tabular}{L{2.3cm}L{6.6cm}L{5.5cm}}
\toprule
\textbf{Dimension} & \textbf{Question to be answered} & \textbf{Mechanically checkable?} \\
\midrule
Functional correctness & Do the outputs and the state changes meet the requirements? &
  For an explicit specification, provable or testable \\
Side-effect correctness & Has the program performed unauthorised I/O, writes, network
  requests or state changes? &
  Constrainable by effect specifications, permission boundaries and runtime isolation \\
Resource correctness & Do time, memory, bandwidth and energy meet the budget? &
  Partly provable statically; partly measurable only, or statistically guaranteed in a
  specified environment \\
System correctness & After components are composed, are concurrency, failure recovery
  and state consistency still correct? &
  Partly formally provable; partly requiring model checking, fault injection and runtime
  verification \\
Correctness of goal expression & Is the specification the machine executes really the goal
  the human wants? &
  \textbf{Not derivable automatically from a program-correctness proof} \\
\bottomrule
\end{tabular}
\end{table}

The last row of Table \ref{tab:five-correctness} is the crux of the whole problem: \emph{a
program may implement a wrong specification correctly}. A compiler that is entirely correct
can compile a wrong program into machine code entirely correctly; a formal proof is equally
unable to rescue a specification that omits a critical requirement. The historical
statement of this distinction is the separation of verification from validation in software
requirements engineering \cite{boehm1984}.

\section{The Verifiable Refinement Chain}\label{sec:chain}

\begin{definition}[Verifiable refinement chain]\label{def:chain}
A \term{verifiable refinement chain} consists of seven stages:
\begin{enumerate}[label=(G\arabic*)]
\item \term{Purpose and intent}: what is to be achieved; which outcomes count as success;
      which side effects are unacceptable; how to trade off conflicting goals.
\item \term{System and environment model}: a representation of runtime states, events,
      operations, resources, operating-system interfaces, the external environment and
      observable behaviour.
\item \term{Formal behaviour contract}: permitted behaviour, forbidden behaviour,
      invariants that must hold, progress conditions, resource budgets and environmental
      assumptions, as in Definition \ref{def:contract}.
\item \term{Realisability decision}: check whether the contract is self-contradictory, and
      whether the goal is genuinely achievable under the specified resources, hardware,
      operating system and environmental conditions.
\item \term{Implementation synthesis and search}: generators, search algorithms or other
      synthesisers produce candidate implementations. A candidate implementation is not
      trusted in advance.
\item \term{Correctness proof from abstract behaviour to binary}: check the candidate
      implementation together with its compilation, linking and other transformations, and
      prove that the generated machine code satisfies the behaviour contract in the
      specified environment.
\item \term{Attested deployment and actual execution}: ascertain that what is running is
      the verified binary, execute it in the specified environment, and monitor, constrain
      and handle failures for the conditions that cannot be guaranteed statically.
\end{enumerate}
\end{definition}

The significance of this chain lies not in adding seven process steps but in \emph{assigning
a different mathematical task to each}. The model (G2) should not be a natural-language
description but an interpretable model of system state, behaviour and environment; the
realisability decision (G4) asks whether the requirements are satisfiable; synthesis (G5)
searches for possible implementations; (G6) establishes a proof from the implementation to
the required behaviour; and (G7) ensures that the object actually executed is the object the
proof is about.

\begin{principle}[What crosses a boundary]\label{pr:boundary}
At every boundary what crosses is not only an artefact but also \emph{what that artefact
stands for}, and \emph{the evidence for why it may be accepted}.
\end{principle}

\begin{remark}[The entry point that cannot be mechanically replaced]
Verification results and counterexamples can feed back into (G2), (G3) and (G4); only the
attestation of intent in (G1) cannot be replaced by a general code-correctness proof. The
reason is given in \S\ref{sec:intent-gap}.
\end{remark}

\section{Transition Edges and Evidence}\label{sec:five-edges}

The seven stages of Definition \ref{def:chain} compress into five transition edges, each
with its typical evidence and its critical gap. Table \ref{tab:edges} states them.

\begin{table}[htbp]
\centering
\footnotesize
\caption{Five transition edges, typical evidence and critical gaps}\label{tab:edges}
\begin{tabular}{L{1.5cm}L{3.4cm}L{4.6cm}L{4.9cm}}
\toprule
\textbf{Edge} & \textbf{Input $\to$ output} & \textbf{Typical proof/evidence} &
  \textbf{Key assumptions and principal gaps} \\
\midrule
I$\to$R & Natural language/dialogue $\to$ structured requirements, domain facts and
  items requiring clarification &
  Stakeholder attestation, links to requirement sources, examples and counterexamples,
  conflict reports &
  Intent has no automatically given mathematical truth value; external facts and conflicts
  of interest cannot be settled by type checking alone \\
R$\to$S & Requirements $\to$ formal properties/contracts &
  Traceability matrices, requirement-level tests, consistency/satisfiability/realisability
  analysis, vacuity and mutation analysis &
  A formal mapping must preserve \emph{prohibitions and positive examples} alike; an
  omitted requirement cannot be derived from internal consistency alone \\
S$\to$P & Specification $\to$ program/system candidates &
  Synthesis algorithms, SMT, model checking, Hoare/type/dependent-type proofs,
  counterexample-guided inductive synthesis &
  Finite DSLs, logic fragments and search budgets; over-strong premises and vacuous
  specifications must be prevented \\
P$\to$C$\to$B & Source program $\to$ stage-by-stage intermediate representations $\to$
  machine code/deployable artefact &
  Verified-compiler theorems, per-translation verification/attestation, proof-carrying code &
  Source and target semantics must be defined; linking, loading, assembly, the runtime,
  third-party libraries and hardware are not necessarily within the proof's scope \\
B$\to$E & Binary and configuration $\to$ behaviour in the real environment &
  Deployment integrity, system verification, runtime monitoring, isolation and capability
  mechanisms, performance and scheduling analysis &
  The model must be aligned with the real CPU, operating system, network, scheduling,
  failure and service assumptions \\
\bottomrule
\end{tabular}
\end{table}

\begin{remark}[Heterogeneity of the relations $R_i$]
By Remark \ref{rmk:source}, the interpretation relation on different transition edges need
not be of the same kind. For I$\to$R and R$\to$S the available ``relation'' may be only a
traceability mapping rather than a formal relation \cite{mucha2024,mokos2020}; for
P$\to$C$\to$B it may be a compiler's semantics-preservation theorem
\cite{leroy2009,cakeml2019}, a single translation verification \cite{necula2000} or a
translation attestation \cite{krijnen2024}; for B$\to$E it may be a decision by runtime
monitoring \cite{taleb2023}. Forcing them into one kind of relation is a mistake.
\end{remark}

\begin{remark}[Existing engineering on each edge]
Table \ref{tab:edges} is not a purely speculative scheme: verified compilation
\cite{leroy2009} and verified compiler back ends \cite{cakeml2019} on P$\to$C$\to$B,
kernel-level refinement proofs \cite{klein2009}, and proof-carrying code \cite{necula1997}
have already achieved substantial progress for specific languages, compilers, kernels and
systems. What remains genuinely unresolved is how to extend these high-assurance local
results to a wider range of targets and a greater variety of runtime environments, and to
systems whose implementations may be generated arbitrarily or adjusted dynamically.
\end{remark}

\section{Composition Conditions and Failure Modes}\label{sec:composition}

Proposition \ref{prop:soundness} already gives sufficient conditions for composition. Three
concrete counterexamples of composition failure are added here; their role is to show that
``locally correct'' does not entail ``globally correct''.

\begin{example}[Three composition failures]\label{ex:compose-fail}
\begin{enumerate}[label=(\alph*)]
\item The source program is semantically correct, but the linker selects the wrong library.
\item Each service satisfies its interface contract on its own, but the contracts'
      assumptions about network retry are mutually incompatible.
\item The functional proof holds, but the latency budget does not cover cache contention.
\end{enumerate}
The common structure of the three is that one of conditions (a) or (b) of
Proposition \ref{prop:soundness} is violated.
\end{example}

Read through the contract algebra (Definition \ref{def:operad},
Proposition \ref{prop:operad-compose}), the three counterexamples are no longer mutually
unrelated; they fall into three classes of failure:

\begin{corollary}[Common structure of the three counterexamples]\label{cor:three-failures}
\begin{enumerate}[label=(\alph*)]
\item Linker selects the wrong library: \emph{colour mismatch}, the colour of the linked
      implementation being unequal to the colour declared at the link point;
\item Contracts' assumptions about network retry are mutually incompatible: \emph{a cycle
      appears in the interface graph}, retry turning an acyclic call graph into a cyclic
      one while the colouring is not re-verified after the cycle appears
      (Remark \ref{rmk:feedback});
\item The latency budget does not cover cache contention: \emph{a missing contract
      component}, the resource component $R$ never entering the colour, so that cost is
      treated as an implementation detail.
\end{enumerate}
The three are, respectively, a colour whose extension is too small, a changed interface
graph, and a missing contract component.
\end{corollary}

The conclusion is that composition rules must explicitly incorporate these interface
conditions together with \emph{non-functional conditions}.

\begin{remark}[The holistic character of assurance]\label{rmk:duhem}
Corollary \ref{cor:three-failures}(c) and Remark \ref{rmk:monitor-limit} together show that
no verification conclusion actually asserts ``property $\varphi$ holds''; what it asserts is
the conjunction
\[
\varphi \;\land\; A_1 \;\land\; A_2 \;\land \cdots \land A_k,
\]
where the $A_i$ are the full set of environmental assumptions. The number $k$ cannot be
enumerated in advance, and an assumption $A_i$ becomes visible only when it is actually
violated.
This is the form taken by the Duhem--Quine thesis \cite{quine1951} in software verification:
there is no decisive test of a single proposition, only a test of the whole conjunction. Its
immediate consequence is that the strength of ``the tests pass'' is strictly limited by the
\emph{set of assumptions that were not enumerated}. This also explains why
Table \ref{tab:evidence} must attach a ``what must not be claimed'' list to every class of
evidence: that list is the known $A_i$.
\end{remark}

\section{The Intent Gap}\label{sec:intent-gap}

This is the section that requires the most care.

\begin{problem}[The intent gap]
How is an intent $I$, first expressed by a human in natural language, guaranteed to be
transformed into a correct formal specification $\Phi$? The hope is that $I \equiv \Phi$,
where the equivalence means that ``the behaviour $\Phi$ requires fully reflects the human's
real goal''.
\end{problem}

A popular argument holds that, by G\"odel's incompleteness theorem and Rice's theorem, the
gap between intent and specification is \emph{principled} and \emph{can never be
mechanically closed}. The argument does not hold.

\begin{proposition}[Why this paper does not accept that argument]\label{prop:no-godel}
G\"odel's incompleteness theorem concerns the limits of \emph{formal provability} in formal
systems satisfying certain conditions; it does not directly state that a machine cannot
decide whether a requirements specification matches a particular person's real intent.
Rice's theorem \cite{rice1953} states that non-trivial semantic properties of general
programs are undecidable; that is a fundamental limit on general static analysis, but it is
not a theorem about the relation between a person's mental intent and a formal
specification. Deriving ``intent can never be fully formalised'' from it requires additional
philosophical and cognitive assumptions. Turing's original result \cite{turing1936}
likewise applies only to explicitly encoded computational problems and general algorithmic
claims.
\end{proposition}

Does the intent gap exist, then? It does, but the grounds are \emph{epistemic}, not
\emph{logical}.

\begin{proposition}[Three epistemic grounds for the intent gap]\label{prop:three-grounds}
\begin{enumerate}[label=(\roman*)]
\item \term{Requirements are generative, not enumerative}. A person cannot list their
      requirements completely before seeing the system's behaviour. The requirement ``do
      not leak privacy'' may not surface at all until the system is seen actually
      collecting data. This is not logical undecidability but a constructive limitation of
      human cognition: \emph{preferences are activated by the environment, not present
      beforehand as a complete list}.
\item \term{A specification is lossy compression of intent}. Intent lives in a causal,
      embodied, evolutionarily shaped mind; a specification is a finite symbolic object.
      Lossy compression means that information is necessarily lost, and which parts are
      lost cannot be known in advance. This is a phenomenon of information theory, not of
      metamathematics.
\item \term{The distinction between verification and validation has not been overturned}.
      Verification answers ``does the implementation conform to the specification'';
      validation answers ``does the system meet the real need''. The distinction exists and
      persists precisely because the latter cannot be mechanically reduced to the former.
\end{enumerate}
\end{proposition}

Combining Proposition \ref{prop:no-godel} with \ref{prop:three-grounds} gives the paper's
precise statement of the problem:

\begin{principle}[The precise statement of the intent gap]\label{pr:gap}
The specification-transformation process is not a step that can be mechanically guaranteed
\emph{in general}; the reason is not logical undecidability but that human preferences are
dynamic, situational, and not fully knowable before they are encountered.
\end{principle}

This correction does not weaken the claim; it replaces the load-bearing wall. It also shows
that machines can help clarify requirements, propose counterexamples, generate scenarios,
find contradictions and check a specification for internal consistency, whereas which
specification is finally adopted, and whether it fully expresses the goal, still requires an
identifiable accountable party.

\begin{remark}[Three questions, not one, and what automation can do for each]\label{rmk:three-questions}
The intent gap is often discussed as a single obstacle. It is three, they need different
instruments, and separating them makes clear that the framework denies mechanisability to a smaller
region than it is usually taken to.

\begin{enumerate}[label=(\alph*)]
\item \term{Semantic expression}: given requirements that have in fact been settled, can they be
      written as a formal specification? In bounded domains this is largely automatable, and the
      instruments are the type system, the theorem prover and the specification-synthesis methods
      discussed in Chapter \ref{ch:literature}. Nothing in Proposition \ref{prop:three-grounds}
      denies this, and treating it as denied would be an over-reading.
\item \term{Requirement discovery}: has the interested party found and expressed everything that
      matters? This is Proposition \ref{prop:three-grounds}(i), and its remedy is interaction:
      counterexamples, scenarios, and the probing loop of Chapter \ref{ch:comprehension}. Partially
      automatable, in the sense that machines can generate the probes; not automatable in the sense
      that the answers must come from someone with something at stake.
\item \term{Normative choice}: when requirements conflict---latency against safety, privacy against
      analytical power, cost against performance---who has the standing to decide? This is not an
      epistemic question at all; it is the authority axis of \S\ref{sec:authority}, and the
      instruments are procedural: a designated accountable body and a record of the decision.
\end{enumerate}
Reading the gap this way also corrects a common over-statement in the opposite direction. Because
(a) is largely mechanisable, ``some intent can be precisely defined and automatically checked'' is
true; what cannot be mechanised is the residue---preferences not yet activated, conflicts not yet
surfaced, and decisions nobody has the standing to make alone. The precise claim this paper defends
is therefore narrower than ``intent cannot be formalised'': \emph{a formal proof of implementation
correctness cannot by itself establish whether the specification serves the goals of all relevant
parties}, and the open-world remainder must be handled by an independent validation process. That
process may use automated tools; it may not treat a tool's internal consistency as evidence that the
goal was right.
\end{remark}

\begin{remark}[The intent gap is the frame problem]\label{rmk:frame}
The three grounds of Proposition \ref{prop:three-grounds} are not three independent
findings; they point jointly to an older problem. To formalise an intent one must first
decide which facts are \emph{relevant}; and deciding relevance requires already knowing
which conclusions are to be drawn, which is exactly the frame problem: without knowing the
answer in advance, one cannot enumerate the relevant premises.

The intent gap is therefore not a defect of formal methods but an instance of the frame
problem. The identification has two consequences:
\begin{enumerate}[label=(\alph*)]
\item It explains why stronger formal tools do not solve the problem: tools improve
      \emph{derivation}, whereas the problem lies in \emph{the enumeration of premises};
\item It explains why \emph{interaction} is the only effective means. The frame problem is
      unsolved in general and is evaded only in systems that possess a perception--action
      loop: an agent need not enumerate relevance in advance, because the environment gives
      feedback on its actions. The probing loop of Chapter \ref{ch:comprehension} and the
      traceable refinement of this section both move that evasion mechanism up to the
      requirements layer.
\end{enumerate}
Interactive refinement is therefore forced, not selected from among candidate approaches.
\end{remark}

\begin{remark}[The oblivious set is where the frame problem is paid for, level by level]\label{rmk:frame-per-level}
Remark \ref{rmk:frame} identifies the intent gap with the frame problem, and stops there. The
identification can be pushed one step further, and doing so joins two parts of the paper that have so
far been separate.

The frame problem, in the form used here, is the difficulty of \emph{deciding relevance}: of saying
which facts bear on a conclusion before knowing what is to be concluded. An oblivious set
(Definition \ref{def:oblivious}) is a relevance judgement of exactly that shape, made explicit and
local: $\Obl(\Lambda_i)$ names the differences that $\Lambda_i$ declares not to bear on it. The
relation between the two is therefore not analogy but scale---and it is worth being precise about
what is and is not obtained.

What is obtained is a \emph{localisation}. A frame problem posed for a system as a whole is not
smaller for being posed for one of its levels; but it is bounded by that level's declared
observation map, and it is discharged in a currency the level can state. The global question, which
facts are relevant to the system's purpose, still has no general solution, and nothing in this paper
claims one. What the level structure buys is that the global question is never asked all at once:
each level answers its own bounded version and publishes the answer.

What is obtained second is \emph{testability}. A relevance judgement that is merely held cannot be
criticised; one that is written down can. Section \ref{sec:neardecomp} gives the two ways a declared
judgement can be wrong---the interface may be insufficient for external behaviour, or the observation
map may be incomplete with respect to what is observable in fact. Both are checks, not opinions. In
this sense the oblivious set is the frame problem's engineering form: not a solution to it, but a
place to pay for it where the payment can be audited.

Two consequences follow, and they are the reason the remark is here rather than in Chapter
\ref{ch:comprehension}. First, it explains why the descent protocol of Principle \ref{pr:descent} has
the stopping rule it has: one descends until the invariant governing the change becomes visible,
because that is the depth at which the relevant oblivious sets have been consulted, and descending
further buys relevance that no obligation depends on. Second, it makes the intent layer's
difficulty legible in the same terms as a module boundary's: the intent gap is a frame problem at
$\Lk{\infty}$, which is exactly the level at which no oblivious set can be declared for the machine,
because declaring one would already be deciding the purpose.
\end{remark}

\section{The Asymmetry Between Verification and Generation}\label{sec:asymmetry}

A second formulation requiring correction is the description of ``grading is easier than
sitting the exam'' as ``the universal dividend of $\mathrm{P}$ versus $\mathrm{NP}$''. The
description does not hold: whether $\mathrm{P}$ equals $\mathrm{NP}$ remains open, and it is
not known in general whether finding a solution is necessarily harder than checking one.

This correction must be distinguished from a second point: the open status of
$\mathrm{P}$ versus $\mathrm{NP}$ does \emph{not} make the phenomenon ``verification is
cheaper than generation'' fictitious.

\begin{proposition}[Structural asymmetry]\label{prop:asym-struct}
The structural asymmetry of polynomial-time certificate checking is \emph{the definition of}
the class $\mathrm{NP}$ itself: that a solution can be verified efficiently is a basic fact
about computation, and it does not depend on whether $\mathrm{P} \neq \mathrm{NP}$ has been
proved.
\end{proposition}

\begin{proof}[Argument]
The definition of $\mathrm{NP}$ is ``there exists a polynomial-time verifier $V$ and a
certificate $w$ of polynomial length such that
$x \in L \iff \exists w: V(x,w) = 1$''. The definition itself asserts the existence and the
efficiency of verification, and asserts nothing about the complexity of finding $w$. The
effectiveness of the verifier architecture therefore rests on the definition itself and is
independent of the eventual answer to $\mathrm{P}$ versus $\mathrm{NP}$.
\end{proof}

This cost scale indicates the precise boundary of ``verification is cheap''. The limit is
not the open status of $\mathrm{P}$ versus $\mathrm{NP}$ but \emph{certificate length}:
checking a proof can be done in polynomial time, while the proof itself can be
astronomically long. Hence:

\begin{principle}[The two components of verification cost]\label{pr:two-costs}
``Verification is cheaper than generation'' must be decomposed into two independent
propositions:
\begin{enumerate}[label=(\alph*)]
\item \term{Checking time is cheap}: given a certificate, the decision is efficient
      (Proposition \ref{prop:asym-struct});
\item \term{Certificate size is cheap}: the certificate itself is short enough to be
      stored, transmitted and re-checked.
\end{enumerate}
(a) is guaranteed by the definition of $\mathrm{NP}$; (b) is \emph{not}, and in many
practical systems it is the real bottleneck. A verification regime that satisfies only (a)
is unusable in engineering, because no one can re-check a certificate they cannot finish
reading.
\end{principle}

Principle \ref{pr:two-costs}(b) is measurable: for a given prover and class of properties it
suffices to measure the distribution of certificate sizes. It also explains why a ``small
kernel plus large proof'' architecture must in practice always be paired with a
\emph{readable} proof-summary layer (invariants, prose statements of contracts,
counterexamples) rather than supplying only proof terms.

\begin{example}[Historical cost scale of formal verification]\label{ex:sel4-cost}
The initial functional-correctness proof of the microkernel seL4 took about 11 person-years
(excluding investment in reusable tools and frameworks); the proof scripts of the time ran
to about 200{,}000 lines and verified about 8{,}500 lines of C code \cite{klein2009}. That
scale indicates the historical cost structure of formal verification. AI may substantially
reduce the cost of writing proofs, finding invariants and repairing proofs, but
specification design, model selection, proof architecture and the trusted base remain real
work.
\end{example}

\section{Grades of Evidence}\label{sec:evidence-levels}

Different kinds of evidence provide support of different strength, and conflating them is
the most common error in this area. Table \ref{tab:evidence} gives the grades of evidence
together with the corresponding ``what must not be claimed'' list.

\begin{table}[htbp]
\centering
\footnotesize
\caption{Grades of evidence}\label{tab:evidence}
\begin{tabular}{L{2.6cm}L{5.6cm}L{6.2cm}}
\toprule
\textbf{Type of evidence} & \textbf{Conclusions it can support} &
  \textbf{Conclusions that must not be claimed} \\
\midrule
Machine-checked formal proof & For the stated proposition, holds under the given logic,
  definitions, axioms, libraries and checker assumptions &
  Does not prove that the proposition fully expresses the intent; does not cover unmodelled
  environments; does not by itself establish latency or energy consumption \\
Verified-compiler theorem/translation certificate & Between specific source and target
  semantics, the defined preservation or refinement relation holds, or a single translation
  result is verified &
  Does not mean that the linker, operating system, hardware and external services are
  verified \\
Static analysis and abstract interpretation & Proves or conservatively estimates the
  properties of the analysed program, relative to a sound approximation &
  Does not guarantee precision or the absence of false alarms; is not a measurement of
  actual performance \\
Testing and mutation testing & Finds counterexamples, assesses the specified behaviour on
  the listed tests, and provides evidence over a finite sample &
  Does not prove that all inputs and all executions are safe; a pass rate is not a general
  proof of semantic equivalence of specifications \\
Statistical performance evidence & An empirical distribution under specified data, sampling
  procedure, hardware and confidence method &
  Is not a hard-real-time worst-case execution time; cannot be extrapolated unconditionally
  to a deployment whose distribution has changed \\
Runtime monitoring & Gives an online decision about the observed traces under the monitor's
  semantics &
  If observation is incomplete, ``not detected'' must not be stated as ``all executions
  satisfy'' \\
\bottomrule
\end{tabular}
\end{table}

\begin{principle}[Three states that must be distinguished]\label{pr:three-states}
``A violation was detected'', ``no violation has yet been observed'' and ``all executions
are proved safe'' are three different states, and none may be substituted for another.
\end{principle}

\begin{definition}[Artefact evidence list]\label{def:evidence-package}
Every artefact should carry ten items of evidence, listed in
Table \ref{tab:evidence-package}. The purpose of the list is to make ``may this binary be
trusted'' an auditable question rather than a question about the identity of its author.
\end{definition}

\begin{table}[htbp]
\centering
\small
\caption{Artefact evidence list}\label{tab:evidence-package}
\begin{tabular}{L{0.8cm}L{13.4cm}}
\toprule
\textbf{No.} & \textbf{Evidence item} \\
\midrule
1 & Specification version and requirement traceability identifiers \\
2 & Formal specification and environmental assumptions \\
3 & Hash of the candidate source code \\
4 & Compiler, flags, target architecture and dependency versions \\
5 & Proof and checker versions \\
6 & Resource-analysis model and results \\
7 & Test and monitoring strategy \\
8 & Trusted computing base inventory \\
9 & \textbf{Known uncovered behaviour} \\
10 & \textbf{Re-verification rules} (which changes trigger recomputation of the affected
  evidence) \\
\bottomrule
\end{tabular}
\end{table}

Items 9 and 10 are the two most easily omitted and the most important: the former makes the
boundary of the guarantee explicit, the latter turns assurance from a one-off check into a
lifecycle mechanism.

\section{Side Effects and Constraint Mechanisms}\label{sec:effects}

The example in Problem \ref{prob:trust} mixes side effects of several kinds. Their
mechanical mechanisms become visible only when the kinds are separated.

\begin{definition}[Four classes of side effect]\label{def:effects}
See Table \ref{tab:effects}.
\end{definition}

\begin{table}[htbp]
\centering
\small
\caption{Classification of side effects and the corresponding mechanical mechanisms}\label{tab:effects}
\begin{tabular}{L{1.8cm}L{4.6cm}L{5.0cm}L{3.0cm}}
\toprule
\textbf{Class} & \textbf{Example} & \textbf{Mechanical mechanism} &
  \textbf{Point of verification} \\
\midrule
Functional & A generated function modifies global state it must not modify &
  Type systems, effect types (which make ``what will be touched'' part of the signature),
  ownership systems &
  Compile time, at zero runtime cost \\
Resource & 1\,ns becomes 1\,$\mu$s; memory blows up &
  Fuel/gas metering, worst-case execution time analysis, quotas &
  Both compile time $+$ runtime \\
Security & Unauthorised access, sandbox escape &
  Capability models, sandboxes, bytecode verifiers &
  At admission $+$ at runtime \\
Semantic drift & The specification itself is wrong &
  \textbf{No mechanical mechanism} & Human attestation and consensus \\
\bottomrule
\end{tabular}
\end{table}

The first three classes are already covered by mechanical mechanisms, and are arranged in
depth (proof $>$ types $>$ testing $>$ metering $>$ isolation $>$ monitoring, each layer
having a simpler trust base than the layer above it). Only the fourth class, a specification
that departs from the intent, always requires humans.

\begin{remark}[Assurance scope of three engineering prototypes]\label{rmk:three-prototypes}
Three common engineering mechanisms cover different kinds of property and must not be
advertised side by side:
\begin{enumerate}[label=(\alph*)]
\item A browser sandbox guarantees \emph{execution isolation and a permission boundary}, not
      functional correctness. It is the largest untrusted-code runtime on earth: every day
      it executes hundreds of millions of pieces of code of unknown provenance that have
      never been reviewed; it never reviews, it isolates.
\item A bytecode verifier verifies \emph{the safety properties its rules cover} (no
      out-of-bounds access, bounded loops, no dangling pointers), not the correctness of
      arbitrary business logic. A deterministic verifier of a few tens of thousands of
      lines adjudicates the admission of millions of lines of free code; this is precisely
      an instance of ``compiler-like assurance that does not depend on anyone's will''.
\item A resource-metering mechanism guarantees \emph{constraints under an abstract resource
      accounting}, not real wall-clock time. It turns a performance side effect from
      something ``a human must notice'' into ``a budget violation automatically refuses
      execution'', but it need not catch a degradation from 1\,ns to 1\,$\mu$s, because it
      meters abstract operation costs rather than actual hardware latency. Real-time
      performance assurance requires timing analysis and a restricted runtime environment,
      a completely different engineering discipline.
\end{enumerate}
\end{remark}

\section{Runtime Assurance}\label{sec:runtime}

Formal proofs depend on environmental assumptions, and the actual environment may change.
Beyond static proof, therefore, the executing system must remain safe when an assumption
fails.

\begin{principle}[Layered defence]\label{pr:layered-defense}
An ideal architecture combines static proof (for known deterministic behaviour) with dynamic
monitoring (for unknown or non-deterministic behaviour), forming a layered defence: whatever
can be proved is proved first; residual risk is detected by independent monitoring; serious
violations have their impact limited by isolation or rollback.
\end{principle}

\begin{remark}[The monitor's own boundary]\label{rmk:monitor-limit}
Even a runtime monitor that is itself correct does not thereby possess the complete facts of
an execution. Missing observation traces, ambiguous events, unreliable transport and
concurrency that cannot be seen all limit its conclusions. The observation mechanism of the
monitoring system must therefore itself be included in the assurance chain.
\end{remark}

\begin{remark}[Four conditions for contract composition]\label{rmk:contract-conds}
The composition of local properties depends on four things: the environment genuinely
satisfies the subcomponent assumptions; the connection topology is well posed; the feedback
condition is checkable; and the observation and time semantics are compatible. ``Individually
correct'' for components and ``globally correct'' for the system are not interchangeable.
\end{remark}

\begin{principle}[The conditional form of an invariant]\label{pr:ag}
A usable assume--guarantee contract should state simultaneously: the environmental
conditions it depends on, the behaviour it guarantees outwards, the side effects it may
produce, the resources it must consume and occupy, and the constraints that must hold when
it is composed with other subsystems.
\end{principle}

\section{Artefact Binding}\label{sec:fourth}

The last conjunct of the assurance condition \eqref{eq:assurance}, namely that the object
actually executed is the object verified, is frequently omitted from existing discussions.
It is not a marginal detail: even if a verifier correctly accepts a binary, the preceding
verification conclusion does not transfer automatically if the deployment system uses a
different file, or if the loading environment changes a critical premise. Supply-chain
attacks aim precisely at this crack.

The complete trust chain must therefore include content-addressed build outputs, a
cryptographic binding between the verification evidence and the binary, and re-verification
at deployment. In the notation of \eqref{eq:binding},
$\mathrm{digest}(B_{\text{deployed}}) = \mathrm{digest}(B_{\text{verified}})$
and $E_{\text{deployed}} \models A_{\text{claimed}}$.

\begin{principle}[A three-way distinction for reproducible builds]\label{pr:reproducible}
Three properties of different strength must be distinguished:
\begin{enumerate}[label=(\alph*)]
\item \term{Reproducible}: the same inputs produce a byte-for-byte identical output;
\item \term{Verifiable}: an independent third party can rebuild from the inputs and compare;
\item \term{Verified}: a proof that a particular binary satisfies its declared properties.
\end{enumerate}
(a) and (b) are preconditions for (c) to take effect; (c) does not imply (a). The three are
not interchangeable.
\end{principle}

\section{The Reliability Bound of the Consensus Layer}\label{sec:consensus}

\S\ref{sec:trust} showed that the intent layer admits no mechanical solution. What does it
admit? This section gives a precise answer and proves a counter-intuitive result.

\begin{theorem}[Consensus reliability bound]\label{th:condorcet}
Let $n$ verifiers judge independently the proposition ``whether a specification faithfully
expresses the intent'', each erring with probability at most $p < 1/2$. Then the probability
that majority voting errs satisfies
\begin{equation}\label{eq:chernoff}
\Pr\bigl[\text{majority errs}\bigr] \;\le\; \exp\Bigl(-2n\Bigl(\tfrac{1}{2} - p\Bigr)^{2}\Bigr).
\end{equation}
\end{theorem}

\begin{proof}
Let $X$ be the number of verifiers who err, $X = \sum_{i=1}^{n} X_i$, with the $X_i$
independent and $\Pr[X_i = 1] \le p$, so $\mathbb{E}[X] \le np$. A majority errs exactly
when $X \ge n/2$. By Hoeffding's inequality \cite{hoeffding1963}, for $t > 0$,
$\Pr[X - \mathbb{E}[X] \ge t] \le \exp(-2t^{2}/n)$. Taking
$t = n(1/2 - p)$ (so that $X \ge n/2$ implies $X - \mathbb{E}[X] \ge t$) gives
\[
\Pr\Bigl[X \ge \tfrac{n}{2}\Bigr] \;\le\;
\exp\Bigl(-\frac{2n^{2}(1/2 - p)^{2}}{n}\Bigr) \;=\;
\exp\Bigl(-2n\bigl(\tfrac{1}{2} - p\bigr)^{2}\Bigr). \qedhere
\]
\end{proof}

This result is the quantitative form of the jury theorem \cite{condorcet1785}. It relies on
independence, however, and independence is a fragile assumption; its failure has an equally
precise form.

\begin{proposition}[Bound dominated by a correlated error mode]\label{th:correlated}
Suppose that the errors of $n$ verifiers have the following structure: with probability $q$
a \emph{common error mode} $C$ occurs (for instance a shared training-data bias, a shared
domain misunderstanding or a shared incentive), in which case all verifiers err
simultaneously; with probability $1-q$ the verifiers err independently, with error rate
$p' < 1/2$. Then
\begin{equation}\label{eq:correlated}
\Pr\bigl[\text{majority errs}\bigr] \;\ge\; q .
\end{equation}
In particular, when $q$ is bounded below by a positive constant, increasing $n$ cannot push
the probability that the majority errs below $q$.
\end{proposition}

\begin{proof}
When $C = 1$ all $n$ verifiers err, so the majority necessarily errs. Hence
$\Pr[\text{majority errs}] \ge \Pr[C = 1] = q$. Taking the limit in $n$ does not change
this lower bound.
\end{proof}

\begin{corollary}[Diversity, not number]\label{cor:diversity}
By Proposition \ref{th:correlated}, the governing variable of consensus reliability is
\emph{the independence of errors} (that is, the size of $q$), not the number of verifiers.
The central goal of consensus design is therefore to reduce $q$: admit participants with
different backgrounds, models with different training provenance, reviewers with different
interests, and explicitly designed \emph{adversarial} roles.
\end{corollary}

\begin{remark}[Use of this result]\label{rmk:diversity-use}
Corollary \ref{cor:diversity} yields a criterion directly usable in engineering judgement:
when someone proposes ``adding a few more agents to review this'', the right question is
``are the new reviewers independent of the existing ones in their error modes''. If they are
not, increasing the number increases cost without increasing reliability. This also explains
why ``reviewing with another model from the same source'' does not constitute a guarantee
structurally.
\end{remark}

\begin{remark}[What the consensus model is, and what it is not]\label{rmk:consensus-status}
Theorem \ref{th:condorcet}, Proposition \ref{th:correlated} and the design-effect approximation
below are statements \emph{within} stated probability models. They should not be read as predictions
of the reliability of a real review process, and the gap between model and practice is worth being
explicit about, because it is where the framework is most likely to be misused.

Three caveats. First, the models assume a common proposition to be judged, independent error sources
modulo $q$, and a fixed individual competence. Real attestation exercises violate all three: what
counts as a faithful specification is itself under negotiation, attestors influence one another by
talking, and competence varies with the clause being judged. The literature on jury theorems is
largely a study of how fragile these premises are---in particular, independence and competence are
not simultaneously justifiable under any single notion of probability \cite{dietrich2008}---and a
shared body of evidence is the most natural mechanism by which apparent independence fails, since it
makes the attestors' errors share a common cause \cite{dietrich2004}. The second of those results is
the formal counterpart of what Principle \ref{pr:q-decomposition} calls information isomorphism.

Second, a single parameter $q$ does not characterise dependence in general. Two attestors can be
correlated in one clause and independent in another, correlated through shared training in one
subdomain and through shared incentives in another, and a one-dimensional summary averages over all
of it. $q$ is the right parameter for the bound and a lossy summary of the phenomenon.

Third, and most consequentially, $p$ and $q$ are not estimable from data unless ``faithful to the
intent'' has a stable standard of evaluation, a designated adjudicator and a specification of what
counts as an error. Where those exist---bounded domains with observable acceptance criteria---the
quantities can be estimated from outcomes, which is the case Hypotheses $H_8$ and $H_{11}$ address.
Where they do not, the bound remains a conditional statement and no quantity of attestation data
supplies the missing standard. The honest summary is that the model tells you what to measure if the
measurement is possible, and tells you the shape of the answer if it is.
\end{remark}

\begin{remark}[Partial correlation: the design effect]\label{rmk:design-effect}
Theorem \ref{th:condorcet} and Proposition \ref{th:correlated} describe two extremes: independent
errors, giving a bound that improves exponentially in $n$, and a common error mode, giving a floor
$q$ that no $n$ can breach. A real attestation network sits between them, with errors partially
correlated in pairs. The standard statistical instrument applies, and is worth writing down because
it makes Corollary \ref{cor:diversity} quantitative rather than rhetorical.

Let $\bar\rho$ be the mean pairwise correlation of attestor errors. The customary
\emph{design-effect} approximation \cite{kish1965} replaces the nominal attestor count by an
effective one,
\begin{equation}\label{eq:design-effect}
n_{\mathrm{eff}} \;=\; \frac{n}{1 + (n-1)\bar\rho},
\end{equation}
and one reads the bound of Theorem \ref{th:condorcet} at $n_{\mathrm{eff}}$ in place of $n$. Two
consequences follow.
\begin{enumerate}[label=(\alph*)]
\item As soon as $n\bar\rho \gg 1$, \eqref{eq:design-effect} gives
      $n_{\mathrm{eff}} \to 1/\bar\rho$, \emph{independent of $n$}. Beyond that point additional
      attestors buy nothing; the ceiling is set by the correlation alone.
\item Lowering $\bar\rho$ raises that ceiling, and the trade is explicit: once $n\bar\rho$ is large,
      halving $\bar\rho$ roughly doubles the achievable $n_{\mathrm{eff}}$. Diversity is a budget
      line, not a slogan.
\end{enumerate}
The relation to Proposition \ref{th:correlated} is that $q$ is the degenerate case
$\bar\rho \to 1$: with complete correlation $n_{\mathrm{eff}} \to 1$ and \eqref{eq:design-effect}
reproduces the floor.

Two caveats belong with the remark. First, \eqref{eq:design-effect} is the standard approximation
for the variance of a mean under exchangeable correlation, not a theorem about majority error; a
rigorous bound for the intermediate regime needs an explicit correlation model, and none is claimed
here. Second---and this is what makes the approximation usable---$\bar\rho$ is \emph{estimable} from
data the framework already collects: the historical pairwise correlation of attestor errors, settled
against subsequently observed defects and rework. This does not answer the open problem
\ref{op:q} of measuring $q$; it supplies a companion quantity that is measurable, while $q$ remains
the limit that is not.
\end{remark}

\subsection{Why the two layers cannot be merged, and what that forces}\label{sec:enrichment}

Definition \ref{def:two-layer} asserts that the chain has a mechanical layer and a consensus layer
and that the two are stitched rather than fused. That assertion has a mathematical reason, and
stating it settles a design question that is otherwise answered by preference: why a single, more
capable verifier cannot replace the pair.

The reason is that the two layers do not merely differ in strength; they assign their guarantees to
\emph{different kinds of object}. A mechanical guarantee is a proof: an arrow either carries one or
it does not, and one proof may refine another. A consensus guarantee is a number: an unreliability
$\varepsilon$, with no question of one $\varepsilon$ refining another. Formally the two layers are
two \emph{enrichments} of the same underlying structure of stages and artefacts---hom-objects drawn
from the two-element order in the first case, from the Lawvere quantale
$V = ([0,\infty],\geq,+,0)$ in the second, where the composite of two stages has unreliability at
most the sum of theirs. That composition law is the union bound, and it deserves emphasis that it
needs no independence assumption: if stages $1$ and $2$ are each faithful with probability at least
$1-\varepsilon_i$, then both are faithful with probability at least $1-\varepsilon_1-\varepsilon_2$
whatever the dependence between them. This is worth contrasting with the bound of
Theorem \ref{th:condorcet}, whose exponential improvement requires independence and which the whole
of \S\ref{sec:consensus} exists to qualify.

The asymmetry between the two enrichments is not symmetric. There is a change of base from proof to
reliability---a discharged obligation has unreliability $0$, an undischarged one has
unreliability $\infty$---so a proof may legitimately be weakened into a reliability claim. The
reverse fails, and it fails for a reason that is two lines long. If
$G : ([0,\infty],\geq,+,0) \to (\{0,1\},\leq,\wedge,1)$ respects the monoidal structure, then for
every $\varepsilon$,
\begin{equation}\label{eq:doubling}
G(2\varepsilon) \;=\; G(\varepsilon+\varepsilon) \;=\; G(\varepsilon)\wedge G(\varepsilon)
\;=\; G(\varepsilon),
\end{equation}
so $G$ cannot distinguish $\varepsilon$ from $2^n\varepsilon$ for any $n$. Nothing is lost in this
computation by weakening it: functoriality alone gives $G(2\varepsilon) \leq G(\varepsilon)$ and the
comparison map gives the reverse inequality, so \eqref{eq:doubling} holds for lax monoidal $G$ as
well as strong ones. The form checked in the companion development is the strict one,
$G(a+b) = G(a)\wedge G(b) \Rightarrow G(2\varepsilon) = G(\varepsilon)$; the lax strengthening is the
two-line argument just given and is not itself checked.

The monoidal maps from $V$ to the two-element order are exactly three: the constant map to the top
element; the map $\varepsilon \mapsto [\varepsilon < \infty]$, which is the left adjoint of the
inclusion of the two-element order and is Lawvere's $\pi_0$; and the map
$\varepsilon \mapsto [\varepsilon = 0]$, which is the right adjoint \cite{lawvere1973}. \emph{Each of the three is constant on the finite positive values}---the first two take the top
element there and the third takes the bottom---so none of them separates two such values. That is
the consequence this chapter uses: \emph{no monoidal map separates two finite positive
unreliabilities}. A single threshold on $\varepsilon$ is therefore not merely a lossy summary; it is
not a change of base at all. With $\theta$ for the threshold and $G(\varepsilon) = [\varepsilon \leq
\theta]$, two stages at $\varepsilon = 0.6\,\theta$ each pass while their composite at
$1.2\,\theta$ does not, so $G$ fails \eqref{eq:doubling} and with it the composition law.

\begin{principle}[Reliability is budgeted, not gated]\label{pr:budget}
A chain of $k$ stages that must \emph{all} hold has composite unreliability at most
$\sum_{i=1}^{k}\varepsilon_i$. Gating each stage against a fixed threshold $\theta$ therefore gives a
composite bound of $k\theta$, which grows without limit in $k$; the bound must instead be
\emph{divided}, with each stage held to $\varepsilon_i \leq \Theta/k$ for a total budget $\Theta$.
\end{principle}

The practice the principle rules out is common enough to be worth stating numerically. Take
``$99\%$ reliable'' as the per-clause gate, $\varepsilon \leq 0.01$, and a specification with $k$
clauses each attested on its own---attested separately, not with independent errors, since the union
bound assumed no independence. The composite bound $1-k\varepsilon$ is $0.90$ at $k = 10$, $0.50$ at
$k = 50$, and \emph{zero at $k = 100$}---no information at all. Real specifications have hundreds of
clauses, so clause-wise gating at a fixed per-clause reliability does not merely weaken the
guarantee; past a hundred clauses it says nothing, while continuing to look like a guarantee. The
bound must be divided: to have any content at $k = 100$ the per-clause requirement is
$\varepsilon < 0.01$, strictly below $1/k$. This is the same shape as the budget line of Remark
\ref{rmk:design-effect}---which already says of diversity that it is a budget---applied to a
different quantity, and it is the shape the design-effect correction takes across stages rather than
within one: the design effect governs correlation \emph{among attestors of one stage}, the union
bound governs composition \emph{across stages}.

Two boundary conditions on the claim. First, it concerns the composition of \emph{reliability}; it
does not say that thresholds are useless, only that a threshold applied per stage is not a
compositional summary, and the correct composite statement is the $\varepsilon$-indexed family
$[\varepsilon \leq \theta]$ taken for each $\theta$ rather than one chosen $\theta$. That family is
how a quantitative statement is expressed in Boolean logic without losing its content, and the
device is standard in quantitative equational logic \cite{bacci2023}. The failure mode of a single
threshold has a standard name in the bisimulation-metric literature as well: approximate
bisimulations are not transitive in general, which is the abstract form of the
$0.6\theta + 0.6\theta > \theta$ computation \cite{gebler2013}.

Second, the principle bounds what a chain can claim, not what a single stage can achieve. The
$\varepsilon_i$ are properties of one attestation round given a design and a correlation structure;
the design effect is how they are estimated. The contribution here is therefore narrow and should be
stated narrowly: not a better bound, but the observation that the sum bound, which is often treated
as a crude fallback, is the composition law \emph{forced} by the enrichment, and that the
alternatives to it are not better bounds but a different kind of object. This is worth saying because
the opposite is the received view at exactly the point where the framework needs it. Rushby, in the
analysis of assurance cases closest to the present chapter, states the mismatch between probabilistic
evidence and deductive argument, gives the sum rule and the product rule, judges both ``too crude and
pessimistic'', and leaves the repair as an open problem \cite{rushby2015}. The refinement this
section offers is that the sum rule is not a crude approximation to be improved on---it is what
composition of unreliabilities \emph{is}---and that the apparent crudeness lies in the per-stage
$\varepsilon$, where the design-effect correction and the incentive analysis of \S\ref{sec:mechanism}
do their work, rather than in the rule that combines them.

One point of contact with the software-engineering literature is worth recording, because it is the
empirical counterpart of the same phenomenon. Attestors drawn from one interest class fail together,
and the standard model for that is design diversity, where independently produced versions are found
to exhibit coincident failures far more often than independence would allow; the review of
Littlewood, Popov and Strigini is the standard reference for those models, and Knight and Leveson's
experiment is the origin of the empirical objection \cite{littlewood2001,knight1986}. Read against
\eqref{eq:doubling}, that literature is measuring $\bar\rho$ in the sense of Remark
\ref{rmk:design-effect} rather than bounding it, which is what makes the framework's reliance on an
estimable companion quantity rather than on $q$ itself a practical choice and not merely a
concession.

\section{The Incentive Structure of the Intent Layer}\label{sec:mechanism}

Proposition \ref{th:correlated} attributes consensus reliability to the parameter $q$ but
gives no method of estimating $q$. This section identifies the nature of that omission: the
principal sources of $q$ are \emph{collusion} and \emph{shared interest}, so it is not a
statistical parameter but a problem of \emph{incentive design}. Appendix \ref{app:protocol} turns
the resulting requirements into a procedure and states its own failure modes; it is a sketch of an
unbuilt mechanism, not a validated one, and is offered so that the claim that this intersection is
unoccupied can be assessed against something concrete.

\begin{principle}[The composition of $q$]\label{pr:q-decomposition}
The common error mode decomposes into three sources:
\begin{enumerate}[label=(\alph*)]
\item \term{Cognitive isomorphism}: the verifiers share training provenance, domain
      assumptions or habits of thought;
\item \term{Interest isomorphism}: the verifiers' payoff is tied to ``passing'' rather than
      to ``being right'';
\item \term{Information isomorphism}: the verifiers read the same material under review, so
      their errors are determined by the same defect.
\end{enumerate}
The three sources correspond to three different interventions: admitting heterogeneous
participants (for (a)), changing the payoff structure (for (b)), and introducing independent
information channels (for (c)). Acting on the first alone is the common practice and the
least effective. \S\ref{sec:authority} identifies a fourth co-occurrence of the same family,
\emph{authority isomorphism}, and explains why it is not a fourth term of this decomposition but a
separate failure axis.
\end{principle}

\begin{principle}[Proper scoring rules]\label{pr:scoring}
For the second source, the correct remedy is to tie a verifier's reward to the
\emph{accuracy} of the judgement rather than to its \emph{lean}. A proper scoring rule
\cite{brier1950} supplies a ready-made instrument: the reward function attains its unique
maximum at the true belief, thereby removing the incentive for both ``going along with the
majority'' and ``pleasing the submitter''. A concrete form is a prediction market on the
proposition ``does this specification faithfully express the intent'', settled on the
defects and rework data observed subsequently. The mechanism has two by-products:
\begin{enumerate}[label=(\alph*)]
\item It provides a means of reducing $q$ other than replacing people;
\item It yields a posterior estimate of $q$: it suffices to measure the historical
      correlation of verifier errors.
\end{enumerate}
\end{principle}

This section closes the gap left by \S\ref{sec:consensus}. The bound there is conditional,
and its condition (independence) is here reduced to a quantity that can be designed,
measured and paid for. Chapter \ref{ch:program} writes it as hypothesis $H_{11}$.

\section{Who May Attest: the Deontic Axis}\label{sec:authority}

Sections \ref{sec:consensus} and \ref{sec:mechanism} treat attestation as an epistemic problem: is
the specification faithful to the intent, and how reliable is the majority that says so. A second
question exists, and the framework's own assurance condition already contains it without having
connected it. The second conjunct of \eqref{eq:assurance} requires that \emph{the contract is
explicit and authorised}. Authorisation is not a matter of reliability: a specification can be
faithfully attested, by a large and genuinely independent panel, and still be one that nobody had
the standing to adopt.

\begin{principle}[Authorisation is a separate failure axis]\label{pr:authority}
The reliability of an attestation network concerns whether its verdict is \emph{likely correct}.
Authorisation concerns whether the verdict \emph{decides anything}. The two are independent.
\begin{enumerate}[label=(\alph*)]
\item A verdict can be reliable and unauthorised: a competent panel drawn entirely from one
      stakeholder class, agreeing on a requirement that another class had the standing to refuse.
\item A verdict can be authorised and unreliable: a single accountable owner making an uninformed
      judgement.
\end{enumerate}
Consequently no improvement in $q$, in $n_{\mathrm{eff}}$, or in the incentive structure can repair an
authorisation failure. The failure is not that everyone is wrong together, but that everyone together
is not the body entitled to decide.
\end{principle}

This is why authority belongs \emph{beside} the three sources of $q$ rather than inside them.
Principle \ref{pr:q-decomposition} decomposes the common error mode into cognitive, interest and
information isomorphism---three ways for attestors to be wrong in the same way. A fourth
co-occurrence is easy to name, and call \term{authority isomorphism}: all attestors drawn from a
single stakeholder class. It does raise $q$, because one class tends to share both premises and
interests. But lowering $q$ does not address it. A perfectly diverse panel of outsiders is exactly
the case in which the authorisation failure is cleanest: they may be right and still lack standing.

The distinction has a well-worn institutional form, worth one sentence and no more. Distributed
ledger systems separate the mechanism that establishes \emph{what is true} from the mechanism that
establishes \emph{who may decide what counts as true}; consensus and governance are different
protocols, and a system with excellent consensus and no governance is not thereby legitimate. The
analogy is offered as an analogy. What it points at is a gap in this paper rather than a result in it.

\begin{remark}[Two consequences for the design of the intent layer]\label{rmk:authority-design}
\begin{enumerate}[label=(\alph*)]
\item The attestation record of \S\ref{sec:consensus} should carry the \emph{stakeholder class} of
      each attestor beside its provenance. That costs one field, makes authority isomorphism
      observable, and lets an authorisation failure be distinguished after the fact from a
      reliability failure---the two have different remedies.
\item The choice of the authorised body is not derivable from the framework. It is an institutional
      fact about who bears the consequences of the decision, and is fixed by the same argument that
      assigns $\Lk{\infty}$ to human bearers in Chapter \ref{ch:levels}---an argument whose premise
      is the prevailing liability regime, as that chapter states explicitly.
\end{enumerate}
\end{remark}

\section{Specification Binding}\label{sec:spec-binding}

The fourth pillar of \eqref{eq:assurance} binds the verified artefact to the executed artefact.
Nothing so far binds the specification to the intent it was meant to express, and the asymmetry
matters, because the two have opposite behaviour over time. A binary is byte-identical or it is not:
its binding is static and decidable once and for all. A specification is faithful to an intent that
moves. Proposition \ref{prop:three-grounds}(i) is the reason---requirements are generative, and
preferences are activated by situations that did not exist when the specification was written---so
fidelity is not a property a specification has but a quantity that decays.

Table \ref{tab:evidence-package} already requires item 10, re-verification rules, of every artefact:
which changes trigger recomputation of the affected evidence. The specification side has no
counterpart, and condition (i) of Proposition \ref{prop:assurance} is stated as though it needed none. The
repair is to make the specification carry its own binding record.

\begin{definition}[Specification binding]\label{def:spec-binding}
A \term{specification binding} is a four-tuple
\[
\mathrm{SB}(\Phi) \;=\; \bigl(\text{specification version},\; \text{attestation record},\;
\text{validity horizon},\; \text{re-confirmation triggers}\bigr),
\]
where the attestation record is the consensus artefact of \S\ref{sec:consensus} that established
$\Phi$'s fidelity when it was written, the validity horizon is the interval over which that
attestation is to be treated as current, and the re-confirmation triggers are the events that
invalidate it before the horizon expires.
\end{definition}

The triggers are the interesting part, and they are not arbitrary. Proposition \ref{prop:three-grounds}
names three sources of drift, and each suggests its own trigger.
\begin{enumerate}[label=(\alph*)]
\item \emph{Generative requirements}: a new class of user, scenario or deployment environment appears,
      which is precisely the situation that activates a preference nobody had expressed.
\item \emph{Lossy compression}: the domain model, ontology or controlled language in which $\Phi$ was
      written changes, so that the same text denotes something else.
\item \emph{Validation}: an observed incident whose resolution is not derivable from $\Phi$, which is
      evidence that $\Phi$ under-specifies the intent rather than that the implementation erred.
\end{enumerate}
Trigger (c) is worth singling out because it is available for free: an incident that required a human
judgement call is, by construction, an incident the specification did not settle.

Condition (i) of Proposition \ref{prop:assurance} should therefore be read with a time index. Writing
$\varepsilon_{\text{cons}}(\Delta t)$ for the consensus error after an interval $\Delta t$ since
attestation,
\begin{equation}\label{eq:spec-decay}
\Pr\bigl[\Phi \text{ is faithful to } I \text{ at time } t\bigr] \;\ge\;
1 - \varepsilon_{\text{cons}}(\Delta t),
\qquad \Delta t = t - t_{\text{attest}},
\end{equation}
which reduces to the original condition at $\Delta t = 0$. This paper does not give the form of
$\varepsilon_{\text{cons}}(\cdot)$: no decay model is proposed, and the quantity is left as a
measurable object rather than a modelled one. What the section establishes is the structural
symmetry---a specification needs a re-confirmation rule exactly as a binary needs a re-verification
rule---and that the symmetry is forced by Proposition \ref{prop:three-grounds}(i) rather than added
for completeness.

\begin{remark}[The specification-side dual of the artefact-binding gap]\label{rmk:spec-h9}
Hypothesis $H_9$ concerns the measurable risk of the artefact-binding gap. Its specification-side
dual is the risk of executing against a stale-but-attested specification: the attestation record is
still inside its horizon, the binary still matches its digest, and both are wrong together because
the intent moved. The two gaps are independent, and a deployment can close either while leaving the
other open.
\end{remark}

\section{The Consolidated Assurance Condition}\label{sec:assurance-final}

Combining the parts of this chapter yields a complete, conditional statement of assurance.

\begin{definition}[Two-layer trust structure]\label{def:two-layer}
The chain from intent to execution is a composite of a two-layer trust structure:
\begin{equation}\label{eq:twolayer}
\underbrace{I \xrightarrow{\;\text{consensus}\;} \Phi}_{\text{non-formalisable layer}}
\;\longrightarrow\;
\underbrace{\Phi \xrightarrow{\;\text{mechanical}\;} P \xrightarrow{\;\text{mechanical}\;} B
\xrightarrow{\;\text{mechanical}+\text{monitoring}\;} E}_{\text{formalisable layer}}
\end{equation}
In the formalisable layer, trust is guaranteed by mechanical verifiers; in the
non-formalisable layer, trust is guaranteed by adversarial consensus; and the boundary
between the two layers, the record of the consensus, and the traceability of the whole chain
are guaranteed mechanically by cryptographically anchored infrastructure.
\end{definition}

\begin{proposition}[Consolidated assurance: sufficient conditions]\label{prop:assurance}
Let $I$ be an intent, $\Phi$ an attested specification, $B$ a candidate implementation
(binary) and $E$ a runtime environment, with $A$ the declared environmental assumptions.
If the following conditions hold simultaneously, then ``property $\varphi$ holds within the
declared scope $\Sigma$'' holds:
\begin{align}
&\text{(i)}\quad \Pr\bigl[\Phi \text{ is faithful to } I \text{ at time } t\bigr] \;\ge\;
1 - \varepsilon_{\text{cons}}(\Delta t);\\
&\text{(ii)}\quad \Tr(B, E) \;\subseteq\; \Allow(\Phi);\\
&\text{(iii)}\quad E \models A;\\
&\text{(iv)}\quad \mathrm{digest}(B_{\text{deployed}}) = \mathrm{digest}(B_{\text{verified}});\\
&\text{(v)}\quad \varphi \text{ is preserved under the interpretation relation } R \text{ in use};\\
&\text{(vi)}\quad \Phi \text{ is authorised by a body with standing to adopt it,}
\end{align}
where $\varepsilon_{\text{cons}}(\cdot)$ is the bound of \eqref{eq:chernoff}, governed by $q$ and
$\bar\rho$, evaluated at $\Delta t = t - t_{\text{attest}}$.
\end{proposition}

The statement is deliberately one-directional. It was written elsewhere as an equivalence, and
that was wrong, for a reason worth stating because it recurs wherever a framework of this shape is
proposed. The six conditions are \emph{evidence} conditions: they say what must be in hand. They are
not truth conditions, and truth does not wait for them. A property can hold although no evidence for
it has been assembled, so the converse direction fails on the elementary ground that a system may be
correct and undocumented. What the proposition asserts is the useful direction---that holding all six
is enough---and what it deliberately does not assert is that holding fewer is impossible.

Each conjunct carries its own caveat, and none of them is mechanical in the same sense as the
others.
\begin{enumerate}[label=(\alph*)]
\item Condition (i) is a probability with no natural definition until faithfulness is defined, an
      adjudicator identified, and a class of errors specified. Until those exist, $\varepsilon$ is a
      parameter of a model rather than a measured frequency. Its reliability is governed by $q$ and
      by the correlation $\bar\rho$ of \S\ref{sec:consensus} and \S\ref{sec:mechanism}, and both are
      conditional-model quantities.
\item Condition (ii) is a semantic property, not a routine check. $\Tr(B,E) \subseteq \Allow(\Phi)$
      is decidable outright only for restricted languages, finite models or programs carrying
      checkable certificates; for general programs in general environments it is a proof obligation,
      discharged---when it is discharged at all---by a combination of typing, model checking,
      contract verification and testing of graded strength.
\item Conditions (iii) and (iv) are the mechanical pair, and their strength is limited by the trusted
      computing base and by the model of the environment, not by anything in this chapter.
\item Condition (v) is where the relational premise of Remark \ref{rmk:relational-compose} is spent.
\item Condition (vi) is deontic and is not a matter of reliability at all (\S\ref{sec:authority}).
\end{enumerate}
Three of the six conjuncts of Definition \ref{def:assurance} are therefore mechanical, one is
probabilistic, one is deontic, and one---the requirement that the contract be explicit---is what the
rest of the chapter is for.

One further distinction the consolidated statement must carry is the classical one between safety and
liveness. Condition (ii) constrains finite traces: nothing forbidden happens. Progress properties---a
request is eventually answered, a lock is eventually released, a queue eventually drains---are
liveness properties, they are not instances of ``no finite prefix violates $\Phi$'', and they are not
obtained by the same argument. Where a system's contract includes liveness, the obligations belong at
the levels and to the devices that can express them: fairness assumptions and well-foundedness
arguments for the protocol model, and monitoring with a bounded horizon for the deployed system. A
consolidated statement that quietly absorbed liveness into trace inclusion would be wrong in a way
that is easy to miss, and this one does not.

Equation \eqref{eq:twolayer} and Proposition \ref{prop:assurance} together give the complete
answer to Problem \ref{prob:trust}:

\begin{principle}[Answer to Problem \ref{prob:trust}]\label{pr:answer}
``Humans no longer read the code'' is attainable; ``humans need no longer answer for the
expression of the goal'' is not. The workable division of labour is: humans establish and
bear responsibility for the goal and the constraints; generators perform large-scale
generation, search and optimisation; mechanised verification and the execution boundary
judge and constrain implementations; and the system always carries traceable assurance
evidence.
\end{principle}

\begin{remark}[Summary of conclusions on the central questions Q1--Q8]\label{rmk:q1q8}
\begin{longtable}{L{5.4cm}L{9.0cm}}
\toprule
\textbf{Question} & \textbf{This paper's judgement} \\
\midrule
\endhead
Q1 Can it be proved that the intent and the specification agree? &
  Not by general automatic methods. Controlled languages, domain ontologies, examples and
  human attestation can narrow the problem, but cannot cover implicit requirements that were
  never expressed (Proposition \ref{prop:three-grounds}) \\
Q2 How can a specification be proved sufficient? &
  Consistency, satisfiability, realisability, vacuity and formal coverage metrics can all be
  checked; none of them proves that the specification contains every undeclared prohibition \\
Q3 How do local guarantees yield an end-to-end guarantee? &
  It requires that relations compose, that environmental assumptions close and that
  observation maps are compatible (Proposition \ref{prop:soundness}); liveness and real-time
  properties usually require additional conditions \\
Q4 Can functional and resource guarantees be unified? &
  Several constraints can be expressed in a joint behaviour--cost semantics, but time,
  memory, energy and information flow need different analysis models; joint expression is
  not unified solving \\
Q5 How can trust in the generative model be reduced? &
  The generator proposes candidates and an independent checker verifies them. This reduces
  dependence on the generator's correctness but does not remove the requirement to trust the
  logical kernel, the theorem statements, the external assumptions and the correspondence of
  artefacts \\
Q6 Can a given specification always be synthesised effectively? &
  No. Finite grammars, restricted logics and specific templates make some problems solvable;
  for arbitrary program semantic properties there are boundaries of undecidability and
  search-space explosion \\
Q7 Can actual latency be strictly guaranteed? &
  A usable worst-case execution time bound is obtainable under specific assumptions about
  hardware, inputs, scheduling and the cost model; benchmarks and asymptotic complexity are
  not by themselves sufficient to prove a bound at the nanosecond scale \\
Q8 Is a unified theory needed? &
  The better-founded choice is to establish shared proof interfaces and evidence formats
  while retaining the different semantics, assumptions and analysis tools of each domain. A
  unified formalism may be one goal, but it is not a necessary precondition \\
\bottomrule
\end{longtable}
\end{remark}

The empirical content of this chapter is directions 1--3 (Chapter \ref{ch:program}) and
hypotheses $H_7$, $H_8$ and $H_{11}$. Propositions \ref{prop:no-godel} and
\ref{prop:three-grounds} are arguments rather than theorems, and what can be challenged in
them is whether the three grounds are exhaustive. Theorem \ref{th:condorcet} and
Proposition \ref{th:correlated} are applications of standard probability theory, and their
conclusions depend on the independence assumption.

% =====================================================================
\chapter{Levels as Units of Verification-Obligation Decomposition}\label{ch:levels}
% =====================================================================

The first seven chapters use levels as a tool for comprehension
(Chapter \ref{ch:comprehension}), for design (Chapter \ref{ch:design}) and for assurance
(Chapter \ref{ch:intent}). This chapter shows that the three are three uses of one object,
and argues for the direction among them most worth pursuing: promoting levels from an aid to
comprehension to the \emph{unit into which verification obligations decompose}.

\section{The Cognitive Scale}\label{sec:levels-def}

\begin{definition}[Cognitive scale]\label{def:cognitive-scale}
Given a study object $S$ and a task $T$, the \term{cognitive scale} of $S$,
$\Lk{0}, \Lk{1}, \Lk{2}, \dots$, is divided by the type of question a level answers rather
than by the physical location of the code:
\begin{description}
\item[$\Lk{0}$ Computational principles] What performs the computation? What constrains it?
      Memory, instruction execution, function calls, concurrency, synchronisation, virtual
      memory, I/O, caching, networking, persistence, types and the memory model.
\item[$\Lk{1}$ Infrastructure and runtime] Which shared capabilities does the system
      provide, and how is their correct cooperation guaranteed? Allocators, task schedulers,
      event loops, object systems, executors, storage-engine interfaces, resource managers,
      lifecycle mechanisms. This is the level at which the six functions of
      Chapter \ref{ch:infrastructure} reside.
\item[$\Lk{2}$ Components and subsystems] Which duties exist? Where are the boundaries? How
      do they cooperate? The transaction subsystem, the rendering engine, the
      entity--component system, the network protocol stack, the replication module, the query
      executor.
\item[$\Lk{3}$ Domain function and business semantics] What does the system promise
      externally? Which behaviour is correct? Query semantics, game rules, layout rules,
      transaction isolation semantics, permissions and business invariants.
\item[$\Lk{4}$ and above: scenario and system-level behaviour] How does the whole system
      behave in real scenarios? End-to-end workflows, failure recovery, system-level
      performance, deployment topology, business objectives.
\item[$\Lk{\infty}$ Purpose and responsibility] What should be wanted? Which trade-offs are
      acceptable? This level belongs to humans; the reason is given in \S\ref{sec:linf}.
\end{description}
\end{definition}

\begin{principle}[The cognitive scale is not a physical hierarchy]\label{pr:not-physical}
The levels of Definition \ref{def:cognitive-scale} denote \emph{cognitive scales}, not
necessarily the physical layers of the code. The $\Lk{3}$ domain function of a game engine
may call the $\Lk{1}$ runtime interface directly; the transaction semantics of a database may
equally impose constraints on the storage layer.
\end{principle}

\begin{principle}[Relativity and recursiveness of levels]\label{pr:relative}
The same component may occupy different levels under different study tasks. An asynchronous
runtime is infrastructure ($\Lk{1}$) to an ordinary application developer, and to its own
developers a complex system requiring deep modification (it has a complete
$\Lk{0} \dots \Lk{4}$ of its own). Levels are therefore defined relative to the study object
and the current task.
\end{principle}

Principle \ref{pr:relative} has an important corollary: \emph{every large system contains a
complete level ladder internally}, and this embedded ladder is a distinct thing from the
external technology-stack ladder. This explains directly the conclusion of
Chapter \ref{ch:infrastructure}: every large system grows a ``small operating system'' inside
its own domain, and that is precisely its embedded $\Lk{1}$.

\begin{remark}[Difference from the causal hierarchy]\label{rmk:pearl}
The levels of this paper resemble the causal hierarchy \cite{pearl2009} in form: an
intervention at a lower level does not automatically transmit to a higher level, just as a
modification of the implementation layer does not automatically become a modification of the
contract layer. The two differ, however. The causal hierarchy distinguishes \emph{types of
query} (observation, intervention, counterfactual); the levels of this paper distinguish
\emph{observable domain and guarantee obligation} (Definition \ref{def:level}). Conflating
them yields a concrete error: the belief that being able to ask a question at a higher level
amounts to being able to intervene at that level, when the latter requires an interface the
former does not provide.
\end{remark}

\section{The Level Contract}\label{sec:four-tuple}

\begin{definition}[The level-contract quadruple]\label{def:level-contract}
A level is recognised only if one can write down the quadruple
\begin{equation}\label{eq:level-contract}
\mathrm{LC}(\Lambda) \;=\;
\bigl(\text{participants},\; \text{guaranteed invariants},\; \text{contract form},\;
\text{verification method}\bigr),
\end{equation}
together with the fifth item this paper additionally requires, the \term{oblivious set}
$\Obl(\Lambda)$ (Definition \ref{def:oblivious}), that is, what the level is permitted never
to care about.
\end{definition}

\begin{principle}[Criterion of level reality]\label{pr:level-real}
A level is real if and only if the five items of Definition \ref{def:level-contract} can be
written down. A ``level'' whose fifth item (the oblivious set) cannot be written down is not
a level but a directory or a namespace.
\end{principle}

Applying Definition \ref{def:level-contract} to the instances of
Table \ref{tab:level-instances} shows the concrete form taken by one and the same set of
questions at different levels.

\begin{table}[htbp]
\centering
\footnotesize
\caption{Instantiation of the level-contract quadruple on a relational database
engine}\label{tab:level-instances}
\begin{tabular}{L{1.2cm}L{2.4cm}L{3.6cm}L{3.4cm}L{3.6cm}}
\toprule
\textbf{Level} & \textbf{Participants} & \textbf{Guaranteed invariants} &
  \textbf{Contract form} & \textbf{May ignore} \\
\midrule
$\Lk{0}$ & pages, cache lines, instructions &
  the locality principle; atomicity of a cache line &
  hardware timing and the memory model &
  upper-level semantics; the notion of a transaction does not exist \\
$\Lk{1}$ & buffer pages, log records, locks, memory contexts &
  a dirty page is not reused before it is flushed; the log reaches disk before the data &
  the \code{palloc} contract, the buffer pin/release protocol &
  SQL semantics; query-plan shape \\
$\Lk{2}$ & parser, planner, executor, access methods &
  the tuples consumed by the executor satisfy the visibility of the current snapshot &
  the executor interface (open/next/close) &
  the physical layout of storage; the log format \\
$\Lk{3}$ & transactions, statements, sessions &
  the visibility and atomicity declared by the isolation level &
  SQL semantics and the definition of isolation levels &
  buffer-management policy; index structures \\
$\Lk{4}$ & applications, operations, clients &
  an end-to-end workflow remains recoverable under failure &
  service-level objectives and recovery-point objectives &
  the shape of the internal execution plan \\
\bottomrule
\end{tabular}
\end{table}

The most informative column of Table \ref{tab:level-instances} is the last. It makes explicit
the other half of ``what this level is responsible for'', namely ``what this level is
permitted not to know''. That column is the one most often missing from the design documents
of large systems, and the most valuable.

\section{Decomposition of Verification Obligations}\label{sec:obligation-unit}

This is the central claim of the chapter.

\begin{principle}[The three identities of a level]\label{pr:three-identities}
One and the same level structure carries three functions at once:
\begin{enumerate}[label=(\roman*)]
\item \term{unit of comprehension}: it fixes the lower bound on ``how much information is
      needed to understand a system'', namely the composite of the level quotients
      (Proposition \ref{prop:oblivious});
\item \term{unit of design}: it fixes the structure of ``the read-set of a typical change''
      (Theorem \ref{th:centralise});
\item \term{unit of verification obligation}: it fixes \emph{which properties can be assured
      mechanically, which must be monitored at runtime, and which must be left to the human
      division of responsibility}.
\end{enumerate}
Item (iii) is the part that existing discussions touch on but do not develop.
\end{principle}

A level is naturally a unit of verification-obligation decomposition because it supplies two
things at once: an \emph{observable domain} ($\Obs_i$, what the level can see) and an
\emph{invariant set} ($\Inv_i$, what the level must guarantee). The form of a verification
obligation is precisely to decide a property over some observation domain, so every
verification obligation attaches naturally to some level.

\begin{definition}[Level attribution of a verification obligation]\label{def:obligation-assign}
Given a property $\varphi$ and a level $\Lambda$, $\varphi$ is \term{attributed to} $\Lambda$
if $\varphi$ can be expressed entirely within the observable domain $\Obs_\Lambda$ of
$\Lambda$. The level of attribution of $\varphi$ is the \emph{lowest} level that can express
it.
\end{definition}

Three classes of obligation follow:

\begin{enumerate}[label=(\alph*)]
\item \term{Mechanically decidable at the level}: $\varphi$ is expressible within the
      observable domain of $\Lambda$ and admits an exhaustible decision procedure. Assigned
      to mechanical assurance.
\item \term{Monitorable at the level but undecidable}: $\varphi$ is expressible within the
      observable domain but can be decided only at execution time, over traces already
      observed. Assigned to runtime assurance, with the three states distinguished by
      Principle \ref{pr:three-states}.
\item \term{Not expressible at the level}: $\varphi$ requires cross-level information or
      concerns purpose itself. Assigned to human attestation or consensus assurance.
\end{enumerate}

\begin{table}[htbp]
\centering
\footnotesize
\caption{Distribution of assurance types across levels (example)}\label{tab:assurance-matrix}
\begin{tabular}{L{4.2cm}L{3.2cm}L{3.2cm}L{3.6cm}}
\toprule
\textbf{Property} & \textbf{Level of attribution} & \textbf{Assurance mode} &
  \textbf{Evidence type} \\
\midrule
memory safety & $\Lk{0}$--$\Lk{1}$ & mechanical (type system / ownership / sandboxing) &
  compile-time checking, verifier \\
no out-of-bounds access and bounded loops & $\Lk{0}$--$\Lk{1}$ &
  mechanical (bytecode verifier) & static proof at admission \\
no unauthorised I/O & $\Lk{1}$--$\Lk{2}$ & mechanical (capability model / effect types) &
  compile-time $+$ runtime boundary \\
functional postconditions & $\Lk{2}$--$\Lk{3}$ &
  mechanical (Hoare / refinement types / contracts) & proof terms, SMT decision \\
resource bounds & $\Lk{1}$--$\Lk{2}$ & mechanical (resource analysis) $+$
  environmental assumptions & symbolic bounds, worst-case execution time analysis \\
real-time deadlines & $\Lk{2}$--$\Lk{4}$ & mechanical (within the hardware model) $+$
  runtime & timing analysis, monitoring \\
isolation-level semantics & $\Lk{3}$ & mechanical (model checking / contracts) &
  model-checking results \\
end-to-end availability & $\Lk{4}$ & monitoring $+$ consensus &
  operational metrics, service-level objectives \\
the specification is faithful to the intent & $\Lk{\infty}$ & consensus &
  adversarial review, attestation records \\
\bottomrule
\end{tabular}
\end{table}

Table \ref{tab:assurance-matrix} concretises the claim of this chapter. For a property, three
questions are asked in order: to which level does it belong; which assurance can that level
provide; and what are the assumptions of that assurance. These three questions replace
questions that admit no answer, such as ``do we need formal verification''.

\section{The Bearer of Responsibility}\label{sec:linf}

\begin{principle}[Responsibility cannot be transferred]\label{pr:responsibility}
Deciding what is wanted and bearing the consequences of that decision are two facets of one
act of agency. Hence $\Lk{\infty}$ does not belong to machines, not because machines are not
intelligent enough, but because responsibility must be borne by a subject capable of bearing
consequences.
\end{principle}

This principle must be made precise by the results of \S\ref{sec:consensus}. The residual
responsibility of humans is \emph{not individual but collective}: not ``one person reviews
the specification'' but ``a cognitive network reaches consensus on the specification through
adversarial debate''. That network may include human experts of different backgrounds, models
of different training provenance, and explicitly designed adversarial roles.

\begin{principle}[Diversity is irreplaceable]\label{pr:human-diversity}
What makes humans irreplaceable in the age of AI is not that individual human judgement is
better than that of an AI, but that humans are an irreducible \emph{source of diversity} in
the cognitive network. By Corollary \ref{cor:diversity}, the reliability of consensus is
governed by the independence of error modes, and the ultimate source of independence is
\emph{different histories of problem formation}. The value of humans therefore lies in the
diversity of their perspectives, not in the precision of their computation.
\end{principle}

\section{Level-Language Mismatch}\label{sec:antipattern}

\begin{principle}[Consistency of level language]\label{pr:level-language}
The gravest error in reading code is to understand the current level in the language of the
wrong level.
\end{principle}

\begin{example}[Concrete forms]\label{ex:wrong-language}
\begin{enumerate}[label=(\alph*)]
\item Understanding business validation logic by pointer analysis, treating ``is this
      transaction permitted'' as a question of memory aliasing.
\item Understanding buffer management in business semantics, looking for commercial rules in
      a page-replacement policy.
\item Understanding database isolation levels in the language of scheduler fairness: both are
      concurrency problems, but the correctness criterion of the former is a throughput and
      latency distribution, and that of the latter is visibility and atomicity.
\item Solving multi-tenant isolation with a type system, because isolation is a cross-cutting
      concern (Principle \ref{pr:crosscutting}) and a type system can express only part of it.
\end{enumerate}
\end{example}

The antipattern of Example \ref{ex:wrong-language}(d) has a more serious form: it is not
merely a wrong language but an attempt to press a cross-cutting concern into the level
structure. By Principle \ref{pr:crosscutting}, such an attempt is structurally doomed to fail,
and its failure is often misdiagnosed as ``layering that is not thorough enough''.

The empirical content of this chapter is the reproducibility of
Definition \ref{def:obligation-assign}: whether different analysts give the same level
attribution to the same property, and whether the assignment of
Table \ref{tab:assurance-matrix} reduces the omission of properties that should be assured
mechanically but are left to manual review. With a falling omission rate as the metric,
the claim is testable.

% =====================================================================
\chapter{The Falsifiable Research Programme}\label{ch:program}
% =====================================================================

Every claim in this paper is no more than an interpretive framework unless it is converted
into testable form. This chapter performs that conversion, following the distinction
between hard core and protective belt in the methodology of research programmes
\cite{lakatos1970}: claims that are not load-bearing belong in the protective belt, not
reinforced into the hard core.

\section{Methodology: From Interpretive Framework to Testable Hypotheses}

\begin{principle}[Demotion procedure]\label{pr:demote}
For every claim, carry out the following operations: write down its \emph{observable
consequences}; write down what observation would \emph{refute} it; write down the
\emph{minimal experiment} that would yield that observation; and write down the
\emph{control condition} of that experiment. Any claim for which these four steps cannot be
completed is to be marked explicitly as an interpretive framework rather than an empirical
assertion.
\end{principle}

By Principle \ref{pr:demote}, the claims of this paper fall into three classes:

\begin{itemize}
\item \term{Proved}: propositions \ref{prop:oblivious}, \ref{prop:overconstrain},
      \ref{prop:stopping}, \ref{prop:leakpair}, \ref{prop:soundness},
      \ref{prop:active}, \ref{prop:asym-struct}, \ref{prop:operad-compose},
      \ref{prop:transport}, \ref{prop:observability}, \ref{prop:exponent}; theorems
      \ref{th:integration}, \ref{th:six}, \ref{th:centralise},
      \ref{th:condorcet}; proposition \ref{th:correlated}; corollaries
      \ref{cor:threshold}, \ref{cor:buildadopt}, \ref{cor:diversity},
      \ref{cor:three-failures}.
\item \term{Argued under explicit assumptions}: principles \ref{th:compression},
      \ref{th:necessity}, \ref{th:separation}, \ref{th:design}; propositions
      \ref{prop:no-godel}, \ref{prop:three-grounds}.
\item \term{Advanced as falsifiable hypotheses}: the hypotheses $H_1$--$H_{14}$ below. In
      addition, the structural proposal of Remark \ref{rmk:fibration-status} (promoting the
      level family to a strict fibration) and the renormalisation analogy of Remark
      \ref{rmk:rg} are not counted among the theorems; their status is that of a proposal
      and a heuristic.
\end{itemize}

\section{Fourteen Falsifiable Hypotheses}

\begin{hypothesis}[Structural convergence above a random baseline]\label{hyp:h1}
After controlling for system size, domain, platform, programming language and framework,
large software systems from different domains exhibit structural similarity on their
abstract skeleton (Definition \ref{def:skeleton}) significantly above a random baseline.
\emph{Operationalisation}: construct a collection of architecture graphs and skeleton maps
by the five-step design of \S\ref{sec:baseline}, take the structural metrics of the third
step as the feature vector, and test whether the similarity distribution of cross-domain
pairs stochastically exceeds a null model, such as degree-sequence-preserving graph
rewiring.
\emph{Falsification condition}: the similarity distribution does not differ significantly
from the null model.
\emph{Risk}: freedom in the choice of skeleton map; the map and the metrics must be
pre-registered.
\end{hypothesis}

\begin{hypothesis}[Convergence strength monotone in domain binding]\label{hyp:h2}
Conjecture \ref{conj:strength} holds: structural similarity decreases monotonically in the
domain binding $\mathrm{Dom}$ (Definition \ref{def:dom}).
\emph{Operationalisation}: annotate $\mathrm{Dom}$ at each level (estimated from the
fraction of domain axioms on which that level's semantics depends) and regress similarity
on $\mathrm{Dom}$.
\emph{Falsification condition}: the monotonicity fails, or the coefficient is not
significant.
\end{hypothesis}

\begin{hypothesis}[Phase-transition threshold determined by $n \times m$]\label{hyp:h3}
Whether a system builds a unified substrate is predicted better by the product of the
number of live entities and the number of shared resources (together with the frequency of
recurrent problems and the frequency of evolution) than by lines of code.
\emph{Operationalisation}: for a set of open-source projects, annotate manually or
semi-automatically whether a self-built implementation of any of the six functions appears;
fit separate logistic regressions on a proxy for $n \times m$ (order of magnitude of
concurrent entities, number of kinds of shared resource) and on lines of code; compare
their explanatory power.
\emph{Falsification condition}: lines of code explain the outcome no worse than
$n \times m$.
\emph{Risk}: automatic estimation of $n$ and $m$ is difficult; the number of concurrent
entities and of mutexes, connection pools, caches and the like can serve as a first proxy.
\end{hypothesis}

\begin{hypothesis}[Read-set scaling outperforms traditional size metrics]\label{hyp:h4}
The change read-set $\mathbb{E}[\rho]$ (Definition \ref{def:readset}) and its \emph{scaling
exponent} $\alpha$ with respect to system size (Proposition \ref{prop:exponent}) correlate
with defect rate and change failure rate better than traditional size metrics such as
module count, lines of code and cyclomatic complexity, and better than the current magnitude
of the read-set $\mathbb{E}[\rho]$ itself.
\emph{Operationalisation}: on several open-source projects with a complete commit history
and defect tracking, estimate the magnitude and the scaling of the read-set over the same
time windows, regress each in turn on defect density in a panel design, and compare their
explanatory power.
\emph{Falsification condition}: the scaling provides no predictive power beyond the
magnitude and the existing metrics.
\end{hypothesis}

\begin{hypothesis}[Layered model plus active probing outperforms exhaustive reading]\label{hyp:h5}
Under a constraint on unit cost (token budget or person-hours), a comprehension process
that uses a layered cognitive model plus active probing attains a higher task success rate
(scored against (U1)--(U5) of Definition \ref{def:understand}) than exhaustive reading.
\emph{Operationalisation}: given the same unfamiliar code base and the same set of tasks
(predict the impact surface, locate a defect, add a feature), compare three strategies:
sequential reading, layered summaries only, and layered summaries plus a probing loop.
Measure task success rate and success rate per unit cost.
\emph{Falsification condition}: the active-probing strategy is not significantly better, or
its advantage comes entirely from the number of probes rather than from the structure of the
strategy.
\end{hypothesis}

\begin{hypothesis}[Leak pairs are automatically discoverable]\label{hyp:h6}
The leak pairs of Proposition \ref{prop:leakpair} can be found by automated differential
testing, and the discovery rate is negatively correlated with module boundary quality.
\emph{Operationalisation}: for several libraries, generate execution pairs that are
``identical at the interface and different in internal state'', detect divergence in
external behaviour, and compare against human-annotated boundary quality scores.
\emph{Falsification condition}: automatically generated execution pairs fail to trigger
leaks effectively, or the discovery rate is unrelated to boundary quality.
\end{hypothesis}

\begin{hypothesis}[Specification quality is the present bottleneck]\label{hyp:h7}
At the present state of the art, the principal bottleneck on the path from intent to
verifiable execution is specification quality, not code generation capability or proof
search capability.
\emph{Operationalisation}: take existing benchmark evidence (specification generation
tasks, evaluations that separate code correctness from specification quality, proof success
rates) as initial evidence, and on a new controlled benchmark compare the marginal
contribution of each of the three links to end-to-end success rate.
\emph{Initial evidence}: a component-wise evaluation on 189 Lean tasks reports, for the
best model, code correctness 72.6\%, specification soundness and completeness 52.3\%, and
proof success rate 4.9\% (single trial) \cite{verina2026}; an evaluation of 245
operating-system kernel specification generation tasks shows that models remain deficient
in long context and specification completeness \cite{osvbench2026}; an evaluation that
instantiates specifications as trusted tests points out that proof difficulty interferes
with the assessment of specification quality \cite{coins2026}; a traceable refinement
method reports accuracy gains on the specification-generation subtask
\cite{verispecgen2026}; a mutation-driven method reports precision and recall
\cite{specsyn2026}.
These are empirical results on specific benchmarks, not general reliability theorems, but
they support the judgement that the three dimensions should be evaluated separately.
\emph{Falsification condition}: with the capability of the other links controlled, the
marginal contribution of specification quality is not significantly greater than that of
the other links.
\end{hypothesis}

\begin{hypothesis}[Consensus reliability is dominated by correlated error modes]\label{hyp:h8}
On intent-attestation tasks, the reliability gain from adding verifiers is bounded above by
the correlated error mode $q$ (Proposition \ref{th:correlated}); introducing heterogeneous
verifiers (people of different backgrounds, models from different sources, adversarial
roles) improves reliability more than merely adding homogeneous verifiers.
\emph{Operationalisation}: construct intent-attestation tasks, configure several verifier
groups (homogeneous or heterogeneous), measure how the error rate of the majority judgement
varies with group size and heterogeneity, and compare with the predictions of Equations
\eqref{eq:chernoff} and \eqref{eq:correlated}.
\emph{Falsification condition}: the rate at which the error rate of homogeneous verifiers
decreases with $n$ does not differ significantly from that of heterogeneous verifiers (that
is, $q$ is in fact small).
\end{hypothesis}

\begin{hypothesis}[The artefact binding gap is a measurable risk]\label{hyp:h9}
The binding gap (Equation \eqref{eq:binding}) between the verified artefact and the
artefact actually executed is a measurable risk in real deployments and can be
significantly reduced, and reducing it yields an observable reliability benefit.
\emph{Operationalisation}: introduce content addressing, reproducible builds and
deploy-time verification into several deployment pipelines, and measure the change in the
proportion of build nondeterminism, the rate of deployment drift events and the rate of
related incidents.
\emph{Falsification condition}: the gap is not measurable in practice, or the cost of
closing it far exceeds the benefit.
\end{hypothesis}

\begin{hypothesis}[The six functions are the surviving set under coarse-graining]\label{hyp:h10}
Under the coarse-graining sequence of Definition \ref{def:coarsegrain}, the structures that
survive are exactly the slots of $\mathcal{K}_6$ and their relational graph (Conjecture
\ref{conj:relevant}).
\emph{Operationalisation}: apply mechanical coarse-graining to several repositories
following Principle \ref{pr:coarsegrain-impl}, and record how the surviving set changes
with granularity.
\emph{Falsification condition}: the surviving set contains structures outside the six
functions, or some of the six functions are washed out already at intermediate granularity.
\emph{Risk}: the choice of coarse-graining threshold introduces an arbitrary parameter; the
threshold should be subjected to a sensitivity analysis.
\end{hypothesis}

\begin{hypothesis}[Incentive structure reduces correlated error]\label{hyp:h11}
On intent-attestation tasks, adopting an appropriate scoring rule (Principle
\ref{pr:scoring}) reduces the error rate of the majority judgement more than merely
increasing the number of verifiers, and a measurable fall in error correlation accompanies
it.
\emph{Operationalisation}: on the same set of specification-attestation tasks, compare
three configurations: few unincentivised verifiers, many unincentivised verifiers, and few
verifiers under a scoring rule; measure the error rate and the correlation coefficient of
errors between verifiers.
\emph{Falsification condition}: neither the error rate nor the error correlation
coefficient of the scoring-rule configuration is better than that of the
many-unincentivised configuration.
\end{hypothesis}

\begin{hypothesis}[Rate--distortion slope predicts refactoring benefit]\label{hyp:h12}
The slope $\partial L_T / \partial \DL$ of the rate--distortion curve (Remark
\ref{rmk:rd-refactor}) predicts which refactorings reduce defect rate or rework rate in
practice.
\emph{Operationalisation}: for a set of completed refactorings, estimate ex ante the
direction of displacement along the curve and the slope, and ex post regress the changes in
defect rate and change failure rate on them.
\emph{Falsification condition}: the slope provides no predictive power beyond naive
variables such as the size of the change.
\end{hypothesis}

\begin{hypothesis}[The trajectory of proof cost per unit of code]\label{hyp:h13}
Chapter \ref{ch:objections} declines to claim that formal verification becomes cheap automatically
and asks instead for a measurable cost metric. The hypothesis is that the ratio
\[
\kappa_{\text{proof}} \;=\; \frac{\text{effort to establish a stated guarantee}}{\text{size of the code carrying it}}
\]
is on a measurable trajectory, and that the trajectory---not its current value---decides whether
mechanised verification becomes routine engineering or remains confined to safety-critical work.
\emph{Initial evidence}: the seL4 functional-correctness proof is reported at about eleven
person-years for about 8{,}500 lines of C, or roughly $1.3$ person-years per thousand lines,
\emph{excluding} investment in reusable tools and frameworks \cite{klein2009}; figures that include
that investment are several times larger. The two figures differ by a factor worth stating, because
the reusable-tools share is precisely the part that should fall if the hypothesis holds.
\emph{Operationalisation}: for a fixed class of properties and a fixed prover, record
$\kappa_{\text{proof}}$ per project and per year, separating the reusable-infrastructure share from
the per-project share; fit the trajectory over the observation window.
\emph{Falsification condition}: no measurable trend, or a trend attributable entirely to the
reusable-infrastructure share while the per-project share is flat. Since the paper declines to
predict which way this goes, both outcomes are informative.
\emph{Risk}: the denominator is contentious---lines of code, of specification, or of proof---and the
choice must be pre-registered or the trajectory is not comparable across projects.
\emph{Interpretation bands}: the thresholds conventionally proposed for reading the trajectory, such
as $0.1$ for ``routine'' and $0.5$--$1$ for ``specialist'', are stipulated decision bands rather
than measured quantities, and are labelled as such here.
\end{hypothesis}

\begin{hypothesis}[Agentic generation can lower the read-set scaling exponent]\label{hyp:h14}
Proposition \ref{prop:exponent} makes the read-set scaling exponent $\alpha$ the quantity that
decides whether a system's governance cost grows linearly or quadratically in its size. Every
hypothesis above concerns systems built by people or by human-directed tools. This one asks the
AI-era question directly: can an agentic generator, given licence to restructure, \emph{discover}
architectures with a lower $\alpha$ than the human baseline, or does it only move constants within
an architecture it inherits? The distinction is not cosmetic---it is the difference between
automation that accelerates work and automation that changes what remains possible at scale.
\emph{Operationalisation}: start an agent and a human control group from the same code base with the
same feature backlog and the same freedom to refactor; track $\mathbb{E}[\rho]$ and the fitted
$\alpha$ against system size over the same window; compare trajectories.
\emph{Falsification condition}: $\alpha$ does not differ significantly between arms, or the apparent
difference is explained entirely by the volume of code written rather than by its structure.
\emph{Risk}: $\alpha$ is estimated through the read-set proxy of Definition \ref{def:readset}, whose
validity is itself hypothesis $H_4$; a failure of $H_4$'s proxy would make this hypothesis
unfalsifiable rather than false, and the two must not be confused.
\end{hypothesis}

\begin{table}[htbp]
\centering
\footnotesize
\caption{Summary of hypotheses and experimental designs}\label{tab:hypotheses}
\begin{tabular}{L{0.9cm}L{4.0cm}L{4.6cm}L{4.6cm}}
\toprule
\textbf{ID} & \textbf{Hypothesis} & \textbf{Principal measurement} & \textbf{Principal risk} \\
\midrule
$H_1$ & Structural convergence above a random baseline & Skeleton similarity vs.\ null
  model & Excessive freedom in the choice of skeleton map \\
$H_2$ & Convergence strength monotone in domain binding &
  Regression of similarity on $\mathrm{Dom}$ &
  $\mathrm{Dom}$ is hard to annotate independently \\
$H_3$ & Threshold determined by $n\times m$ & Comparison of logistic-regression explanatory
  power & Automatic estimation of $n$ and $m$ \\
$H_4$ & Read-set scaling outperforms traditional size metrics &
  Panel regression of read-set scaling on defect density & Accuracy of the read-set proxy \\
$H_5$ & Layered model plus active probing outperforms exhaustive reading &
  Task success rate per unit cost & Subjectivity of tasks and scoring \\
$H_6$ & Leak pairs are automatically discoverable & Differential-testing discovery rate &
  Difficulty of triggering valid execution pairs \\
$H_7$ & Specification quality is the bottleneck & Marginal contribution of the three links &
  Ground truth and coverage of the benchmark \\
$H_8$ & Consensus dominated by $q$ & Error rates of homogeneous and heterogeneous verifiers &
  Construction and measurement of heterogeneity \\
$H_9$ & Artefact binding gap is measurable & Build nondeterminism and drift event rate &
  Availability of real deployments \\
$H_{10}$ & The six functions are the coarse-graining surviving set &
  Surviving set under the coarse-graining sequence &
  Arbitrariness of the coarse-graining threshold \\
$H_{11}$ & Incentive structure reduces correlated error &
  Error rate and error correlation coefficient &
  Experimental ethics and cost of incentive schemes \\
$H_{12}$ & Rate--distortion slope predicts refactoring benefit &
  Regression of defect rate before and after refactoring & Estimability of the slope \\
$H_{13}$ & Proof cost per unit of code follows a measurable trajectory &
  Fitted trajectory of $\kappa_{\text{proof}}$, tools share separated &
  Contentious denominator; comparability across projects \\
$H_{14}$ & Agentic generation can lower the read-set exponent $\alpha$ &
  Trajectory of fitted $\alpha$ against system size, two arms &
  Depends on the read-set proxy of $H_4$ \\
\bottomrule
\end{tabular}
\end{table}

\section{Three Priority Research Directions}

The three directions below connect requirements engineering with formal methods,
compilation correctness with resource analysis, and program contracts with runtime
security. They share one characteristic: each has a precisely stated question, a
falsifiable hypothesis, a method of testing and a principal risk.

\subsection*{Direction One: Specification-Quality Evidence Chain for Real Intent}

\paragraph{The precise question.}
Given a set $R$ of natural-language requirements, a domain model $D$ and a generated
specification $S$, how is an independent evaluation process to be built that checks
separately: requirement coverage; exclusion of forbidden behaviour; retention of permitted
behaviour; explicitness of environmental assumptions; and the satisfiability and
implementability of the specification?

\paragraph{Existing foundations.}
Requirements traceability research, requirement formalisation and verification, vacuity
detection, and recent specification-generation and joint-evaluation benchmarks. The key
connection is \emph{semantic traceability} in requirements engineering and
\emph{specification-versus-program counterexamples} in formal methods.

\paragraph{Falsifiable hypothesis.}
Compared with a single ``proof passed'' signal or a single ``similarity to a reference
specification'' score, atomic requirement traceability plus positive and negative testing
plus specification-strength and mutation checks plus independent domain review identifies
the \emph{under-constraint}, \emph{over-constraint} and \emph{intent omission} of a
specification earlier, that is, it is the concretisation of Hypothesis $H_7$.

\paragraph{Method of testing.}
On small benchmarks such as payments, permission control and idempotent interfaces,
establish requirement atoms annotated by multiple stakeholders, permitted and forbidden
trace sets, and environmental preconditions; compare four approaches: pure generation,
requirement decomposition, traceability testing, and mutation refinement with expert
assistance. Measure precision and recall, error categories, the rate of missed
prohibitions, the over-constraint rate, expert agreement and cost.

\paragraph{Falsification condition.}
If adding the traceability and mutation layers does not significantly reduce the rate at
which hidden negative requirements are missed, or if the gain comes only from leaked
reference test answers, the hypothesis fails.

\paragraph{Principal risk.}
The ``gold standard'' for real intent may itself be contested.

\subsection*{Direction Two: Joint Functional--Resource Translation Certificates for Actual Artefacts}

\paragraph{The precise question.}
Given a source program $P$, a functional contract $\varphi$, a resource budget $B$ and a
fixed target platform model $H$, can a certificate $C$ be generated for \emph{each
compilation} such that an independent verifier checks semantic preservation and analysable
resource constraints at the same time?

\paragraph{Existing foundations.}
Semantic-preservation proofs for verified compilers, per-translation validation,
translation certification, proof-carrying code, automatic amortised resource analysis and
worst-case execution time analysis. The substantive connection is to place the
``behaviour relation'' and the ``cost relation'', currently verified \emph{separately},
inside a \emph{cost-carrying behaviour semantics}.

\paragraph{Working model.}
The program semantics is extended to a pair $(\tau, c)$, where $\tau$ is an observable
trace and $c$ a cost; the translation proof should guarantee that target traces are
explicable by the source semantics, and should prove $c_{\text{target}} \le B$ under an
explicit cost model, or else give a conservative relation relative to the source cost.

\begin{remark}[An assumption that must be avoided]
It must not be presupposed that optimised compilation always lowers cost monotonically.
Optimisation may improve average performance while worsening the worst case; the cost
relation must therefore be a conservative relation comparable in both directions, not an
assumption of monotonicity.
\end{remark}

\paragraph{Method of testing.}
Choose a reproducible sublanguage and target architecture, and compare three routes:
verified compilation, translation validation, and ordinary compilation with post-hoc
verification. Record certificate size, verification time, coverage of program categories,
tightness of resource bounds, false-rejection rate, and robustness to changes in compiler
optimisation.

\paragraph{Falsification condition.}
If the resource proofs do not remain sound once calibrated against a real platform, or if
certificate generation and verification cost far more than the baseline with no safety
benefit, the scope should be narrowed to specific safety-critical domains.

\paragraph{Principal risk.}
The target processor microarchitecture and scheduling model are not precise enough.

\subsection*{Direction Three: Composition Contracts and Runtime Assurance for Untrusted Components}

\paragraph{The precise question.}
When a service is composed of verified components, generated components that are not fully
verified, and external services, how are assumption--guarantee contracts, information-flow
policies, resource budgets and runtime monitoring to be used to give an auditable overall
safety proposition together with an upper bound on failure impact?

\paragraph{Existing foundations.}
Modular assumption--guarantee reasoning, contract composition for dynamical systems,
information-flow non-interference, verified kernels, and studies of uncertainty in runtime
verification. The substantive connection is to place the formal contracts of subcomponents
and the runtime detection and enforcement of contract premises and guarantees within a
single system argument.

\paragraph{Falsifiable hypothesis.}
When every component has explicit assumptions, entry and exit events are observable, and
contract violations can trigger isolation or rollback from an independent monitor, a
composable conditional guarantee can be given for a specified safety property without
proving the internal implementation of each component correct.

\paragraph{Method of testing.}
Build a multi-component payment flow: an idempotent request service, an identity and
permission service, a generative risk-decision component, an external database and a
network simulator. Under fault injection, compare three groups (no contracts, static
contracts, contracts plus monitoring and rollback) and measure duplicate charges,
cross-user data leakage, detection latency, isolation success rate and performance cost.

\paragraph{Falsification condition.}
If the events observable by the monitor are insufficient to decide the critical properties,
or if the rollback mechanism cannot itself remain safe, the architecture cannot provide the
guarantee it claims.

\paragraph{Principal risk.}
For subtle semantic properties that go unobserved, runtime assurance can constrain only the
detectable and enforceable part.

\section{Six-Stage Implementation Roadmap}

\begin{table}[htbp]
\centering
\small
\caption{Research implementation roadmap}\label{tab:roadmap}
\begin{tabular}{L{1.0cm}L{6.3cm}L{7.0cm}}
\toprule
\textbf{Stage} & \textbf{Core tasks} & \textbf{Deliverables} \\
\midrule
1 & Define the benchmark and the threat model: choose a payment service and an idempotent
  interface as the running case; fix the assets, roles, observable events, forbidden
  behaviour, environmental assumptions and time limits &
  Set of requirement atoms, positive and negative trace sets, list of environmental
  assumptions (attested by stakeholders) \\
2 & Specification quality assessment: compare pure generation, traceability refinement, and
  mutation- and counterexample-driven approaches; check for under-constrained
  specifications, over-constrained specifications, vacuous satisfaction, unsatisfiability
  and unimplementability &
  Specification quality assessor; for each requirement identifier, the corresponding
  specification clause, test and attestation record \\
3 & Program synthesis in a restricted language: choose an explicit sublanguage and a logic
  of controlled decidability; separate candidate generation from proof checking &
  Candidate programs, proof obligations, counterexample classification; distinguish
  unsatisfiable, unknown, timeout and inexpressible \\
4 & Artefact and resource proofs: choose a mature compilation or translation-validation
  route; specify reproducible experimental hardware or a simulation model &
  Compilation-chain metadata, semantic-preservation or translation certificates, resource
  bounds and model boundaries \\
5 & Composition and runtime assurance: write assumption--guarantee contracts for external
  services and not-fully-verified modules; add independent monitoring, sandboxing, auditing
  and safe rollback &
  System-level contract composition results, fault-injection report, residual risk list \\
6 & End-to-end evaluation: measure assurance coverage, false negatives, false positives,
  development cost, proof cost, artefact verification cost, resource budgets and failure
  recovery &
  A public benchmark capable of falsifying the research hypotheses, a tool prototype and an
  experimental report \\
\bottomrule
\end{tabular}
\end{table}

\section{Success Criteria}

The criteria below assess whether the roadmap succeeds. They are designed to be
independently checkable:

\begin{enumerate}[label=(S\arabic*)]
\item \term{Explicit object}: the object of each guarantee (which artefact, which
      semantics, which environment) can be uniquely identified.
\item \term{Real assumptions}: the environmental assumptions relied upon can be tested, not
      assumed by default.
\item \term{Independent evidence}: the judgement of evidence does not depend on the
      generator's self-assessment.
\item \term{Distinguishable failure categories}: unsatisfiable, unknown, timeout,
      inexpressible, and candidate refuted by counterexample must be reported separately and
      not merged into a single ``failure''.
\item \term{Measurable cost}: generation cost, proof cost, verification cost and runtime
      cost can be measured separately.
\item \term{Falsifiability}: the refutation of at least one $H_i$ would weaken the core
      claims of the programme.
\end{enumerate}

\begin{remark}[Self-examination of the programme]
One core claim of this paper is that ``generation is untrusted and verification must be
independent''. The principle applies to the paper itself: its load-bearing propositions
must be independently checkable and must not depend on any unproved conjecture. Remark
\ref{rmk:q1q8} and the list of non-claims in Chapter \ref{ch:objections} are the concrete
results of that self-examination.
\end{remark}

\section{Minimum Viable Subset and the Adoption Risk}\label{sec:mvs}

The hypotheses above share a failure mode that none of them covers. They are each refuted or
confirmed by experiment, and the framework can also simply fail to be adopted, in which case no
experiment is run at all. The risk is worth stating explicitly because it is orthogonal to
falsifiability, and because it has a specific shape: the formal objects cost something to
maintain---declarations must be written, kept current and defended---and if that cost exceeds the
benefit, the practice dies of neglect rather than of refutation.

\begin{principle}[Minimum viable subset]\label{pr:mvs}
If only three elements of the framework can be adopted, they should be:
\begin{enumerate}[label=(\alph*)]
\item \term{the oblivious-set declaration}: each level or module states what it is permitted never to
      care about (Definition \ref{def:oblivious});
\item \term{the leak-pair test}: for each claimed boundary, a differential test seeking two executions
      with equal interface traces and different external behaviour (Proposition \ref{prop:leakpair},
      Corollary \ref{cor:diagnose});
\item \term{the level attribution of verification obligations}: each property is assigned to the
      lowest level that can express it, which decides whether it is mechanically decidable,
      monitorable, or a matter for attestation (Definition \ref{def:obligation-assign}).
\end{enumerate}
These three form the smallest closed loop in the framework: a declaration, a test of the declaration,
and a routing decision that consumes both. The contract operad of \S\ref{sec:operad} and the consensus
bound of \S\ref{sec:consensus} are refinement layers---valuable where the corresponding difficulty is
present, dispensable where it is not. Principles \ref{pr:descent} and \ref{pr:two-criteria} are what
make the loop usable: the first says how deep to look, the second says what the declaration has to
satisfy.
\end{principle}

The measurable form of the adoption claim is deliberately modest, and it is the metric Chapter
\ref{ch:levels} already proposes for its own empirical content: the rate at which properties that
should be assured mechanically are instead left to manual review, and the rate at which claimed
boundaries fail their own tests. If adopting the three elements above does not move those rates in a
real project, then the framework's problem is not that a hypothesis was false but that the remedy
costs more than the disease.

Stating the adoption risk as a fifteenth hypothesis would also be defensible, with adoption rate and
cost-to-benefit ratio as its measurements. This paper does not, for a reason that is itself part of
the framework's discipline: the falsification of a hypothesis should turn on an observation about the
world, whereas non-adoption turns on a decision by practitioners, and letting one word cover both
would blur the distinction between verification and validation that Chapter \ref{ch:intent} works to
keep sharp. The choice is recorded here rather than made silently.

% =====================================================================
\chapter{Literature Map and Positioning}\label{ch:literature}
% =====================================================================

This chapter places the paper's claims in their literature context, states their relation
to existing theory, and marks where the gaps lie.

\section{Six Groups of Literature}\label{sec:six-groups}

\paragraph{First Group: Modern Specification Generation and Joint Evaluation.}
This is the frontier closest to the intent $\to$ specification step. The relevant
literature reports measurable results: a specification-generation benchmark for
operating-system kernel verification tasks still shows model limits on long context and
specification completeness \cite{osvbench2026}; a benchmark that evaluates code
correctness, specification quality and proof generation separately reports a significant
gap among the three dimensions (on a particular task set, code correctness, specification
soundness and completeness, and single-pass proof success rate are 72.6\%, 52.3\% and
4.9\% respectively) \cite{verina2026}; an evaluation framework that instantiates
specifications as trusted tests points out that proof difficulty interferes with the
assessment of specification quality \cite{coins2026}; a traceable refinement method
decomposes natural-language requirements into atomic requirements and establishes
requirement-level traceability \cite{verispecgen2026}; a mutation-driven method tests
specification strength with semantically non-equivalent variants \cite{specsyn2026}; and
the route of generating specifications from code should be studied separately from the
route of generating specifications from intent \cite{specgen2025}. Vacuity detection
\cite{kupferman2003} supplies a theoretical route to identifying the case in which
verification passes while the specification holds vacuously. Together these results support
one important judgement: \emph{code generation ability, specification expressiveness and
formal proof ability are distinct research variables and should be evaluated separately}.

\paragraph{Second Group: Program Semantics, Refinement and Observable Behaviour.}
Operational semantics, denotational semantics, the lattice model of abstract interpretation
\cite{cousot1977}, monadic semantics \cite{moggi1991,plotkin2002}, domain theory,
refinement relations \cite{hoare1969} and separation logic \cite{appel2007} together supply
the mathematical tools for the question of when two systems at different levels implement
the same behavioural specification. Semantic refinement is not simple information
compression: it must define the \emph{behavioural relation} between levels.

\paragraph{Third Group: Type Theory, Constructive Mathematics and Proof Extraction.}
Constructive type theory \cite{hott2013}, the correspondence between proofs and programs
\cite{howard1980}, and the mechanisms that connect proof objects with constructive
computation supply the mathematical foundation for ``an AI proposes candidate proofs and a
mechanical kernel checks them''. One inference must be avoided: \emph{if a specification
has already been translated into a type, then the specification necessarily expresses human
intent faithfully}. The correspondence guarantees only the stipulated structural
correspondence among propositions, types, proofs and terms within a chosen logic and type
system; how human intent is turned into a proposition is a separate process that requires
verification. Refinement types \cite{rondon2008} and SMT solving \cite{z32008} are the
engineering representatives of this direction; automated theorem-proving environments and
retrieval-augmented proof search \cite{leandojo2023,leancopilot2024} are its continuation
in the age of AI.

\paragraph{Fourth Group: From Proof to Machine Code.}
This is the most mature group. A verified compiler turns the compilation transformation
itself into a proved object \cite{leroy2009}; the back ends of a verified compiler cover
several target architectures \cite{cakeml2019}; a refinement proof for an operating-system
kernel shows how far the formal verification of large-scale low-level system software can
reach \cite{klein2009}; translation validation checks a \emph{single} compilation result
independently \cite{necula2000}; translation certification establishes a verifiable
relational certificate for a \emph{specific artefact} \cite{krijnen2024}; and proof-carrying
code lets a code provider attach a proof that a machine can check independently
\cite{necula1997}.

The key lesson of this group is that \emph{verifying the correctness of a compiler and
verifying the correctness of a particular binary are two different strategies}. The former
proves that the compiler satisfies a universal property; the latter verifies one build
output independently. The two can be combined, without staking all trust on a single tool.

\paragraph{Fifth Group: Resource Correctness.}
Automatic amortised resource analysis can derive symbolic resource bounds for a program
under a particular language and cost model, with a soundness proof; worst-case execution
time analysis derives timing bounds for particular hardware and assumptions; static cache
analysis shows the influence of the hardware model on bound inference. Three notions of
time must be distinguished: algorithmic complexity gives a bound on the growth of cost as
input size varies; worst-case execution time analysis gives an upper bound under a
particular platform and assumptions; measurement provides latency data under particular
experimental conditions. \emph{None of the three substitutes for another}. Equally, linear
use is not optimal use: a program may consume each input resource exactly once and still use
a poor algorithm or introduce excessive latency.

\paragraph{Sixth Group: Runtime Verification, Contract Composition and the Emergence of Levels.}
This group contains two classes of material of different kinds, which must be treated
separately.

The first class is \emph{runtime assurance and contracts}: assume--guarantee reasoning
requires a subsystem to give the corresponding guarantee when its input assumption is
satisfied \cite{henzinger1998}; contract composition for dynamical systems gives composition
rules and checking methods under arbitrary interconnection topologies
\cite{sharf2021,sharf2024,shali2026}; work on the uncertainty of runtime verification shows
that observation, event transport and trace integrity are components of monitoring
reliability \cite{taleb2023}; online monitoring of temporal logic \cite{leucker2009} and
information-flow security \cite{goguen1982} address trace properties and access
restrictions respectively.

The second class is \emph{explanatory frameworks}: the near-decomposability of complex
systems provides an explanation for level structures \cite{simon1962}; minimum description
length supplies the statistical and information-theoretic background for ``abstraction is
compression'' \cite{rissanen1978,grunwald2006,grunwald2019}; computational thermodynamics
gives a physical lower bound on computation \cite{landauer1961}. This group can supply an
explanatory framework for architecture, compression, model selection and complex systems,
but it \emph{cannot automatically supply a correctness proof for a particular program}.

\section{The Five Layers Rearranged by Mathematical Function}\label{sec:five-layers}

Rearranged by which segment of the chain it solves, this literature yields a structure
deeper than ``an AI writes code and a verifier checks it'':

\begin{table}[htbp]
\centering
\small
\caption{The five layers rearranged by mathematical function}\label{tab:five-layers}
\begin{tabular}{L{1.6cm}L{5.4cm}L{7.2cm}}
\toprule
\textbf{Layer} & \textbf{Function} & \textbf{Representative lines of work} \\
\midrule
A & Expressibility of the objective: formal specification generation, intent traceability,
  requirement decomposition, counterexample generation &
  Specification-generation benchmarks, traceable refinement, mutation-driven strength
  testing, vacuity detection, requirement formalisation \\
B & Realisability of the objective: logical consistency, program synthesis, constraint
  solving, realisability decision &
  Counterexample-guided inductive synthesis, syntax-guided synthesis, program-skeleton
  completion, constructive type theory, SMT solving \\
C & Cross-layer semantic preservation: operational semantics, refinement relations, proof
  checking, compiler correctness, binary verification &
  Verified compilers, verified compiler back ends, kernel refinement proofs, proof-carrying
  code, translation validation \\
D & Resource and environment assurance: complexity, resource bounds, runtime monitoring,
  assume--guarantee contracts, execution constraints &
  Automatic amortised resource analysis, worst-case execution time analysis, runtime
  verification, contract composition, information-flow security \\
E & Structure and physical limits: model compression, levels of complex systems,
  thermodynamics and physical realisation &
  Minimum description length, near-decomposability of complex systems, computational
  thermodynamics, molecular self-assembly \\
\bottomrule
\end{tabular}
\end{table}

The most important finding is that these areas \emph{do not share one ready-made, fully
unified mathematical object}. Each has nonetheless solved a different problem along this
chain. The research opportunity therefore lies in defining the composition relations among
them, not in inventing the entire mathematical toolkit from scratch.

\section{The Unoccupied Intersection}\label{sec:gaps}

\begin{principle}[Shape of the literature map]\label{pr:map-shape}
Every step from specification to binary has a mature literature; automated synthesis from
proof to implementation has an active frontier; but the gap from \emph{intent} to
\emph{specification} has no formal solution, only consensus structures, and the literatures
of social epistemology, Byzantine fault tolerance and adversarial review exist in entirely
different disciplines and have not been systematically coupled with formal verification
systems.
\end{principle}

On these grounds, the most valuable position to occupy is the following.

\begin{principle}[Proposed research position]\label{pr:position}
Integrate adversarial consensus in the style of social epistemology with mechanical
admission in the style of proof-carrying code inside one framework: let a
consensus-attested specification be the direct input to a proof-carrying artefact, and
anchor consensus records, specification versions and artefact identity together
cryptographically. As far as the review conducted for this paper extends, few works address
both literatures at once.
\end{principle}

This position is valuable because it satisfies three conditions at once: it corresponds to a
real and still unsolved technical gap (no mechanical solution at the intent layer); both of
its components already have mature theoretical foundations; and its composition rules can
be written out explicitly (Definition \ref{def:two-layer} and Theorem \ref{th:assurance})
and can therefore be tested.

\section{Relation to Existing Theory}\label{sec:relations-lit}

\begin{longtable}{L{3.6cm}L{10.8cm}}
\caption{Relation of this paper to existing theory}\label{tab:relations}\\
\toprule
\textbf{Theory} & \textbf{How this paper uses it, and how it does not use it} \\
\midrule
\endfirsthead
\toprule
\textbf{Theory} & \textbf{How this paper uses it, and how it does not use it} \\
\midrule
\endhead
\bottomrule
\endlastfoot
Near-decomposability of complex systems &
  \emph{Used}: as the explanation of why level structures are stable, and recast into a
  testable sufficient-statistic condition (Definition \ref{def:neardecomp}, Proposition
  \ref{prop:leakpair}).
  \emph{Not used}: as evidence that ``all complex systems necessarily have the same
  levels''; the source work does not prove this, nor does it give a general theorem about
  verification cost. \\
Abstract interpretation and Galois connections &
  \emph{Used}: as the formal tool for semantic preservation between levels (Definition
  \ref{def:galois}).
  \emph{Not used}: no claim that all inter-level relations can be expressed as Galois
  connections; where they cannot, the paper falls back on simulation relations or on an
  explicit human argument. \\
Minimum description length &
  \emph{Used}: as the formal background for ``comprehension is compression'' and the source
  of the stopping rule for descent (Definition \ref{def:mdl}, Proposition
  \ref{prop:stopping}).
  \emph{Not used}: no claim that the shortest description is the most correct description.
  MDL is a model-selection principle; its conclusions depend on the chosen model family,
  encoding and data, and it cannot by itself determine what purpose a system should
  achieve. \\
Information hiding and deep modules &
  \emph{Used}: as the normative basis for module boundaries.
  \emph{Not used} as it stands: the paper adds a testable version of it, namely whether an
  interface is a sufficient statistic of behaviour (Definition \ref{def:neardecomp}),
  because how well information is hidden is not itself measurable. \\
Correspondence between organisational structure and system structure &
  \emph{Used}: as one of the mechanisms of high-level convergence and one of the causes of
  divergence.
  \emph{Not used}: not as a technical necessity; it is a source of constraints, not a
  source of uniqueness for solutions. \\
Program logic and refinement &
  \emph{Used}: as the mathematical interface for the specification--implementation
  relation.
  \emph{Not used}: no claim that safety refinement automatically guarantees liveness,
  timing and fairness properties. \\
Proof-carrying code &
  \emph{Used}: as the classic implementation of a trust structure in which generation is
  free and admission is strict, and as the predecessor of the fourth pillar of this paper.
  \emph{Not used}: its premise is not ignored: the receiver must fix a safety policy in
  advance, and the correctness of the policy itself is an external problem. \\
Assume--guarantee reasoning &
  \emph{Used}: as the source of the rules by which local guarantees compose into a global
  guarantee (Proposition \ref{prop:soundness}).
  \emph{Not used}: its four composition conditions are not ignored; a component's being
  ``correct in isolation'' and a system's being ``correct as a whole'' are not
  interchangeable. \\
Resource analysis &
  \emph{Used}: as evidence that resource correctness can enter formal analysis.
  \emph{Not used}: a symbolic bound is not equated with a time bound on a particular
  processor; complexity order, worst-case execution time, statistical quantiles and
  measurement are not conflated. \\
Computational thermodynamics &
  \emph{Used}: as the basis for the boundary condition that the physical cost of computation
  is incompressible.
  \emph{Not used}: it is not used to predict the power consumption of a piece of ordinary
  code, nor folded into semantic correctness.
  (The paper mentions it only in Chapter \ref{ch:objections}, and not as a load-bearing
  structure.) \\
The halting problem and Rice's theorem &
  \emph{Used}: as the boundary at which a general synthesis algorithm and a general decision
  algorithm cannot exist.
  \emph{Not used}: they are not used to prove that ``intent can never be formalised''
  (Proposition \ref{prop:no-godel}). \\
$\mathrm{P}$ and $\mathrm{NP}$ &
  \emph{Used}: only the structural asymmetry entailed by the \emph{definition} of the class
  $\mathrm{NP}$ (Proposition \ref{prop:asym-struct}).
  \emph{Not used}: no open assertion of ``$\mathrm{P} \neq \mathrm{NP}$'' or
  ``$\mathrm{P} = \mathrm{NP}$'' serves as a load-bearing structure. \\
\bottomrule
\end{longtable}

\begin{principle}[Positioning of this paper]\label{pr:self-position}
The contribution of this paper lies at the \emph{interface layer}, not in a new semantic
theory: it promotes levels from a comprehension aid to the unit into which verification
obligations decompose; it formalises oblivious sets, near-decomposability and read-sets as
computable or measurable objects; it attributes the necessity of internal
infrastructuralisation to the completeness of the operational semantics rather than to code
volume; and it writes the trust structure as the composite of the two modes of trust
generation and gives a bound on consensus reliability. These interfaces connect existing
theory rather than replace it.
\end{principle}

% =====================================================================
\chapter{Objections and Limits}\label{ch:objections}
% =====================================================================

This chapter performs a self-audit: it removes claims that cannot bear load from the
structure, and records explicitly what has been removed.

\section{List of Non-Claims}

The following list is binding: no other part of the paper may be read as supporting the
claims below.

\begin{enumerate}[label=(N\arabic*)]
\item The paper does \emph{not} claim that all large software eventually converges on one
      and the same high-level architecture (Hypothesis \ref{hyp:strong} is explicitly
      rejected).
\item The paper does \emph{not} claim that code volume by itself causes internal
      infrastructuralisation; it claims that the cause is the product of entity count and
      resource count, together with the recurrence of problems and the frequency of
      evolution.
\item The paper does \emph{not} claim that the six functions must be \emph{built
      in-house}; it claims that they must be \emph{defined}, and may be adopted from a host.
\item The paper does \emph{not} claim that Simon's near-decomposability is a universal law.
      The paper recasts it as a conditional, testable sufficient-statistic criterion
      (Definition \ref{def:neardecomp}) and concedes that the original source provides no
      general theorem about verification cost.
\item The paper does \emph{not} claim that minimum description length can decide what
      purpose a system ought to achieve. MDL is a model-selection principle, not a theory
      of intent.
\item The paper does \emph{not} claim that ``comprehension is compression'' is a theorem in
      every domain. Within inductive inference it is a defensible theoretical position;
      outside that range it is a strong heuristic.
\item The paper does \emph{not} claim that a mechanical system never errs. It claims that
      mechanical decision does not depend on the generator's subjective will, and that a
      mechanical system can still be defective, can fail, and can embody an incomplete
      model.
\item The paper does \emph{not} claim that trust can be fully mechanised. The \TCB{} can be
      reduced, verified and isolated, but it cannot be eliminated outright.
\item The paper does \emph{not} claim that the intent layer can be closed mechanically. It
      can be \emph{approached} by adversarial consensus, and the reliability of that
      approximation is governed by the correlated error mode.
\item The paper does \emph{not} claim that ``adding more agents to review'' raises
      reliability. By Corollary \ref{cor:diversity}, reliability rises only when the
      correlated error mode falls.
\item The paper does \emph{not} claim that layering can accommodate everything. A
      cross-cutting concern cannot be localised structurally into any single layer
      (Principle \ref{pr:crosscutting}).
\item The paper does \emph{not} claim that formal verification will become cheap
      \emph{automatically} in the AI era. That is a hypothesis worth investigating, not an
      established fact (\S\ref{sec:asymmetry}).
\item The paper does \emph{not} claim that compiler correctness equals system correctness.
      Linking, loading, the runtime, external libraries, hardware and the operating system
      need not lie within the scope of the proof.
\item The paper does \emph{not} claim that resource guarantees can be unified into a single
      method. Time complexity, symbolic resource bounds, worst-case execution time,
      statistical quantiles and direct measurement are evidence of different grades.
\item The paper does \emph{not} claim that ``no violation detected'' equals ``every
      execution satisfies''. Three states must be distinguished (Principle
      \ref{pr:three-states}).
\item The paper does \emph{not} claim to determine who is authorised to adopt a specification.
      Principle \ref{pr:authority} establishes that authorisation is a failure axis independent of
      reliability; which body holds the standing is an institutional fact that the framework
      locates but does not settle.
\item The paper does \emph{not} claim a model of how specification fidelity decays over time.
      Section \ref{sec:spec-binding} argues that the quantity must be time-indexed and that a
      re-confirmation rule is forced by Proposition \ref{prop:three-grounds}(i); the form of
      $\varepsilon_{\text{cons}}(\cdot)$ is left open.
\item The paper does \emph{not} claim that the design-effect expression \eqref{eq:design-effect} is a
      theorem about majority error. It is the standard approximation for a mean under exchangeable
      correlation, and is used here to make Corollary \ref{cor:diversity} quantitative, not to
      establish a bound.
\end{enumerate}

\section{Arguments Removed from the Load-Bearing Structure}

The paper deliberately rejects several arguments that are common in discussions of this
kind but do not hold. They are rejected not because they are useless, but because they
state heuristics as theorems, and thereby make the whole framework depend on unproved
conjectures.

\begin{longtable}{L{3.4cm}L{5.2cm}L{5.6cm}}
\caption{Arguments removed from the load-bearing structure}\label{tab:overreach}\\
\toprule
\textbf{Rejected argument} & \textbf{Why it fails} & \textbf{Replacement in this paper} \\
\midrule
\endfirsthead
\toprule
\textbf{Rejected argument} & \textbf{Why it fails} & \textbf{Replacement in this paper} \\
\midrule
\endhead
\bottomrule
\endlastfoot
Deriving from G\"odel's incompleteness theorem that ``a permanent gap between intent and
  specification necessarily exists'' &
  The incompleteness theorem concerns limits of provability in particular formal systems;
  it is not a theorem about the relation between human mental intent and formal
  specification &
  Replace it with three epistemic grounds: requirements are generative; specification is
  lossy compression; and the distinction between verification and validation
  (Proposition \ref{prop:three-grounds}) \\
Deriving from Rice's theorem that ``intent cannot be formalised'' &
  Rice's theorem concerns the undecidability of non-trivial semantic properties of
  programs; it has nothing to do with human intent &
  Use it only as a boundary on general synthesis and general decision algorithms (Q6 of
  Remark \ref{rmk:q1q8}) \\
Calling ``grading is easier than answering'' a dividend of $\mathrm{P}\neq\mathrm{NP}$ &
  Whether $\mathrm{P}$ equals $\mathrm{NP}$ is still open; the argument depends on an
  unproved conjecture &
  Use only the structural asymmetry implied by the definition of the class $\mathrm{NP}$
  (Proposition \ref{prop:asym-struct}) \\
Asserting as a mathematical theorem that ``the verifier must be minimal, deterministic and
  never isomorphic'' &
  A small trusted kernel is a valuable architectural principle; small does not mean free of
  defects &
  \TCB{} minimisation: make the untrusted part inspectable, and concentrate the unavoidable
  trust on a basis that is as small, and as explicit in its properties, as possible \\
Asserting that ``a dead thing does not betray'', hence mechanical assurance is absolutely
  reliable &
  A mechanical system can be defective, can fail, and can embody an incomplete model &
  Distinguish ``does not betray'' (does not depend on the generator's will) from ``cannot
  err'' (impossible) \\
Asserting that Simon's law of levels proves that every complex system necessarily organises
  itself into levels &
  The original source offers an explanation, not a universal law, and no general theorem
  about verification cost &
  Recast as a conditional, testable sufficient-statistic criterion
  (Definition \ref{def:neardecomp}) \\
Asserting that linear logic makes wasted resources structurally impossible &
  The exponential structure of linear logic still permits particular forms of duplication
  and discarding; linear use is not optimal use &
  Use a composite structure: usage constraints $+$ resource-bound analysis $+$ performance
  models $+$ runtime instrumentation $+$ qualified worst-case execution time analysis \\
Asserting that formal verification must become cheap in the AI era &
  Historical cost magnitudes show that specification design, model selection and proof
  architecture remain real work (Example \ref{ex:sel4-cost}) &
  Downgrade it to a hypothesis and give a measurable cost metric (criterion S5 of Chapter
  \ref{ch:program}) \\
Asserting that ``comprehension is compression'' is a universal law &
  Its rigorous formalisation applies to inductive inference, and beyond that range it
  degenerates into a heuristic &
  Restrict the scope of application and tie it to task-relative sufficiency
  (Definition \ref{def:sufficient}) \\
Asserting a substantive connection between this framework and the Langlands programme &
  That programme studies the functorial duality between arithmetic objects and automorphic
  representations; the ``free generation / deterministic admission'' this framework needs
  is only the \emph{smell of a duality}, not a functor. To write it in would substitute
  rhetoric for structure &
  Adopt a duality that pays its way: the LP/SMT duality. Farkas certificates, unsatisfiable
  cores and unrealisability certificates give a theorem-level correspondence between a
  synthesis problem and its infeasibility proof, and are already used in engineering \\
Asserting that the Landauer principle constrains software design &
  The scale of software is some twelve orders of magnitude away from that bound; the
  erasure bound is not a constraint surface for architecture &
  Adopt the memory wall and locality: moving one bit costs about $10^{9}$ times the erasure
  bound in energy, and this is the physics that actually fixes the memory hierarchy, cache
  policy and data layout \\
Asserting that homotopy type theory is a necessary foundation for this framework &
  That theory extends the foundations of mathematics, but gives no software guarantee that
  resource analysis or verified compilation does not already give &
  Restrict the role of type theory to what has been proved: dependent types expressing
  postconditions and invariants \cite{hott2013} \\
\bottomrule
\end{longtable}

\section{Possible Objections and Responses}

\paragraph{Objection 1: read-sets and oblivious sets cannot be measured precisely in real
systems, so the formalisation is decorative.}
Response: the read-set of Definition \ref{def:readset} can be counted directly from version
history, using ``the set of modules touched by each change'' as a proxy; it is a computable
quantity (Hypothesis $H_4$ exists precisely to test the validity of that proxy). The
oblivious set is not required here to be extracted automatically; it is required to be
\emph{explicitly declared}. It is a design artefact, not an analysis artefact. Once
declared, it can be checked in part by the leak-pair diagnostic (Corollary
\ref{cor:diagnose}).

\paragraph{Objection 2: the core conclusion that ``a large system grows a kernel'' is too
broad and almost unfalsifiable.}
Response: Theorem \ref{th:six} has explicit premises (more than one concurrently live
entity, finite contended resources, reachable faults), and its conclusion is falsifiable: a
well-defined system that satisfies the premises but lacks one of the six functions suffices
to refute it. The empirical part rests on Hypothesis $H_3$, whose falsifier is that ``lines
of code explain the data no worse than $n \times m$''.

\paragraph{Objection 3: the two-layer trust structure places too much responsibility on
consensus, and consensus often fails in practice.}
Response: this is exactly the problem Proposition \ref{th:correlated} addresses. The
proposition does not defend consensus; it gives the form of its failure and its design
variables. It also explains a common phenomenon: a ``review panel'' composed of models of
common origin is structurally not a guarantee, however large it is.

\paragraph{Objection 4: does treating the level as the unit into which verification
obligations decompose merely rename the existing practice of ``layered verification''?}
Response: in part, yes, and the paper concedes it. The difference lies in three points:
(a) the assignment rule given here (Definition \ref{def:obligation-assign}) is an
operational decision, not a slogan; (b) the paper explicitly introduces the class of
``obligations this layer cannot express'' and requires that it be assigned to consensus or
to human approval rather than ignored; (c) the paper gives a quantitative bound on the
reliability of the consensus layer, which makes ``assign it to consensus'' a decision that
can be assessed rather than merely trusted.

\paragraph{Objection 5: defining technical debt as a description-length gap is
uncomputable and therefore useless.}
Response: the value of the definition lies not in numerical computation but in the
direction and the criterion it gives (Principle \ref{pr:debt}(c)): deciding whether a
refactoring is worthwhile is equivalent to deciding whether it makes the system easier to
generate from a shorter description. That is a question that can be used in review, even
though $\DL^{\star}$ is uncomputable.

\paragraph{Objection 6: the paper mixes engineering, mathematics and social epistemology;
does it exceed the defensible range of a single discipline?}
Response: the criticism holds. Two mitigations are in place. First, every class of claim is
labelled with its strength (the strength trichotomy of Chapter \ref{ch:program}); second,
the citations that are \emph{used} are explicitly separated from those that are
\emph{rejected} (Table \ref{tab:relations}). Cross-disciplinarity is not itself a defect,
provided that no discipline's conclusion is used outside its scope.

\section{Open Problems}

The following questions are open:

\begin{enumerate}[label=(O\arabic*)]
\item \term{Independent measurability of $\Dom$ and $\Conv$}. For Conjecture
      \ref{conj:strength} to become a testable proposition, both quantities must be
      measurable independently of the conclusion. The paper gives only the direction, not a
      reproducible annotation protocol.
\item \term{Formalisation of the completeness argument for the six functions}. The proof of
      Theorem \ref{th:six} is a structural enumeration, not a derivation in a formal system.
      Writing it as a completeness theorem for a given operational-semantics framework
      remains unfinished work.
\item \term{Whether cross-layer semantics always compose}. Proposition \ref{prop:soundness}
      gives sufficient conditions, but whether those conditions can be satisfied
      simultaneously in systems containing nondeterminism, real-time constraints and
      reflexivity is not known.
\item \term{How to measure $q$ for the consensus layer}\label{op:q}. Proposition \ref{th:correlated}
      attributes reliability to the correlated error mode, but gives no operational method
      for estimating $q$. This is the crux of whether the direction can be put into
      practice. Section \ref{sec:consensus} supplies a companion quantity, the mean pairwise
      error correlation $\bar\rho$, which is estimable from data the framework already collects;
      $q$ remains the limit that is not. Appendix \ref{app:protocol} makes the same quantities
      addressable at the level of procedure---attestor roster, evidence provenance and dissent
      record are the fields from which $\bar\rho$ can be estimated and against which selection
      capture becomes visible---but it neither measures $q$ nor claims to.
\item \term{That the consensus protocol of Appendix \ref{app:protocol} works}. The appendix
      specifies a procedure and states its failure modes; it has not been built, run or
      evaluated, and no parameter in it was measured. Its role is to make the intent layer's
      obligations concrete enough to criticise, not to discharge them.
\item \term{Whether a joint semantics of resources and function is decidable}. A costed
      trace semantics can express joint constraints, but the size of its decidable fragment
      and the tightness of its bounds are unsettled.
\item \term{Compatibility of artefact binding with dynamically generated implementations}.
      If an implementation can be regenerated at runtime in response to load, the static
      identity assertion of Equation \eqref{eq:binding} must be replaced by a stronger
      question: which changes preserve the existing contract, which require re-verification,
      and how to fall back to a known safe state when verification fails.
\item \term{Consistency of level assignment}. Whether Definition
      \ref{def:obligation-assign} yields consistent assignments across different analysts
      requires empirical testing.
\end{enumerate}

\section{Methodological Limitations}

\begin{enumerate}[label=(L\arabic*)]
\item For every reference cited, the depth of checking is limited to the formal publication
      page, the author's or project's home page, and preprint metadata; the full text was
      \emph{not read through in full for each entry}. Every concrete number given in the
      text should therefore be treated as a \emph{second-hand report awaiting re-checking};
      several entries with recent nominal years are retained as recorded in the original
      source, and no independent checking was done here.
      A verification round was carried out for the first revision. The five empirical
      citations supporting Hypothesis $H_7$ were checked against the primary abstract pages:
      the four figures quoted (189 Lean tasks; $72.6\%$, $52.3\%$ and $4.9\%$ for the best
      model) and the 245 operating-system tasks in the second benchmark were all found
      \emph{in the abstracts themselves}, not merely in secondary summaries. Two venue
      corrections were made as a result: the traceable-refinement work is a workshop poster
      rather than a main-conference paper, and the mutation-driven method is an unrefereed
      preprint. Newly added citations were checked to the same depth, and the two entries
      whose full text was not read are marked ``metadata only'' in the bibliography.
      What the round did \emph{not} establish: no number was re-derived from a paper's tables
      or figures, and no full text was read end to end, so the grade of the numeric claims is
      still ``quoted from a primary abstract'', one step above the original second-hand
      report and short of independent replication.
\item The formalisation in this paper consists of \emph{working definitions established for
      the purposes of this paper}, and should not be mistaken for an existing unified
      standard. In particular, the three items ``oblivious set'', ``level contract
      four-tuple'' and ``level assignment of verification obligations'' are working
      concepts proposed here.
\item The three epistemic grounds (Proposition \ref{prop:three-grounds}) are arguments
      rather than theorems; the challengeable point is whether the three are exhaustive.
\item No experiments were performed. Every hypothesis in Chapter \ref{ch:program} remains
      in the ``to be tested'' state.
\end{enumerate}

\begin{principle}[Self-audit conclusion]\label{pr:self-audit}
After this scrutiny, three minimal independently testable propositions remain in the
load-bearing structure:
\begin{enumerate}[label=(\roman*)]
\item a structural asymmetry between verification and generation, implied by the definition
      of the class $\mathrm{NP}$;
\item the level is the unit into which assurance obligations decompose, because every
      obligation belongs to some observable domain;
\item the gap between intent and specification has epistemic grounds, and therefore cannot
      be closed mechanically in general.
\end{enumerate}
Every remaining claim (including ``compression'', ``emergence'' and ``internal
infrastructuralisation'') has either been downgraded to a falsifiable hypothesis or
confined to an explicit scope of application.
\end{principle}

% =====================================================================
\chapter{Conclusion}\label{ch:conclusion}
% =====================================================================

\section{Answers to the four questions}

\begin{problem}[Generality of low-level determinism]
The answer has three parts. The low level is nearly unique because it is
\emph{over-constrained}: the number of constraints is at least the number of design
degrees of freedom, and the solution set degenerates into a discrete point set
(Proposition \ref{prop:overconstrain}). The upper levels admit many solutions because
problem-level constraints are far fewer than the degrees of freedom, and the solution
set is a high-dimensional manifold. The upper levels nevertheless converge, because a
\emph{second set of constraints}, imposed by the human organisations that build and
maintain the system, compresses the manifold into a few sharp peaks
(Principle \ref{pr:org}).

The accurate answer is therefore this: convergence exists, but not towards the
engineering optimum; it is convergence towards a \emph{structure that a bounded
organisation can navigate}. What converges is carried by \emph{relations} rather than by
components (Principle \ref{pr:relations}); its strength decreases monotonically with
domain binding (Conjecture \ref{conj:strength}).
\end{problem}

\begin{problem}[Coercion by scale]
The conclusion is a \emph{conditional necessity} result. The independent variable that
decides whether a unified substrate is required is the product of the number of live
entities and the number of shared resources, not the line count
(Theorem \ref{th:integration}, Corollary \ref{cor:threshold}). The six functions
(identity and lifecycle, allocation and reclamation, scheduling, naming and addressing,
communication and synchronisation, and fault isolation and recovery) are necessary
conditions for the operational semantics to be well defined (Theorem \ref{th:six}). They
may be built or adopted from a host; the criterion is whether the invariant they carry
is domain-critical (Corollary \ref{cor:buildadopt}).
\end{problem}

\begin{problem}[Comprehension and design]
The conclusion is that comprehension is inverse compression and design is forward
compression, the two reading directions of one structure. The optimal descent depth of
comprehension is given by the stationarity condition of rate--distortion
(Proposition \ref{prop:stopping}); the measurable objective of design is the expected
read-set (Definition \ref{def:readset}), and minimising it requires the enforcement of
each invariant at a single site (Theorem \ref{th:centralise}). Whether an abstraction
leaks is measurable, not merely felt (Proposition \ref{prop:leakpair},
Corollary \ref{cor:diagnose}). For an AI the advantage lies not in reading more but in
probing faster (Proposition \ref{prop:active}).
\end{problem}

\begin{problem}[Correctness without review]
The conclusion is this: it is attainable that no human reviews the implementation line
by line, and unattainable that no human remains responsible for the expression of the
goal. ``Correctness'' must be decomposed into five dimensions
(Definition \ref{def:five-correctness}); the intermediate process must be defined as a
seven-stage verifiable refinement chain (Definition \ref{def:chain}); trust must be a
composite of the two modes of trust generation (Definition \ref{def:two-layer}); and
artefact binding is the fourth pillar and cannot be omitted (Eq. \eqref{eq:binding}). The
consolidated assurance condition (Theorem \ref{th:assurance}) gives the complete answer.
\end{problem}

\section{Compression as the through-line}

Gathering the whole paper together yields a single thread: \term{compression}.

\begin{itemize}
\item Complex systems compress themselves by \emph{levels}, so that every large system
      grows a kernel internal to its domain (Chapter \ref{ch:infrastructure});
\item Bounded intelligence compresses information by \emph{comprehension}, so that a
      finite cognition accommodates an unbounded world (Chapter \ref{ch:comprehension});
\item Engineering compresses the cost of change by \emph{architecture}, so that a local
      modification does not require global comprehension (Chapter \ref{ch:design});
\item Civilisation compresses trust by \emph{rules}, so that a million lines of code of
      unknown provenance can be run safely (Chapter \ref{ch:intent}).
\end{itemize}

The mechanism of compression is the same in all four cases: quotienting by the kernel of
the abstraction map (Proposition \ref{prop:oblivious}). \emph{All the information a level
exposes to the level above is the quotient of lower-level behaviour by that kernel.}

\section{The core proposition}

If the whole paper is to be compressed into one testable sentence, this paper chooses the
following:

\begin{principle}[Core proposition]\label{pr:core}
A complex system is comprehensible to a bounded intelligence if and only if it is
compressible; and it is compressible if and only if it admits a separation of scales in
which every level has an explicit, declared \emph{oblivious set}.
\end{principle}

Three consequences follow directly.

First, it turns ``comprehensibility'' from a subjective impression into a structural
property: comprehensibility is equivalent to the existence of a quotient map under which
task-relevant predictive power is preserved. This makes comprehensibility an
\emph{objective function} of design rather than a retrospective assessment.

Second, it redefines ``architecture'' as \emph{manufactured separation of scales}. Nature
obtains levels free from separation of scales and renormalisation; software levels do not
emerge spontaneously and must be paid for in interfaces, encapsulation and invariants,
and they are liable to leak. The core deliverable of architecture is therefore not a
module list but the \emph{declaration of the oblivious set of every level}.

Third, it gives a concrete rule of action for an AI: do not attempt to exhaust low-level
information; maintain a layered, falsifiable compression model that counterexamples can
update; obtain gradient by probing rather than by reading; and take task sufficiency as
the stopping condition.

\section{Closing remarks}

The paper begins with three engineering intuitions and concludes with a joint proposition
about structure and assurance. What is established is not the answers but \emph{which
questions have no answer}, and why:

\begin{itemize}
\item the gap between intent and specification has no mechanical solution, and its
      grounds are epistemological rather than logical; it is in substance the frame
      problem, and is therefore avoidable only in interaction;
\item the basis of trust for mechanical assurance can be narrowed but not eliminated;
      and ``verification is cheap'' must be decomposed into two independent propositions,
      cheap checking time and cheap certificate size, of which the latter is the usual
      bottleneck;
\item the reliability of consensus assurance is dominated by correlated error modes, so
      that \emph{diversity} rather than \emph{number} is the design variable, and
      diversity is ultimately a problem of incentive design, not of statistics;
\item the structural convergence of large systems is probabilistic, domain-specific and
      shaped by organisational constraints, rather than unique;
\item the duty of architecture is not to lower the cost of a single change but to lower
      the \emph{exponent} with which that cost grows with size.
\end{itemize}

Set side by side, these items share one shape: \emph{complexity compression is a
necessity driven by size, and the form the compression takes is a contingency driven by
domain; while the boundary of assurance falls exactly between what a machine can decide
and what must be validated by the party that bears the consequences.}

The price is the surrender of a picture expressible as a single theorem; the gain is that
every claim sits where it can be tested. For a framework about how to establish
trustworthy conclusions on top of untrusted complexity, this is the only self-consistent
ending: \emph{it must pass its own test.}


\appendix
% =====================================================================
\appendix

\chapter{Supplementary Proofs and Arguments}\label{app:proofs}

This appendix supplies complete proofs, or fuller arguments, for several results in the
main text. Where a complete proof is already given there, it is not repeated.

\section{Notes on Proposition \ref{prop:overconstrain}}

The proof given in the main text rests on the transversality theorem. Two points are added
here.

\paragraph{Why ``independent constraints'' are necessary.}
If two constraints are not independent (for instance, when the second is a consequence of
the first), their contributions to the codimension do not add,
$\mathrm{codim}\,\mathcal{S} < c$, and the solution set has dimension greater than
$d - c$. The conclusion of Proposition \ref{prop:overconstrain} must therefore be read
under the premise that the constraints are independent and transversal. That premise
usually holds in low-level design, where each constraint originates in a distinct
pre-existing object (binary compatibility, the debugger, exception unwinding, the
system-call interface), and often fails in high-level design, where several
``constraints'' are in fact different formulations of one preference.

\paragraph{Why ``solution-set degeneracy'' does not mean ``unique solution''.}
When $c \ge d$ the solution set is finite or empty, and a finite set may contain more than
one element. Low-level design is nonetheless ``nearly unique'' in practice because, over
and above the formal constraints, an implicit economic constraint is at work: maintaining
two incompatible calling conventions at once carries labour and toolchain costs, and that
cost induces a strict preference order among the admissible options.

\section{Complete Argument for Theorem \ref{th:integration}}

The main text gives the principal argument. The choice of cost function and the
reasonableness of the assumptions are supplied here.

\paragraph{The definition of cost.}
Let the unit interaction cost $c$ be the whole of the work required to handle one
interaction between a particular entity and a particular resource correctly: design,
implementation, testing and long-term maintenance. In the no-substrate case every entity
must do this work independently, and must handle the resource invariants independently as
well.

\paragraph{The magnitude of the consistency obligation.}
For a resource $r$, $n$ independent implementations must agree on the invariants of $r$.
If consistency is not enforced, the deviation of any one implementation breaks the global
invariant. If it is enforced, the enforcement itself requires that all $n$ sites know the
invariant, so the number of enforcement sites is $\Omega(n)$. Neither case improves on the
conclusion ``read-set $\ge k$'' (Theorem \ref{th:centralise}). In other words, the absence
of a substrate is not merely more costly: it forces invariants to be scattered, and
thereby \emph{permanently raises} the read-set lower bound of the system. This is the
second benefit of a unified substrate, and it is independent of amortisation.

\paragraph{Relaxing assumption (b).}
Assumption (b), that every entity participates in all $m$ resources, is the worst case. If
entity $e$ participates in only $|\mathcal{R}(e)|$ resources, the no-substrate cost is
$\sum_e |\mathcal{R}(e)|$. As long as $\mathbb{E}_e[|\mathcal{R}(e)|]$ grows with $m$
(that is, resource use is diffuse and broad), the product-form growth still holds. The
qualitative conclusion of \eqref{eq:threshold} therefore does not depend on the worst-case
assumption.

\section{Supplement to the Item-by-Item Argument for Theorem \ref{th:six}}

The main text argues item by item, showing that the omission of any one function yields
undefined behaviour. Two remarks on the strength of that argument are added here.

\paragraph{Strength of the (F3) scheduling argument.}
The argument relies on the claim that $\to$ must be a function, or that its
non-determinism must be explicitly specified. If the semantics of a system is defined as
the set of \emph{all} interleavings, that is, if non-determinism is taken as part of the
semantics and no selection among it is made, then the selection function may indeed be
absent. In that case, however, the \emph{implementation} of the system must still select
one interleaving in every concrete execution, so the scheduling function remains present
at the implementation level even though it does not appear in the abstract semantics. The
formulation used in this paper, ``either in the system's own code or in its host'', covers
this case.

\paragraph{Strength of the (F6) isolation argument.}
The argument relies on fault reachability. If the execution environment of a system
guarantees the absence of faults (pure functional computation, for instance, or execution
in a restricted environment verified to be fault-free), then (F6) may be absent or may
degenerate to the trivial. This is precisely why some minimal systems can do without the
complete six functions, and it shows that the premises of Theorem \ref{th:six} are not
redundant.

\section{The Read-Set Lower Bound of Theorem \ref{th:centralise}}

The main text proves $|\rho(c)| \ge |\mathrm{S}(\iota)|$. Why that inequality is tight is
explained here.

If $\mathrm{S}(\iota) = \{M_0\}$ and the interface of $M_0$ is sufficient
(Definition \ref{def:neardecomp}), then every change that could break $\iota$ must be
expressed at the interface of $M_0$, so an agent performing the change need only
understand $M_0$ and its call sites. Then $|\rho(c)| = 1 + O(\text{call sites})$, and the
lower bound $1$ is attained (when the change itself occurs inside $M_0$).

If $\mathrm{S}(\iota) = \{M_0, M_1\}$ and both modules bear partial responsibility for
$\iota$, the agent must understand both \emph{and the consistency rule between them},
whose comprehension cost is strictly greater than that of the two modules together. The
read-set therefore grows superlinearly in $|\mathrm{S}(\iota)|$, and the linear relation is
only a lower bound. The bound given in the main text is consequently conservative.

\section{Exhaustiveness of the Three Grounds of Proposition \ref{prop:three-grounds}}

This argument has limits. The three grounds take the following form:

\begin{enumerate}[label=(g\arabic*)]
\item Preferences are generative: there exist an intent $I$ and a situation $s$ such that
      some component of $I$ is fixed only after the bearer of the intent encounters $s$.
\item Specifications are lossy compressions: there exist an intent $I$ and a specification
      $\Phi$ such that $\Phi$ cannot carry all information relevant to $I$.
\item Validation does not reduce to verification: there exist a requirement $R$ and an
      implementation $P$ such that ``$P$ satisfies $\Phi$'' is true while ``$P$ satisfies
      $R$'' is false, and the procedure that decides the latter cannot be reduced to the
      procedure that decides the former.
\end{enumerate}

The three grounds are not independent of one another: (g1) is a source of (g2), and (g3)
is a consequence of (g2). A more accurate formulation may therefore be that \emph{the root
of the intent gap is the composite of (g1) and (g2)}. This paper retains the three
parallel formulations because they correspond to different engineering remedies: (g1) to
interactive refinement and counterexample generation, (g2) to traceability and coverage
analysis, and (g3) to an independent validation process and an allocation of
responsibility.

Whether a fourth ground exists is an open problem (Chapter \ref{ch:objections}, items
(O)).

\section{Tightness of the Bound in Theorem \ref{th:condorcet}}

Equation \eqref{eq:chernoff} is the Hoeffding bound, which is not tight for large $n$. For
$p < 1/2$ the probability of a majority error is in fact given by the central limit
theorem, $\Pr[\text{majority errs}] \approx \Phi\bigl(-\sqrt{n}\,(1/2-p)/\sigma\bigr)$,
whose exponent is $\Theta(n (1/2-p)^{2})$ with a better constant.

This paper chooses the Hoeffding bound over a tighter estimate for three reasons:

\begin{enumerate}[label=(\alph*)]
\item it requires no assumption about the distribution of individual errors, only
      independence, boundedness and a mean;
\item its form is simpler, and it is therefore less easily misread as an exact value;
\item the central conclusion of this paper (Proposition \ref{th:correlated}) does not
      depend on the tightness of the bound: the lower bound $q$ given by a correlated
      error mode is independent of the constant in the upper bound.
\end{enumerate}

\section{A Stronger Form of Proposition \ref{th:correlated}}

The main text gives $\Pr[\text{majority errs}] \ge q$. A finer form is available. Let the
common error mode $C$ and the individual independent errors $I_i$ satisfy: verifier $i$
errs if and only if $C = 1$ or $I_i = 1$. The event that the majority errs is then
\[
\{C = 1\} \;\cup\; \Bigl\{C = 0,\; \sum_{i=1}^{n} I_i \ge \tfrac{n}{2}\Bigr\},
\]
so that
\[
\Pr[\text{majority errs}] \;=\; q + (1-q)\,\Pr\Bigl[\sum_{i=1}^n I_i \ge \tfrac{n}{2}\Bigr]
\;\ge\; q .
\]
For fixed $q > 0$, moreover,
$\lim_{n \to \infty} \Pr[\text{majority errs}] = q$. This yields a statement stronger than
a lower bound of $q$: \emph{consensus reliability has an asymptotic floor determined by
the correlated error mode; increasing the number of verifiers can only approach that
floor and cannot break through it}. This is the quantitative basis of
Corollary \ref{cor:diversity}.

\section{Notes on Proposition \ref{prop:active}}

The argument in the main text uses the definition of information gain. One clarification
is added here: Proposition \ref{prop:active} does not say that passive reading is without
value. Passive reading supplies the \emph{prior}, in the Bayesian sense, whereas active
probing supplies the \emph{likelihood}. A prober with no prior does not know what to
probe; a reader with only a prior cannot update. The first four of the seven steps of the
protocol in \S\ref{sec:protocol} are therefore reading (establishing the prior), and the
last three are probing (acquiring the gradient). Their order is not arbitrary.

\chapter{Templates}\label{app:templates}

This appendix provides three templates for direct use. They convert the formal objects of the
main text into engineering documents.

\section{Architecture Map Template}

The template records the level ladder of a system. Each row is one level.

{\small\begin{longtable}{L{1.1cm}L{2.2cm}L{2.9cm}L{2.5cm}L{2.6cm}L{2.3cm}}
\caption{Architecture map}\label{tab:tpl-architecture}\\
\toprule
\textbf{Level} & \textbf{Participant} & \textbf{Guaranteed invariants} &
\textbf{Assumed of the level below} & \textbf{May be ignored} &
\textbf{Verification method} \\
\midrule
\endfirsthead
\toprule
\textbf{Level} & \textbf{Participant} & \textbf{Guaranteed invariants} &
\textbf{Assumed of the level below} & \textbf{May be ignored} &
\textbf{Verification method} \\
\midrule
\endhead
\bottomrule
\endlastfoot
$\Lk{4}$ & & & & & \\
$\Lk{3}$ & & & & & \\
$\Lk{2}$ & & & & & \\
$\Lk{1}$ & & & & & \\
$\Lk{0}$ & & & & & \\
\end{longtable}}

The following requirements apply when filling in the table:

\begin{enumerate}[label=(\arabic*)]
\item The ``may be ignored'' column must not be left blank. A blank means the level is not a
      level (Principle \ref{pr:level-real}).
\item ``Guaranteed invariants'' must be written as propositions that can be checked, or at
      least decided.
\item ``Verification method'' must state which class of assurance in
      Table \ref{tab:assurance-matrix} the invariant falls under (mechanical, monitoring,
      consensus).
\end{enumerate}

\section{Module Contract Template}

\begin{description}
\item[Identity] Module name, the level it belongs to, unique identifier.
\item[Responsibility] What problem it solves (one sentence, no longer than one line).
\item[Boundary] What it is responsible for, and what it is explicitly not responsible for.
\item[External interface] The functions, types or interfaces through which the level above
      uses it.
\item[Precondition $P$] What the caller must satisfy.
\item[Postcondition $Q$] What is guaranteed on return.
\item[Invariant $I$] Where it is enforced; why circumvention is impossible.
\item[Side effects $E^{+}$ / $E^{-}$] What is permitted to occur; what is explicitly
      forbidden.
\item[Resource budget $R$] Upper bounds on time, space, bandwidth and concurrency, and their
      cost model.
\item[Failure semantics $F$] What happens on error, cancellation, retry and recovery.
\item[Environmental assumptions $A$] The hardware, operating system, load and external-service
      conditions on which the conclusions depend.
\item[Ownership and lifecycle] How key objects are created, initialised, closed and destroyed;
      who owns them.
\item[Concurrency model] Thread-safety guarantees, lock ordering, reentrancy.
\item[Dependencies] Which modules it depends on; whether the dependency is on an abstract
      interface or on a concrete implementation.
\item[\textbf{Oblivious set}] What this module is permitted never to care about (required; see
      Definition \ref{def:oblivious}).
\item[Verification requirements] Which tests or checks can demonstrate the above conclusions;
      whether there is an isolable unit of verification.
\item[Leak-pair risk] Known scenarios in which ``the interface is the same but the external
      behaviour differs''.
\item[Open questions] Design questions not yet settled.
\end{description}

\section{Decision Record Template}

\begin{description}
\item[Identifier and title] \code{ADR-XXXX}: decision title.
\item[Status] Draft / accepted / superseded / deprecated.
\item[Context] What situation forces the decision to be made.
\item[Decision drivers] The constraints and goals that bear on the decision (performance,
      consistency, team size, rate of evolution, cost).
\item[Candidate options] Listed one by one, each with its assumptions and its cost.
\item[Decision] Which one is chosen, and the reason in one sentence.
\item[Impact] What the decision makes easier and what it makes harder.
\item[Full argument for the excluded options] Why the other options were rejected. This item
      is the most valuable part of a decision record: it records the information and judgement
      available at the time, so that future readers need not derive it again.
\item[Verification] How the consequences of the decision are checked; which external evidence
      would show that it is wrong.
\item[Reassessment conditions] What circumstances should prompt this decision to be revisited.
\end{description}

\begin{principle}
If an architecture is a physical fact like a low-level ABI, there is no need to write a
decision record.
\emph{The ubiquity of decision records is itself evidence that high-level architecture has no
single correct answer.}
\end{principle}

\chapter{Glossary}\label{app:glossary}

\begin{longtable}{L{3.4cm}L{11.0cm}}
\toprule
\textbf{Term} & \textbf{Definition and location} \\
\midrule
\endfirsthead
\toprule
\textbf{Term} & \textbf{Definition and location} \\
\midrule
\endhead
\bottomrule
\endlastfoot
Level & The six-tuple $(X, \Beh, \Obs, \Env, \Cost, \Inv)$; a cognitive scale, not a
  physical directory. Definition \ref{def:level}. \\
Abstraction map & The map $\Obs_i$ projecting lower-level behaviour onto upper-level
  behaviour. Definition \ref{def:galois}. \\
Concretisation map & The map $\Conc_i$ expanding an upper-level behaviour into the set of
  its lower-level implementations. Definition \ref{def:galois}. \\
Galois connection & $\Obs_i \dashv \Conc_i$, stating what abstraction discards and what it
  retains. Definition \ref{def:galois}. \\
Inter-level soundness & $\Obs_i(\Beh_{i+1}\!\upharpoonright\!A) \subseteq \Beh_i$.
  Definition \ref{def:soundness}. \\
Oblivious set & The kernel $\Ker \Obs_i$ of the abstraction map, that is, the set of
  differences the level can in principle never distinguish. Definition \ref{def:oblivious}. \\
Compression ratio & The information that cannot be recovered under the abstraction map.
  Definition \ref{def:compression}. \\
Architecture relational graph & A typed multigraph whose edges fall into five classes:
  abstraction, dependency, composition, control and contract. Definition
  \ref{def:relgraph}. \\
Abstract skeleton & The graph obtained by replacing each component with its duty role,
  retaining the type of each relation and erasing its instance. Definition
  \ref{def:skeleton}. \\
Task & The pair $T = (Q, d)$, where $Q$ is a query family and $d$ a loss on answers.
  Definition \ref{def:task}. \\
$T$-sufficiency & The model answers every query as the real system does. Definition
  \ref{def:sufficient}. \\
Description length & The length $\DL(M)$ of a model under a fixed encoding. Definition
  \ref{def:mdl}. \\
Behaviour contract & The eight-tuple $(P, Q, E^{+}, E^{-}, R, I, F, A)$. Definition
  \ref{def:contract}. \\
Observable traces & The set $\Tr(B,E)$ of observable execution traces an artefact may
  produce in an environment. Definition \ref{def:traces}. \\
Assurance & The four-tuple (property, object, premises, evidence); the condition under
  which it holds is \eqref{eq:assurance}. Definition \ref{def:assurance}. \\
Trusted computing base (TCB) & The set of components that must be trusted unconditionally
  for the assurance to hold. Definition \ref{def:tcb}. \\
Mechanical assurance & Assurance given by a single deterministic, exhaustively inspectable
  decision procedure. Definition \ref{def:twotrust}. \\
Consensus assurance & Assurance given by the adversarial scrutiny of multiple independent
  participants. Definition \ref{def:twotrust}. \\
Internal infrastructuralisation & The three-stage process: identification of commonality
  $\to$ unification of mechanisms $\to$ upper-layer composition. Definition
  \ref{def:infra}. \\
Participant-resource structure & The five-tuple
  $(E, \mathcal{R}, \to, \mathrm{Ready}, \mathrm{Fault})$. Definition
  \ref{def:pr-structure}. \\
Unified substrate & The set of mechanisms through whose declared interfaces every
  interaction of an entity with a resource passes. Definition \ref{def:substrate}. \\
The six functions & Identity and lifecycle; allocation and reclamation; scheduling; naming
  and addressing; communication and synchronisation; fault isolation and recovery.
  Definition \ref{def:six}. \\
Read-set of a change & The set $\rho(c)$ of modules that must be read to complete a change
  correctly. Definition \ref{def:readset}. \\
Enforcement sites of an invariant & The set $\mathrm{S}(\iota)$ of modules that must be
  implemented correctly for the invariant to hold. Definition \ref{def:enforce}. \\
Near-decomposability & A module's interface observations are a sufficient statistic for
  its external behaviour. Definition \ref{def:neardecomp}. \\
Leak pair & Two executions with equal interface trace and different external behaviour.
  Proposition \ref{prop:leakpair}. \\
Technical debt & The difference between the description length of the code and its shortest
  generating description. Definition \ref{def:debt}. \\
Verifiable refinement chain & The seven-stage chain from purpose to certified deployment,
  $(\mathrm{G}1)\dots(\mathrm{G}7)$. Definition \ref{def:chain}. \\
The five dimensions of correctness & Functional, side-effect, resource, system and
  goal-expression correctness. Definition \ref{def:five-correctness}. \\
Four classes of side effect & Functional, resource, security and semantic drift.
  Definition \ref{def:effects}. \\
Artefact evidence package & The ten items of evidence every artefact should carry.
  Definition \ref{def:evidence-package}. \\
Cognitive scale & $\Lk{0}$ principles of computation, $\Lk{1}$ infrastructure and runtime,
  $\Lk{2}$ components and subsystems, $\Lk{3}$ domain function and business semantics,
  $\Lk{4}$ and above: scenarios and system-level behaviour, $\Lk{\infty}$ purpose and
  responsibility. Definition \ref{def:cognitive-scale}. \\
Level contract four-tuple & (participant, invariant assured, contract form, mode of
  verification), plus the oblivious set. Definition \ref{def:level-contract}. \\
Level attribution of a verification obligation & A property $\varphi$ belongs to the lowest
  level that can express it. Definition \ref{def:obligation-assign}. \\
Domain binding & The fraction $\mathrm{Dom}$ of the level's axioms that are
  domain-specific. Definition \ref{def:dom}. \\
Two-layer trust structure & A composite in which the intent layer is assured by consensus
  and the formalisable layer by mechanical assurance. Definition \ref{def:two-layer}. \\
\bottomrule
\end{longtable}

\section*{Symbol Index}

\begin{longtable}{L{3.0cm}L{11.4cm}}
\toprule
\textbf{Symbol} & \textbf{Meaning} \\
\midrule
\endfirsthead
\toprule
\textbf{Symbol} & \textbf{Meaning} \\
\midrule
\endhead
\bottomrule
\endlastfoot
$\Lambda_i$ & level $i$ \\
$\Beh_i$ & behaviour set of level $i$ \\
$\Obs_i, \Conc_i$ & abstraction and concretisation maps \\
$\Env_i, \Cost_i, \Inv_i$ & environmental assumptions, cost model, invariant set \\
$\Obl(\Lambda)$ & oblivious set \\
$\kappa$ & compression ratio \\
$\mathcal{G}, \phi$ & architecture relational graph, abstract skeleton \\
$T, Q, d$ & task, query family, loss functional \\
$\DL, L_T$ & description length, task loss \\
$S$ & behaviour contract \\
$\Tr(B,E)$, $\Allow(S)$ & set of observable traces, set of permitted traces \\
$\mathfrak{A}, \TCB$ & assurance four-tuple, trusted computing base \\
$\Pi$ & participant-resource structure \\
$\mathcal{K}_6$ & the six functions \\
$n, m, k, F$ & number of entities, number of resources, interfaces per entity, fixed cost
  of the substrate \\
$\rho(c), \mathrm{S}(\iota)$ & read-set, set of enforcement sites \\
$D$ & technical debt \\
$\mathrm{Dom}, \mathrm{Conv}$ & domain binding, convergence strength \\
$q, p, \varepsilon_{\text{cons}}$ & correlated error probability, individual error rate,
  consensus error bound \\
\bottomrule
\end{longtable}


\backmatter
% =====================================================================
%  References
%  Note: this list contains four classes of entry---
%  (i) literature directly concerned with intent-to-verifiable-execution;
%  (ii) classical theoretical literature used in the argument;
%  (iii) engineering and organisational theory;
%  (iv) philosophy, methodology, composition, and incentive theory.
%  The complete metadata of every entry is governed by its official publication
%  page; entries with recent nominal years are retained as originally recorded.
% =====================================================================

\begin{thebibliography}{99}

\bibitem{mucha2024}
J. Mucha, A. Kaufmann, and D. Riehle.
\newblock A systematic literature review of pre-requirements specification traceability.
\newblock \emph{Requirements Engineering}, 29:119--141, 2024.
\newblock \url{https://doi.org/10.1007/s00766-023-00412-z}

\bibitem{mokos2020}
K. Mokos and P. Katsaros.
\newblock A survey on the formalisation of system requirements and their validation.
\newblock \emph{Array}, 7:100030, 2020.
\newblock \url{https://doi.org/10.1016/j.array.2020.100030}

\bibitem{osvbench2026}
S. Li, J. Jiang, T. Zhao, and J. Shen.
\newblock OSVBench: Benchmarking LLMs on specification generation tasks for operating
system verification.
\newblock In \emph{Proceedings of AAAI 2026}, 40(37):31698--31707, 2026.
\newblock \url{https://doi.org/10.1609/aaai.v40i37.40437}

\bibitem{coins2026}
F. Yang, X. Li, S. Wang, J. An, Z. Sun, S. Feng, W. Wang, W. Wang, N. Zhan, and F. Xu.
\newblock How powerful are LLMs in generating formal program specifications?
\newblock In \emph{Proceedings of ICML 2026}, PMLR 306:145221--145249, 2026.
\newblock \url{https://proceedings.mlr.press/v306/yang26am.html}

\bibitem{verispecgen2026}
Z. Ye, A. Z. H. Yang, H. Su, Z. Liao, S. Tenka, Z. Qin, U. Ghai, D. Song, and S. Kong.
\newblock Intent-aligned formal specification synthesis via traceable refinement
(VeriSpecGen).
\newblock arXiv:2604.10392, 2026.
\newblock \url{https://arxiv.org/abs/2604.10392}

\bibitem{specsyn2026}
L. Ma, S. Liu, Y. Li, Q. Wu, H. Wang, and L. Bu.
\newblock SpecSyn: LLM-based synthesis and refinement of formal specifications for
real-world program verification.
\newblock arXiv:2604.21570, 2026.
\newblock \url{https://arxiv.org/abs/2604.21570}

\bibitem{specgen2025}
L. Ma, S. Liu, Y. Li, X. Xie, and L. Bu.
\newblock SpecGen: Automated generation of formal program specifications via large
language models.
\newblock In \emph{Proceedings of ICSE 2025}, 2025.
\newblock \url{https://doi.org/10.1109/ICSE55347.2025.00129}

\bibitem{verina2026}
Z. Ye, Z. Yan, J. He, T. Kasriel, K. Yang, and D. Song.
\newblock VERINA: Benchmarking verifiable code generation.
\newblock In \emph{Proceedings of ICLR 2026}, 2026.

\bibitem{kupferman2003}
O. Kupferman and M. Y. Vardi.
\newblock Vacuity detection in temporal model checking.
\newblock In \emph{CHARME 2003}, 2003.
\newblock \url{https://doi.org/10.1007/3-540-48153-2_8}

\bibitem{solarlezama2008}
A. Solar-Lezama.
\newblock \emph{Program Synthesis by Sketching}.
\newblock PhD thesis / Technical Report EECS-2008-177, UC Berkeley, 2008.
\newblock \url{https://www2.eecs.berkeley.edu/Pubs/TechRpts/2008/EECS-2008-177.html}

\bibitem{alur2013}
R. Alur et al.
\newblock Syntax-guided synthesis.
\newblock In \emph{Proceedings of CAV 2013}, 2013.
\newblock \url{https://www.sygus.org/}

\bibitem{kurabayashi2023}
T. Kurabayashi, H. Tanno, H. Kirinuki, and T. Oshima.
\newblock A survey of program synthesis.
\newblock \emph{Computer Software}, 40(2):95--128, 2023.
\newblock \url{https://doi.org/10.11309/jssst.40.2_95}

\bibitem{leandojo2023}
K. Yang, A. Swope, A. Gu, R. Chalamala, P. Song, S. Yu, S. Godil, R. J. Prenger, and
A. Anandkumar.
\newblock LeanDojo: Theorem proving with retrieval-augmented language models.
\newblock In \emph{NeurIPS 2023 Datasets and Benchmarks}, 2023.

\bibitem{leancopilot2024}
P. Song, K. Yang, and A. Anandkumar.
\newblock Towards large language models as copilots for theorem proving in Lean.
\newblock arXiv:2404.12534, 2024.
\newblock \url{https://arxiv.org/abs/2404.12534}

\bibitem{necula1997}
G. C. Necula.
\newblock Proof-carrying code.
\newblock In \emph{Proceedings of POPL 1997}, pages 106--119, 1997.
\newblock \url{https://doi.org/10.1145/263699.263712}

\bibitem{hoare1969}
C. A. R. Hoare.
\newblock An axiomatic basis for computer programming.
\newblock \emph{Communications of the ACM}, 12(10):576--580, 1969.
\newblock \url{https://doi.org/10.1145/363235.363259}

\bibitem{moggi1991}
E. Moggi.
\newblock Notions of computation and monads.
\newblock \emph{Information and Computation}, 93(1):55--92, 1991.
\newblock \url{https://doi.org/10.1016/0890-5401(91)90052-4}

\bibitem{plotkin2002}
G. D. Plotkin and J. Power.
\newblock Notions of computation determine monads.
\newblock In \emph{FOSSACS 2002}, LNCS 2303:342--356, 2002.
\newblock \url{https://doi.org/10.1007/3-540-45931-6_24}

\bibitem{hott2013}
The Univalent Foundations Program.
\newblock \emph{Homotopy Type Theory: Univalent Foundations of Mathematics}.
\newblock Institute for Advanced Study, 2013.
\newblock \url{https://homotopytypetheory.org/book/}

\bibitem{cousot1977}
P. Cousot and R. Cousot.
\newblock Abstract interpretation: A unified lattice model for static analysis of
programs by construction or approximation of fixpoints.
\newblock In \emph{Proceedings of POPL 1977}, pages 238--252, 1977.
\newblock \url{https://doi.org/10.1145/512950.512973}

\bibitem{appel2007}
A. W. Appel and S. Blazy.
\newblock Separation logic for small-step Cminor.
\newblock In \emph{TPHOLs 2007}, LNCS 4732:5--21, 2007.
\newblock \url{https://doi.org/10.1007/978-3-540-74591-4_3}

\bibitem{leroy2009}
X. Leroy.
\newblock Formal verification of a realistic compiler.
\newblock \emph{Communications of the ACM}, 52(7):107--115, 2009.
\newblock \url{https://doi.org/10.1145/1538788.1538814}

\bibitem{necula2000}
G. C. Necula.
\newblock Translation validation for an optimizing compiler.
\newblock In \emph{Proceedings of PLDI 2000}, 2000.
\newblock \url{https://doi.org/10.1145/349299.349314}

\bibitem{cakeml2019}
Y. K. Tan, M. O. Myreen, R. Kumar, A. Fox, S. Owens, and M. Norrish.
\newblock The verified CakeML compiler backend.
\newblock \emph{Journal of Functional Programming}, 29, 2019.
\newblock \url{https://doi.org/10.1017/S0956796818000229}

\bibitem{klein2009}
G. Klein et al.
\newblock seL4: Formal verification of an OS kernel.
\newblock In \emph{Proceedings of SOSP 2009}, pages 207--220, 2009.
\newblock \url{https://sel4.systems/Research/publications.html}

\bibitem{krijnen2024}
J. O. G. Krijnen, M. M. T. Chakravarty, G. Keller, and W. Swierstra.
\newblock Translation certification for smart contracts.
\newblock \emph{Science of Computer Programming}, 233:103051, 2024.
\newblock \url{https://doi.org/10.1016/j.scico.2023.103051}

\bibitem{turing1936}
A. M. Turing.
\newblock On computable numbers, with an application to the Entscheidungsproblem.
\newblock \emph{Proceedings of the London Mathematical Society}, 42:230--265, 1936.
\newblock \url{https://doi.org/10.1112/plms/s2-42.1.230}

\bibitem{rice1953}
H. G. Rice.
\newblock Classes of recursively enumerable sets and their decision problems.
\newblock \emph{Transactions of the AMS}, 74(2):358--366, 1953.

\bibitem{hoffmann2022}
J. Hoffmann and S. Jost.
\newblock Two decades of automatic amortized resource analysis.
\newblock \emph{Mathematical Structures in Computer Science}, 32(6), 2022.
\newblock \url{https://doi.org/10.1017/S0960129521000487}

\bibitem{wilhelm2008}
R. Wilhelm et al.
\newblock The worst-case execution-time problem---overview of methods and survey of tools.
\newblock \emph{ACM Transactions on Embedded Computing Systems}, 7(3), 2008.
\newblock \url{https://doi.org/10.1145/1347375.1347389}

\bibitem{lv2016}
M. Lv, N. Guan, J. Reineke, R. Wilhelm, and W. Yi.
\newblock A survey on static cache analysis for real-time systems.
\newblock \emph{LITES}, 3(1), 2016.
\newblock \url{https://doi.org/10.4230/LITES-v003-i001-a005}

\bibitem{landauer1961}
R. Landauer.
\newblock Irreversibility and heat generation in the computing process.
\newblock \emph{IBM Journal of Research and Development}, 5:183--191, 1961.
\newblock \url{https://doi.org/10.1147/rd.53.0183}

\bibitem{georgescu2021}
I. Georgescu.
\newblock 60 years of Landauer's principle.
\newblock \emph{Nature Reviews Physics}, 3:770, 2021.
\newblock \url{https://doi.org/10.1038/s42254-021-00400-8}

\bibitem{taleb2023}
R. Taleb, S. Hall\'e, and R. Khoury.
\newblock Uncertainty in runtime verification: A survey.
\newblock \emph{Computer Science Review}, 50:100594, 2023.
\newblock \url{https://doi.org/10.1016/j.cosrev.2023.100594}

\bibitem{henzinger1998}
T. A. Henzinger, S. Qadeer, and S. K. Rajamani.
\newblock You assume, we guarantee: Methodology and case studies.
\newblock In \emph{Proceedings of CAV 1998}, 1998.

\bibitem{sharf2021}
M. Sharf, B. Besselink, K. H. Johansson, et al.
\newblock Assume/guarantee contracts for dynamical systems: Theory and computational
tools.
\newblock \emph{IFAC-PapersOnLine}, 54(5):25--30, 2021.
\newblock \url{https://doi.org/10.1016/j.ifacol.2021.08.469}

\bibitem{sharf2024}
M. Sharf, B. Besselink, and K. H. Johansson.
\newblock Contract composition for dynamical control systems: Definition and
verification using linear programming.
\newblock \emph{Automatica}, 164:111637, 2024.
\newblock \url{https://doi.org/10.1016/j.automatica.2024.111637}

\bibitem{shali2026}
B. Shali, A. J. van der Schaft, and B. Besselink.
\newblock Design and control for implementation of simulation-based assume-guarantee
contracts.
\newblock \emph{IEEE Transactions on Automatic Control}, 71(1):65--80, 2026.
\newblock \url{https://doi.org/10.1109/TAC.2025.3587874}

\bibitem{simon1962}
H. A. Simon.
\newblock The architecture of complexity.
\newblock \emph{Proceedings of the American Philosophical Society}, 106(6):467--482,
1962.

\bibitem{grunwald2006}
P. D. Gr\"unwald.
\newblock The minimum description length principle in coding and modeling.
\newblock \emph{IEEE Transactions on Information Theory}, 52(2), 2006.
\newblock \url{https://doi.org/10.1109/TIT.2005.862128}

\bibitem{grunwald2019}
P. Gr\"unwald and T. Roos.
\newblock Minimum description length revisited.
\newblock arXiv:1908.08484, 2019.
\newblock \url{https://arxiv.org/abs/1908.08484}

\bibitem{goguen1982}
J. A. Goguen and J. Meseguer.
\newblock Security policies and security models.
\newblock In \emph{IEEE Symposium on Security and Privacy 1982}, pages 11--20, 1982.
\newblock \url{https://doi.org/10.1109/SP.1982.10014}

\bibitem{feser2015}
M. C. Feser, S. Chaudhuri, and I. Dillig.
\newblock Synthesizing data structure traversals from local specifications.
\newblock In \emph{Proceedings of PLDI 2015}, 2015.
\newblock \url{https://doi.org/10.1145/2737924.2738007}

\bibitem{rondon2008}
P. M. Rondon, M. Kawaguchi, and R. Jhala.
\newblock Liquid types.
\newblock In \emph{Proceedings of PLDI 2008}, 2008.
\newblock \url{https://doi.org/10.1145/1375581.1375602}

\bibitem{z32008}
L. de Moura and N. Bj{\o}rner.
\newblock Z3: An efficient SMT solver.
\newblock In \emph{TACAS 2008}, LNCS 4963, 2008.
\newblock \url{https://doi.org/10.1007/978-3-540-78800-3_24}

\bibitem{leucker2009}
M. Leucker and C. Schallhart.
\newblock Runtime verification for LTL and TLTL.
\newblock In \emph{TACAS 2009}, 2009.
\newblock \url{https://doi.org/10.1007/978-3-642-00768-2_17}

\bibitem{howard1980}
W. A. Howard.
\newblock The formulae-as-types notion of construction.
\newblock In \emph{To H. B. Curry: Essays on Combinatory Logic, Lambda Calculus and
Formalism}, 1980.

\bibitem{reqgen2009}
On the generation of requirements specifications from software engineering models:
A systematic literature review.
\newblock \emph{Information and Software Technology}, 51(9):1291--1307, 2009.
\newblock \url{https://doi.org/10.1016/j.infsof.2009.04.001}

% ---------- Classical theoretical literature ----------

\bibitem{parnas1972}
D. L. Parnas.
\newblock On the criteria to be used in decomposing systems into modules.
\newblock \emph{Communications of the ACM}, 15(12):1053--1058, 1972.

\bibitem{conway1968}
M. E. Conway.
\newblock How do committees invent?
\newblock \emph{Datamation}, 14(4):28--31, 1968.

\bibitem{brooks1975}
F. P. Brooks.
\newblock \emph{The Mythical Man-Month}.
\newblock Addison-Wesley, 1975.

\bibitem{boehm1984}
B. W. Boehm.
\newblock Verifying and validating software requirements and design specifications.
\newblock \emph{IEEE Software}, 1(1):75--88, 1984.

\bibitem{ousterhout2018}
J. Ousterhout.
\newblock \emph{A Philosophy of Software Design}.
\newblock Yaknyam Press, 2018.

\bibitem{lamport2002}
L. Lamport.
\newblock \emph{Specifying Systems: The TLA+ Language and Tools for Hardware and
Software Engineers}.
\newblock Addison-Wesley, 2002.

\bibitem{shannon1948}
C. E. Shannon.
\newblock A mathematical theory of communication.
\newblock \emph{Bell System Technical Journal}, 27:379--423, 1948.

\bibitem{rissanen1978}
J. Rissanen.
\newblock Modeling by shortest data description.
\newblock \emph{Automatica}, 14(5):465--471, 1978.

\bibitem{kolmogorov1965}
A. N. Kolmogorov.
\newblock Three approaches to the quantitative definition of information.
\newblock \emph{Problems of Information Transmission}, 1(1):1--7, 1965.

\bibitem{hoeffding1963}
W. Hoeffding.
\newblock Probability inequalities for sums of bounded random variables.
\newblock \emph{Journal of the American Statistical Association}, 58(301):13--30,
1963.

\bibitem{condorcet1785}
M. de Condorcet.
\newblock \emph{Essai sur l'application de l'analyse \`a la probabilit\'e des
d\'ecisions rendues \`a la pluralit\'e des voix}.
\newblock Paris, 1785.

\bibitem{paulson1989}
L. C. Paulson.
\newblock The foundation of a generic theorem prover.
\newblock \emph{Journal of Automated Reasoning}, 5(3):363--397, 1989.

\bibitem{nygard2018}
M. T. Nygard.
\newblock \emph{Release It! Design and Deploy Production-Ready Software}.
\newblock Second edition, Pragmatic Bookshelf, 2018.

\bibitem{evans2003}
E. Evans.
\newblock \emph{Domain-Driven Design: Tackling Complexity in the Heart of Software}.
\newblock Addison-Wesley, 2003.

% ---------- Philosophy and methodology ----------

\bibitem{dennett1991}
D. C. Dennett.
\newblock Real patterns.
\newblock \emph{The Journal of Philosophy}, 88(1):27--51, 1991.

\bibitem{quine1951}
W. V. O. Quine.
\newblock Two dogmas of empiricism.
\newblock \emph{The Philosophical Review}, 60(1):20--43, 1951.

\bibitem{lakatos1970}
I. Lakatos.
\newblock Falsification and the methodology of scientific research programmes.
\newblock In \emph{Criticism and the Growth of Knowledge}, Cambridge University Press,
1970.

\bibitem{pearl2009}
J. Pearl.
\newblock \emph{Causality: Models, Reasoning, and Inference}.
\newblock Second edition, Cambridge University Press, 2009.

% ---------- Composition, coordination and incentives ----------

\bibitem{leinster2004}
T. Leinster.
\newblock \emph{Higher Operads, Higher Categories}.
\newblock London Mathematical Society Lecture Note Series 298, Cambridge University
Press, 2004.

\bibitem{jacobs1999}
B. Jacobs.
\newblock \emph{Categorical Logic and Type Theory}.
\newblock North-Holland, 1999.

\bibitem{calm2011}
P. Alvaro, N. Conway, J. M. Hellerstein, and W. R. Marczak.
\newblock Consistency analysis in Bloom: A CALM and collected approach.
\newblock In \emph{Proceedings of CIDR 2011}, 2011.

\bibitem{brier1950}
G. W. Brier.
\newblock Verification of forecasts expressed in terms of probability.
\newblock \emph{Monthly Weather Review}, 78(1):1--3, 1950.

\bibitem{shannon1959}
C. E. Shannon.
\newblock Coding theorems for a discrete source with a fidelity criterion.
\newblock \emph{IRE National Convention Record}, 7(4):142--163, 1959.

\bibitem{kalman1960}
R. E. Kalman.
\newblock On the general theory of control systems.
\newblock In \emph{Proceedings of the First International Congress on Automatic
Control}, 1960.

\end{thebibliography}


\end{document}
